\input{preamble}

% OK, start here.
%
\begin{document}

\title{Chow Homology and Chern Classes}

\maketitle

\phantomsection
\label{section-phantom}


\tableofcontents


\section{Introduction}
\label{section-introduction}

\noindent
In this chapter we discuss Chow homology groups and the construction
of Chern classes of vector bundles as elements of operational
Chow cohomology groups (everything with $\mathbf{Z}$-coefficients).

\medskip\noindent
We start this chapter by giving the shortest possible
algebraic proof of the Key Lemma \ref{lemma-milnor-gersten-low-degree}.
We first define the Herbrand quotient
(Section \ref{section-periodic-complexes})
and we compute it in some cases
(Section \ref{section-calculation}).
Next, we prove some simple algebra lemmas on
existence of suitable factorizations after modifications
(Section \ref{section-preparation-tame-symbol}).
Using these we construct/define the tame symbol in
Section \ref{section-tame-symbol}.
Only the most basic properties of the tame symbol
are needed to prove the Key Lemma, which we do
in Section \ref{section-key-lemma}.

\medskip\noindent
Next, we introduce the basic setup we work with in the rest of this
chapter in Section \ref{section-setup}. To make the material a little
bit more challenging we decided to treat a somewhat more general case
than is usually done. Namely we assume our schemes $X$ are locally of
finite type over a fixed locally Noetherian base scheme which is universally
catenary and is endowed with a dimension function. These assumptions suffice
to be able to define the Chow homology groups $\CH_*(X)$ and the action of
capping with Chern classes on them. This is an indication that we should
be able to define these also for algebraic stacks locally of finite type
over such a base.

\medskip\noindent
Next, we follow the first few chapters of \cite{F} in order to define
cycles, flat pullback, proper pushforward, and rational equivalence,
except that we have been less precise about the supports of the cycles
involved.

\medskip\noindent
We diverge from the presentation given in \cite{F} by using the
Key lemma mentioned above to prove a basic commutativity relation in
Section \ref{section-key}. Using this we prove that the operation
of intersecting with an invertible sheaf passes through rational
equivalence and is commutative, see Section \ref{section-commutativity}.
One more application of the Key
lemma proves that the Gysin map of an effective Cartier divisor
passes through rational equivalence, see Section \ref{section-gysin}.
Having proved this, it is straightforward to define Chern
classes of vector bundles, prove additivity, prove the splitting principle,
introduce Chern characters, Todd classes, and state the
Grothendieck-Riemann-Roch theorem.

\medskip\noindent
There are two appendices. In Appendix A (Section \ref{section-appendix-A})
we discuss an alternative (longer) construction of the
tame symbol and corresponding proof of the Key Lemma.
Finally, in Appendix B (Section \ref{section-appendix-chow})
we briefly discuss the relationship with $K$-theory of coherent
sheaves and we discuss some blowup lemmas.
We suggest the reader look at their introductions for
more information.

\medskip\noindent
We will return to the Chow groups $\CH_*(X)$ for smooth projective varieties
over algebraically closed fields in the next chapter. Using a moving
lemma as in \cite{Samuel}, \cite{ChevalleyI}, and \cite{ChevalleyII}
and Serre's Tor-formula
(see \cite{Serre_local_algebra} or \cite{Serre_algebre_locale})
we will define a ring structure on $\CH_*(X)$. See
Intersection Theory, Section \ref{intersection-section-introduction} ff.








\section{Periodic complexes and Herbrand quotients}
\label{section-periodic-complexes}

\noindent
Of course there is a very general notion of periodic complexes.
We can require periodicity of the maps, or periodicity of the objects.
We will add these here as needed. For the moment we only need
the following cases.

\begin{definition}
\label{definition-periodic-complex}
Let $R$ be a ring.
\begin{enumerate}
\item A {\it $2$-periodic complex} over $R$ is given
by a quadruple $(M, N, \varphi, \psi)$ consisting of
$R$-modules $M$, $N$ and $R$-module maps $\varphi : M \to N$,
$\psi : N \to M$ such that
$$
\xymatrix{
\ldots \ar[r] &
M \ar[r]^\varphi &
N \ar[r]^\psi &
M \ar[r]^\varphi &
N \ar[r] & \ldots
}
$$
is a complex. In this setting we define the {\it cohomology modules}
of the complex to be the $R$-modules
$$
H^0(M, N, \varphi, \psi) = \Ker(\varphi)/\Im(\psi)
\quad\text{and}\quad
H^1(M, N, \varphi, \psi) = \Ker(\psi)/\Im(\varphi).
$$
We say the $2$-periodic complex is {\it exact} if the cohomology
groups are zero.
\item A {\it $(2, 1)$-periodic complex} over $R$ is given
by a triple $(M, \varphi, \psi)$ consisting of an $R$-module $M$ and
$R$-module maps $\varphi : M \to M$, $\psi : M \to M$
such that
$$
\xymatrix{
\ldots \ar[r] &
M \ar[r]^\varphi &
M \ar[r]^\psi &
M \ar[r]^\varphi &
M \ar[r] & \ldots
}
$$
is a complex. Since this is a special case of a $2$-periodic complex
we have its {\it cohomology modules} $H^0(M, \varphi, \psi)$,
$H^1(M, \varphi, \psi)$ and a notion of exactness.
\end{enumerate}
\end{definition}

\noindent
In the following we will use any result proved for $2$-periodic
complexes without further mention for $(2, 1)$-periodic complexes.
It is clear that the collection of $2$-periodic complexes forms a
category with morphisms
$(f, g) : (M, N, \varphi, \psi) \to (M', N', \varphi', \psi')$
pairs of morphisms $f : M \to M'$ and $g : N \to N'$ such
that $\varphi' \circ f = g \circ \varphi$ and $\psi' \circ g = f \circ \psi$.
We obtain an abelian category, with kernels and cokernels as in
Homology, Lemma \ref{homology-lemma-cat-chain-abelian}.

\begin{definition}
\label{definition-periodic-length}
Let $(M, N, \varphi, \psi)$ be a $2$-periodic complex
over a ring $R$ whose cohomology modules have finite length.
In this case we define the {\it multiplicity} of $(M, N, \varphi, \psi)$
to be the integer
$$
e_R(M, N, \varphi, \psi) =
\text{length}_R(H^0(M, N, \varphi, \psi))
-
\text{length}_R(H^1(M, N, \varphi, \psi))
$$
In the case of a $(2, 1)$-periodic complex $(M, \varphi, \psi)$,
we denote this by $e_R(M, \varphi, \psi)$ and we will sometimes call this
the {\it (additive) Herbrand quotient}.
\end{definition}

\noindent
If the cohomology groups of $(M, \varphi, \psi)$
are finite abelian groups, then it is customary to call the
{\it (multiplicative) Herbrand quotient}
$$
q(M, \varphi, \psi) =
\frac{\# H^0(M, \varphi, \psi)}{\# H^1(M, \varphi, \psi)}
$$
In words: the multiplicative Herbrand quotient is the number of elements of
$H^0$ divided by the number of elements of $H^1$. If $R$ is local and if
the residue field of $R$ is finite with $q$ elements, then we see that
$$
q(M, \varphi, \psi) = q^{e_R(M, \varphi, \psi)}
$$

\medskip\noindent
An example of a $(2, 1)$-periodic complex over a ring $R$ is any triple of
the form $(M, 0, \psi)$ where $M$ is an $R$-module and $\psi$ is an
$R$-linear map. If the kernel and cokernel of $\psi$ have finite length,
then we obtain
\begin{equation}
\label{equation-multiplicity-coker-ker}
e_R(M, 0, \psi) = \text{length}_R(\Coker(\psi)) - \text{length}_R(\Ker(\psi))
\end{equation}
We state and prove the obligatory lemmas on these notations.

\begin{lemma}
\label{lemma-additivity-periodic-length}
Let $R$ be a ring. Suppose that we have a short exact sequence of
$2$-periodic complexes
$$
0 \to (M_1, N_1, \varphi_1, \psi_1)
\to (M_2, N_2, \varphi_2, \psi_2)
\to (M_3, N_3, \varphi_3, \psi_3)
\to 0
$$
If two out of three have cohomology modules of finite length so does
the third and we have
$$
e_R(M_2, N_2, \varphi_2, \psi_2) =
e_R(M_1, N_1, \varphi_1, \psi_1) +
e_R(M_3, N_3, \varphi_3, \psi_3).
$$
\end{lemma}

\begin{proof}
We abbreviate $A = (M_1, N_1, \varphi_1, \psi_1)$,
$B = (M_2, N_2, \varphi_2, \psi_2)$ and $C = (M_3, N_3, \varphi_3, \psi_3)$.
We have a long exact cohomology sequence
$$
\ldots
\to H^1(C)
\to H^0(A)
\to H^0(B)
\to H^0(C)
\to H^1(A)
\to H^1(B)
\to H^1(C)
\to \ldots
$$
This gives a finite exact sequence
$$
0 \to I
\to H^0(A)
\to H^0(B)
\to H^0(C)
\to H^1(A)
\to H^1(B)
\to K \to 0
$$
with $0 \to K \to H^1(C) \to I \to 0$ a filtration. By additivity of
the length function (Algebra, Lemma \ref{algebra-lemma-length-additive})
we see the result.
\end{proof}

\begin{lemma}
\label{lemma-finite-periodic-length}
Let $R$ be a ring. If $(M, N, \varphi, \psi)$ is a $2$-periodic complex
such that $M$, $N$ have finite length, then
$e_R(M, N, \varphi, \psi) = \text{length}_R(M) - \text{length}_R(N)$.
In particular, if $(M, \varphi, \psi)$ is a $(2, 1)$-periodic complex
such that $M$ has finite length, then
$e_R(M, \varphi, \psi) = 0$.
\end{lemma}

\begin{proof}
Observe that on the category of $2$-periodic complexes with $M$, $N$
of finite length the quantity ``$\text{length}_R(M) - \text{length}_R(N)$''
is additive in short exact sequences (precise statement left to the reader).
Consider the short exact sequence
$$
0 \to (M, \Im(\varphi), \varphi, 0) \to
(M, N, \varphi, \psi) \to (0, N/\Im(\varphi), 0, 0) \to 0
$$
The initial remark combined with the additivity of
Lemma \ref{lemma-additivity-periodic-length}
reduces us to the cases (a) $M = 0$ and (b) $\varphi$ is surjective.
We leave those cases to the reader.
\end{proof}

\begin{lemma}
\label{lemma-compare-periodic-lengths}
Let $R$ be a ring. Let $f : (M, \varphi, \psi) \to (M', \varphi', \psi')$
be a map of $(2, 1)$-periodic complexes whose cohomology modules
have finite length. If $\Ker(f)$ and $\Coker(f)$ have finite length,
then $e_R(M, \varphi, \psi) = e_R(M', \varphi', \psi')$.
\end{lemma}

\begin{proof}
Apply the additivity of Lemma \ref{lemma-additivity-periodic-length}
and observe that $(\Ker(f), \varphi, \psi)$ and
$(\Coker(f), \varphi', \psi')$ have vanishing multiplicity by
Lemma \ref{lemma-finite-periodic-length}.
\end{proof}




\section{Calculation of some multiplicities}
\label{section-calculation}

\noindent
To prove equality of certain cycles later on we need to
compute some multiplicities. Our main tool, besides the
elementary lemmas on multiplicities given in the previous section,
will be Algebra, Lemma \ref{algebra-lemma-order-vanishing-determinant}.

\begin{lemma}
\label{lemma-length-multiplication}
Let $R$ be a Noetherian local ring.
Let $M$ be a finite $R$-module. Let $x \in R$. Assume that
\begin{enumerate}
\item $\dim(\text{Supp}(M)) \leq 1$, and
\item $\dim(\text{Supp}(M/xM)) \leq 0$.
\end{enumerate}
Write
$\text{Supp}(M) = \{\mathfrak m, \mathfrak q_1, \ldots, \mathfrak q_t\}$.
Then
$$
e_R(M, 0, x) =
\sum\nolimits_{i = 1, \ldots, t}
\text{ord}_{R/\mathfrak q_i}(x)
\text{length}_{R_{\mathfrak q_i}}(M_{\mathfrak q_i}).
$$
\end{lemma}

\begin{proof}
We first make some preparatory remarks.
The result of the lemma holds if $M$ has finite length, i.e., if $t = 0$,
because both the left hand side and the right hand side are zero
in this case, see Lemma \ref{lemma-finite-periodic-length}.
Also, if we have a short exact sequence $0 \to M \to M' \to M'' \to 0$
of modules satisfying (1) and (2), then lemma for 2 out of 3
of these implies the lemma for the third by the
additivity of length (Algebra, Lemma \ref{algebra-lemma-length-additive}) and
additivty of multiplicities (Lemma \ref{lemma-additivity-periodic-length}).

\medskip\noindent
Denote $M_i$ the image of $M$ in $M_{\mathfrak q_i}$, so
$\text{Supp}(M_i) = \{\mathfrak m, \mathfrak q_i\}$.
The kernel and cokernel of the map $M \to \bigoplus M_i$
have support $\{\mathfrak m\}$ and hence have finite length.
By our preparatory remarks, it follows that it suffices to
prove the lemma for each $M_i$. Thus we may assume that
$\text{Supp}(M) = \{\mathfrak m, \mathfrak q\}$.
In this case we have a finite filtration
$M \supset \mathfrak qM \supset \mathfrak q^2M \supset \ldots \supset
\mathfrak q^nM = 0$ by Algebra, Lemma
\ref{algebra-lemma-Noetherian-power-ideal-kills-module}.
Again additivity shows that it suffices to prove the lemma
in the case $M$ is annihilated by $\mathfrak q$.
In this case we can view $M$ as a $R/\mathfrak q$-module,
i.e., we may assume that $R$ is a Noetherian local domain
of dimension $1$ with fraction field $K$.
Dividing by the torsion submodule, i.e., by the
kernel of $M \to M \otimes_R K = V$ (the torsion has
finite length hence is handled by our preliminary remarks)
we may assume that $M \subset V$ is a lattice
(Algebra, Definition \ref{algebra-definition-lattice}).
Then $x : M \to M$ is injective and
$\text{length}_R(M/xM) = d(M, xM)$
(Algebra, Definition \ref{algebra-definition-distance}). Since
$\text{length}_K(V) = \dim_K(V)$
we see that $\det(x : V \to V) = x^{\dim_K(V)}$ and
$\text{ord}_R(\det(x : V \to V)) = \dim_K(V) \text{ord}_R(x)$.
Thus the desired equality follows from
Algebra, Lemma \ref{algebra-lemma-order-vanishing-determinant}
in this case.
\end{proof}

\begin{lemma}
\label{lemma-additivity-divisors-restricted}
Let $R$ be a Noetherian local ring.
Let $x \in R$. If $M$ is a finite Cohen-Macaulay module over $R$
with $\dim(\text{Supp}(M)) = 1$ and $\dim(\text{Supp}(M/xM)) = 0$, then
$$
\text{length}_R(M/xM)
=
\sum\nolimits_i \text{length}_R(R/(x, \mathfrak q_i))
\text{length}_{R_{\mathfrak q_i}}(M_{\mathfrak q_i}).
$$
where $\mathfrak q_1, \ldots, \mathfrak q_t$ are the
minimal primes of the support of $M$. If $I \subset R$ is an ideal
such that $x$ is a nonzerodivisor on $R/I$ and $\dim(R/I) = 1$, then
$$
\text{length}_R(R/(x, I))
=
\sum\nolimits_i \text{length}_R(R/(x, \mathfrak q_i))
\text{length}_{R_{\mathfrak q_i}}((R/I)_{\mathfrak q_i})
$$
where $\mathfrak q_1, \ldots, \mathfrak q_n$ are the minimal
primes over $I$.
\end{lemma}

\begin{proof}
These are special cases of Lemma \ref{lemma-length-multiplication}.
\end{proof}

\noindent
Here is another case where we can determine the value of a multiplicity.

\begin{lemma}
\label{lemma-powers-period-length-zero}
Let $R$ be a ring. Let $M$ be an $R$-module.
Let $\varphi : M \to M$ be an endomorphism and $n > 0$
such that $\varphi^n = 0$ and such that $\Ker(\varphi)/\Im(\varphi^{n - 1})$
has finite length as an $R$-module.
Then
$$
e_R(M, \varphi^i, \varphi^{n - i}) = 0
$$
for $i = 0, \ldots, n$.
\end{lemma}

\begin{proof}
The cases $i = 0, n$ are trivial as $\varphi^0 = \text{id}_M$ by convention.
Let us think of $M$ as an $R[t]$-module where multiplication by $t$
is given by $\varphi$. Let us write
$K_i = \Ker(t^i : M \to M)$ and
$$
a_i = \text{length}_R(K_i/t^{n - i}M),\quad
b_i = \text{length}_R(K_i/tK_{i + 1}),\quad
c_i = \text{length}_R(K_1/t^iK_{i + 1})
$$
Boundary values are $a_0 = a_n = b_0 = c_0 = 0$.
The $c_i$ are integers for $i < n$ as $K_1/t^iK_{i + 1}$
is a quotient of $K_1/t^{n - 1}M$ which is assumed to have finite length.
We will use frequently that $K_i \cap t^jM = t^jK_{i + j}$.
For $0 < i < n - 1$ we have an exact sequence
$$
0 \to
K_1/t^{n - i - 1}K_{n - i} \to
K_{i + 1}/t^{n - i - 1}M \xrightarrow{t} K_i/t^{n - i}M
\to K_i/tK_{i + 1} \to 0
$$
By induction on $i$ we conclude that $a_i$ and $b_i$ are
integers for $i < n$ and that
$$
c_{n - i - 1} - a_{i + 1} + a_i - b_i = 0
$$
For $0 < i < n - 1$ there is a short exact sequence
$$
0 \to
K_i/tK_{i + 1} \to
K_{i + 1}/tK_{i + 2} \xrightarrow{t^i}
K_1/t^{i + 1}K_{i + 2} \to
K_1/t^iK_{i + 1} \to 0
$$
which gives
$$
b_i - b_{i + 1} + c_{i + 1} - c_i = 0
$$
Since $b_0 = c_0$ we conclude that $b_i = c_i$ for $i < n$.
Then we see that
$$
a_2 = a_1 + b_{n - 2} - b_1,\quad
a_3 = a_2 + b_{n - 3} - b_2,\quad \ldots
$$
It is straightforward to see that this implies $a_i = a_{n - i}$ as desired.
\end{proof}

\begin{lemma}
\label{lemma-multiply-period-length}
Let $(R, \mathfrak m)$ be a Noetherian local ring. Let
$(M, \varphi, \psi)$ be a $(2, 1)$-periodic complex over $R$
with $M$ finite and with cohomology groups of finite length over $R$.
Let $x \in R$ be such that $\dim(\text{Supp}(M/xM)) \leq 0$. Then
$$
e_R(M, x\varphi, \psi) = e_R(M, \varphi, \psi) - e_R(\Im(\varphi), 0, x)
$$
and
$$
e_R(M, \varphi, x\psi) = e_R(M, \varphi, \psi) + e_R(\Im(\psi), 0, x)
$$
\end{lemma}

\begin{proof}
We will only prove the first formula as the second is proved
in exactly the same manner.
Let $M' = M[x^\infty]$ be the $x$-power torsion submodule of $M$.
Consider the short exact sequence $0 \to M' \to M \to M'' \to 0$.
Then $M''$ is $x$-power torsion free (More on Algebra, Lemma
\ref{more-algebra-lemma-divide-by-torsion}).
Since $\varphi$, $\psi$ map $M'$ into $M'$
we obtain a short exact sequence
$$
0 \to (M', \varphi', \psi') \to (M, \varphi, \psi) \to
(M'', \varphi'', \psi'') \to 0
$$
of $(2, 1)$-periodic complexes. Also, we get a short exact sequence
$0 \to M' \cap \Im(\varphi) \to \Im(\varphi) \to \Im(\varphi'') \to 0$.
We have
$e_R(M', \varphi, \psi) = e_R(M', x\varphi, \psi) =
e_R(M' \cap \Im(\varphi), 0, x) = 0$
by Lemma \ref{lemma-compare-periodic-lengths}.
By additivity (Lemma \ref{lemma-additivity-periodic-length})
we see that it suffices to prove the lemma for $(M'', \varphi'', \psi'')$.
This reduces us to the case discussed in the next paragraph.

\medskip\noindent
Assume $x : M \to M$ is injective.
In this case $\Ker(x\varphi) = \Ker(\varphi)$.
On the other hand we have a short exact sequence
$$
0 \to \Im(\varphi)/x\Im(\varphi) \to
\Ker(\psi)/\Im(x\varphi) \to \Ker(\psi)/\Im(\varphi) \to 0
$$
This together with (\ref{equation-multiplicity-coker-ker}) proves the formula.
\end{proof}







\section{Preparation for tame symbols}
\label{section-preparation-tame-symbol}

\noindent
In this section we put some lemma that will help us define the
tame symbol in the next section.

\begin{lemma}
\label{lemma-glue-at-max}
Let $A$ be a Noetherian ring. Let $\mathfrak m_1, \ldots, \mathfrak m_r$
be pairwise distinct maximal ideals of $A$. For $i = 1, \ldots, r$ let
$\varphi_i : A_{\mathfrak m_i} \to B_i$ be a ring map whose
kernel and cokernel are annihilated by a power
of $\mathfrak m_i$. Then there exists a ring map $\varphi : A \to B$ such
that
\begin{enumerate}
\item the localization of $\varphi$ at $\mathfrak m_i$ is
isomorphic to $\varphi_i$, and
\item $\Ker(\varphi)$ and $\Coker(\varphi)$ are annihilated
by a power of $\mathfrak m_1 \cap \ldots \cap \mathfrak m_r$.
\end{enumerate}
Moreover, if each $\varphi_i$ is finite, injective, or
surjective then so is $\varphi$.
\end{lemma}

\begin{proof}
Set $I = \mathfrak m_1 \cap \ldots \cap \mathfrak m_r$. Set
$A_i = A_{\mathfrak m_i}$ and $A' = \prod A_i$.
Then $IA' = \prod \mathfrak m_i A_i$ and $A \to A'$
is a flat ring map such that $A/I \cong A'/IA'$.
Thus we may use More on Algebra, Lemma
\ref{more-algebra-lemma-application-formal-glueing}
to see that there exists an $A$-module map $\varphi : A \to B$
with $\varphi_i$ isomorphic to the localization of $\varphi$
at $\mathfrak m_i$. Then we can use the discussion in
More on Algebra, Remark \ref{more-algebra-remark-formal-glueing-algebras}
to endow $B$ with an $A$-algebra structure
matching the given $A$-algebra structure on $B_i$.
The final statement of the lemma follows easily from
the fact that $\Ker(\varphi)_{\mathfrak m_i} \cong \Ker(\varphi_i)$
and $\Coker(\varphi)_{\mathfrak m_i} \cong \Coker(\varphi_i)$.
\end{proof}

\noindent
The following lemma is very similar to
Algebra, Lemma \ref{algebra-lemma-nonregular-dimension-one}.

\begin{lemma}
\label{lemma-Noetherian-domain-dim-1-two-elements}
Let $(R, \mathfrak m)$ be a Noetherian local ring of dimension $1$.
Let $a, b \in R$ be nonzerodivisors.
There exists a finite ring extension $R \subset R'$
with $R'/R$ annihilated by a power of $\mathfrak m$
and nonzerodivisors $t, a', b' \in R'$ such that
$a = ta'$ and $b = tb'$ and $R' = a'R' + b'R'$.
\end{lemma}

\begin{proof}
If $a$ or $b$ is a unit, then the lemma is true with $R = R'$.
Thus we may assume $a, b \in \mathfrak m$.
Set $I = (a, b)$. The idea is to blow up $R$ in $I$.
Instead of doing the algebraic argument we work geometrically.
Let $X = \text{Proj}(\bigoplus_{d \geq 0} I^d)$.
By Divisors, Lemma
\ref{divisors-lemma-blowing-up-gives-effective-Cartier-divisor}
the morphism $X \to \Spec(R)$ is an isomorphism over
the punctured spectrum $U = \Spec(R) \setminus \{\mathfrak m\}$.
Thus we may and do view $U$ as an open subscheme of $X$.
The morphism $X \to \Spec(R)$ is projective by
Divisors, Lemma \ref{divisors-lemma-blowing-up-projective}.
Also, every generic point of $X$ lies in $U$, for example
by Divisors, Lemma \ref{divisors-lemma-blow-up-and-irreducible-components}.
It follows from Varieties, Lemma \ref{varieties-lemma-finite-in-codim-1}
that $X \to \Spec(R)$ is finite. Thus $X = \Spec(R')$ is
affine and $R \to R'$ is finite. We have $R_a \cong R'_a$ as $U = D(a)$.
Hence a power of $a$ annihilates the finite $R$-module $R'/R$.
As $\mathfrak m = \sqrt{(a)}$ we see that $R'/R$ is annihilated
by a power of $\mathfrak m$. By
Divisors, Lemma \ref{divisors-lemma-blowing-up-gives-effective-Cartier-divisor}
we see that $IR'$ is a locally principal ideal.
Since $R'$ is semi-local we see that $IR'$ is principal,
see Algebra, Lemma \ref{algebra-lemma-locally-free-semi-local-free},
say $IR' = (t)$. Then we have $a = a't$ and $b = b't$ and everything is
clear.
\end{proof}

\begin{lemma}
\label{lemma-not-infinitely-divisible}
Let $(R, \mathfrak m)$ be a Noetherian local ring of dimension $1$.
Let $a, b \in R$ be nonzerodivisors with $a \in \mathfrak m$.
There exists an integer $n = n(R, a, b)$ such that for a finite ring
extension $R \subset R'$ if $b = a^m c$ for some $c \in R'$, then $m \leq n$.
\end{lemma}

\begin{proof}
Choose a minimal prime $\mathfrak q \subset R$. Observe that
$\dim(R/\mathfrak q) = 1$, in particular $R/\mathfrak q$ is not a field.
We can choose a discrete valuation ring $A$ dominating $R/\mathfrak q$
with the same fraction field, see
Algebra, Lemma \ref{algebra-lemma-dominate-by-dimension-1}. Observe that
$a$ and $b$ map to nonzero elements of $A$ as nonzerodivisors in $R$
are not contained in $\mathfrak q$. Let $v$ be the discrete valuation on $A$.
Then $v(a) > 0$ as $a \in \mathfrak m$.
We claim $n = v(b)/v(a)$ works.

\medskip\noindent
Let $R \subset R'$ be given. Set $A' = A \otimes_R R'$.
Since $\Spec(R') \to \Spec(R)$ is surjective
(Algebra, Lemma \ref{algebra-lemma-integral-overring-surjective})
also $\Spec(A') \to \Spec(A)$ is surjective
(Algebra, Lemma \ref{algebra-lemma-surjective-spec-radical-ideal}).
Pick a prime $\mathfrak q' \subset A'$ lying over $(0) \subset A$.
Then $A \subset A'' = A'/\mathfrak q'$ is a finite extension of rings
(again inducing a surjection on spectra).
Pick a maximal ideal $\mathfrak m'' \subset A''$
lying over the maximal ideal of $A$ and a discrete valuation ring
$A'''$ dominating $A''_{\mathfrak m''}$ (see lemma cited above).
Then $A \to A'''$ is an extension of discrete valuation rings
and we have $b = a^m c$ in $A'''$. Thus $v'''(b) \geq mv'''(a)$.
Since $v''' = ev$ where $e$ is the ramification index
of $A'''/A$, we find that $m \leq n$ as desired.
\end{proof}

\begin{lemma}
\label{lemma-prepare-tame-symbol}
Let $(A, \mathfrak m)$ be a Noetherian local ring of dimension $1$.
Let $r \geq 2$ and let $a_1, \ldots, a_r \in A$ be nonzerodivisors
not all units.
Then there exist
\begin{enumerate}
\item a finite ring extension $A \subset B$ with
$B/A$ annihilated by a power of $\mathfrak m$,
\item for each maximal ideal $\mathfrak m_j \subset B$
a nonzerodivisor $\pi_j \in B_j = B_{\mathfrak m_j}$, and
\item factorizations $a_i = u_{i, j} \pi_j^{e_{i, j}}$ in $B_j$
with $u_{i, j} \in B_j$ units and $e_{i, j} \geq 0$.
\end{enumerate}
\end{lemma}

\begin{proof}
Since at least one $a_i$ is not a unit we find that $\mathfrak m$
is not an associated prime of $A$. Moreover, for any $A \subset B$
as in the statement $\mathfrak m$ is not an associated prime of $B$
and $\mathfrak m_j$ is not an associate prime of $B_j$.
Keeping this in mind will help check the arguments below.

\medskip\noindent
First, we claim that it suffices to prove the lemma for $r = 2$.
We will argue this by induction on $r$; we suggest the reader
skip the proof. Suppose we are given $A \subset B$ and $\pi_j$ in
$B_j = B_{\mathfrak m_j}$ and factorizations
$a_i = u_{i, j} \pi_j^{e_{i, j}}$ for $i = 1, \ldots, r - 1$ in $B_j$
with $u_{i, j} \in B_j$ units and $e_{i, j} \geq 0$.
Then by the case $r = 2$ for $\pi_j$ and $a_r$ in $B_j$
we can find extensions $B_j \subset C_j$ and for every maximal ideal
$\mathfrak m_{j, k}$ of $C_j$ a nonzerodivisor
$\pi_{j, k} \in C_{j, k} = (C_j)_{\mathfrak m_{j, k}}$
and factorizations
$$
\pi_j = v_{j, k} \pi_{j, k}^{f_{j, k}}
\quad\text{and}\quad
a_r = w_{j, k} \pi_{j, k}^{g_{j, k}}
$$
as in the lemma. There exists a unique finite extension $B \subset C$
with $C/B$ annihilated by a power of $\mathfrak m$ such
that $C_j \cong C_{\mathfrak m_j}$ for all $j$, see
Lemma \ref{lemma-glue-at-max}.
The maximal ideals of $C$ correspond $1$-to-$1$
to the maximal ideals $\mathfrak m_{j, k}$ in the localizations
and in these localizations we have
$$
a_i = u_{i, j} \pi_j^{e_{i, j}} =
u_{i, j} v_{j, k}^{e_{i, j}} \pi_{j, k}^{e_{i, j}f_{j, k}}
$$
for $i \leq r - 1$. Since $a_r$ factors correctly too the
proof of the induction step is complete.

\medskip\noindent
Proof of the case $r = 2$. We will use induction on
$$
\ell = \min(\text{length}_A(A/a_1A),\ \text{length}_A(A/a_2A)).
$$
If $\ell = 0$, then either $a_1$ or $a_2$ is a unit and
the lemma holds with $A = B$. Thus we may and do assume $\ell > 0$.

\medskip\noindent
Suppose we have a finite extension of rings $A \subset A'$ such that
$A'/A$ is annihilated by a power of $\mathfrak m$ and such that
$\mathfrak m$ is not an associated prime of $A'$.
Let $\mathfrak m_1, \ldots, \mathfrak m_r \subset A'$
be the maximal ideals and set $A'_i = A'_{\mathfrak m_i}$.
If we can solve the problem for $a_1, a_2$ in each $A'_i$,
then we can apply Lemma \ref{lemma-glue-at-max}
to produce a solution for $a_1, a_2$ in $A$.
Choose $x \in \{a_1, a_2\}$ such that $\ell = \text{length}_A(A/xA)$.
By Lemma \ref{lemma-compare-periodic-lengths} 
and (\ref{equation-multiplicity-coker-ker})
we have $\text{length}_A(A/xA) = \text{length}_A(A'/xA')$.
On the other hand, we have
$$
\text{length}_A(A'/xA') =
\sum [\kappa(\mathfrak m_i) : \kappa(\mathfrak m)]
\text{length}_{A'_i}(A'_i/xA'_i)
$$
by Algebra, Lemma \ref{algebra-lemma-pushdown-module}.
Since $x \in \mathfrak m$ we see that each term on the right hand side
is positive. We conclude that the induction hypothesis applies
to $a_1, a_2$ in each $A'_i$ if $r > 1$ or if $r = 1$ and
$[\kappa(\mathfrak m_1) : \kappa(\mathfrak m)] > 1$.
We conclude that we may assume each $A'$ as above is local with
the same residue field as $A$.

\medskip\noindent
Applying the discussion of the previous paragraph,
we may replace $A$ by the ring constructed in
Lemma \ref{lemma-Noetherian-domain-dim-1-two-elements}
for $a_1, a_2 \in A$. Then since $A$ is local we find,
after possibly switching $a_1$ and $a_2$, that $a_2 \in (a_1)$.
Write $a_2 = a_1^m c$ with $m > 0$ maximal. In fact, by
Lemma \ref{lemma-not-infinitely-divisible}
we may assume $m$ is maximal even after replacing $A$
by any finite extension $A \subset A'$ as in the previous paragraph.
If $c$ is a unit, then we are done. If not, then we replace
$A$ by the ring constructed in
Lemma \ref{lemma-Noetherian-domain-dim-1-two-elements}
for $a_1, c \in A$. Then either (1) $c = a_1 c'$ or
(2) $a_1 = c a'_1$. The first case cannot happen since
it would give $a_2 = a_1^{m + 1} c'$ contradicting the
maximality of $m$. In the second case we get
$a_1 = c a'_1$ and $a_2 = c^{m + 1} (a'_1)^m$.
Then it suffices to prove the lemma for $A$ and $c, a'_1$.
If $a'_1$ is a unit we're done and if not, then
$\text{length}_A(A/cA) < \ell$ because $cA$ is a strictly
bigger ideal than $a_1A$. Thus we win by induction hypothesis.
\end{proof}






\section{Tame symbols}
\label{section-tame-symbol}

\noindent
Consider a Noetherian local ring $(A, \mathfrak m)$ of dimension $1$.
We denote $Q(A)$ the total ring of fractions of $A$, see
Algebra, Example \ref{algebra-example-localize-at-prime}.
The {\it tame symbol} will be a map
$$
\partial_A(-, -) : Q(A)^* \times Q(A)^* \longrightarrow \kappa(\mathfrak m)^*
$$
satisfying the following properties:
\begin{enumerate}
\item $\partial_A(f, gh) = \partial_A(f, g) \partial_A(f, h)$
\label{item-bilinear}
for $f, g, h \in Q(A)^*$,
\item $\partial_A(f, g) \partial_A(g, f) = 1$
\label{item-skew}
for $f, g \in Q(A)^*$,
\item $\partial_A(f, 1 - f) = 1$
\label{item-1-x}
for $f \in Q(A)^*$ such that $1 - f \in Q(A)^*$,
\item $\partial_A(aa', b) = \partial_A(a, b)\partial_A(a', b)$
\label{item-bilinear-better}
and $\partial_A(a, bb') = \partial_A(a, b)\partial_A(a, b')$
for $a, a', b, b' \in A$ nonzerodivisors,
\item $\partial_A(b, b) = (-1)^m$
\label{item-skew-better}
with $m = \text{length}_A(A/bA)$
for $b \in A$ a nonzerodivisor,
\item $\partial_A(u, b) = u^m \bmod \mathfrak m$
\label{item-normalization}
with $m = \text{length}_A(A/bA)$ for $u \in A$ a unit and
$b \in A$ a nonzerodivisor, and
\item
\label{item-1-x-better}
$\partial_A(a, b - a)\partial_A(b, b) = \partial_A(b, b - a)\partial_A(a, b)$
for $a, b \in A$ such that $a, b, b - a$ are nonzerodivisors.
\end{enumerate}
Since it is easier to work with elements of $A$ we will
often think of $\partial_A$ as a map defined on pairs of
nonzerodivisors of $A$ satisfying (\ref{item-bilinear-better}),
(\ref{item-skew-better}), (\ref{item-normalization}),
(\ref{item-1-x-better}). It is an exercise to see that
setting
$$
\partial_A(\frac{a}{b}, \frac{c}{d}) =
\partial_A(a, c) \partial_A(a, d)^{-1} \partial_A(b, c)^{-1} \partial_A(b, d)
$$
we get a well defined map $Q(A)^* \times Q(A)^* \to \kappa(\mathfrak m)^*$
satisfying (\ref{item-bilinear}), (\ref{item-skew}), (\ref{item-1-x})
as well as the other properties.

\medskip\noindent
We do not claim there is a unique map with these properties.
Instead, we will give a recipe for constructing such a map.
Namely, given $a_1, a_2 \in A$ nonzerodivisors, we choose
a ring extension $A \subset B$ and local factorizations
as in Lemma \ref{lemma-prepare-tame-symbol}.
Then we define
\begin{equation}
\label{equation-tame-symbol}
\partial_A(a_1, a_2) = \prod\nolimits_j
\text{Norm}_{\kappa(\mathfrak m_j)/\kappa(\mathfrak m)}
((-1)^{e_{1, j}e_{2, j}}u_{1, j}^{e_{2, j}}u_{2, j}^{-e_{1, j}}
\bmod \mathfrak m_j)^{m_j}
\end{equation}
where $m_j = \text{length}_{B_j}(B_j/\pi_j B_j)$ and the product
is taken over the maximal ideals $\mathfrak m_1, \ldots, \mathfrak m_r$ of $B$.

\begin{lemma}
\label{lemma-well-defined-tame-symbol}
The formula (\ref{equation-tame-symbol}) determines a
well defined element of $\kappa(\mathfrak m)^*$. In other words, the
right hand side does not depend on the choice of the
local factorizations or the choice of $B$.
\end{lemma}

\begin{proof}
Independence of choice of factorizations. Suppose we have
a Noetherian $1$-dimensional local ring $B$, elements $a_1, a_2 \in B$,
and nonzerodivisors $\pi, \theta$ such that we can write
$$
a_1 = u_1 \pi^{e_1} = v_1 \theta^{f_1},\quad
a_2 = u_2 \pi^{e_2} = v_2 \theta^{f_2}
$$
with $e_i, f_i \geq 0$ integers and $u_i, v_i$ units in $B$.
Observe that this implies
$$
a_1^{e_2} = u_1^{e_2}u_2^{-e_1}a_2^{e_1},\quad
a_1^{f_2} = v_1^{f_2}v_2^{-f_1}a_2^{f_1}
$$
On the other hand, setting
$m = \text{length}_B(B/\pi B)$ and $k = \text{length}_B(B/\theta B)$
we find $e_2 m = \text{length}_B(B/a_2 B) = f_2 k$.
Expanding $a_1^{e_2m} = a_1^{f_2 k}$ using the above we find
$$
(u_1^{e_2}u_2^{-e_1})^m =  (v_1^{f_2}v_2^{-f_1})^k
$$
This proves the desired equality up to signs. To see the signs
work out we have to show $me_1e_2$ is even if and only if
$kf_1f_2$ is even. This follows as both $me_2 = kf_2$ and
$me_1 = kf_1$ (same argument as above).

\medskip\noindent
Independence of choice of $B$. Suppose given two extensions
$A \subset B$ and $A \subset B'$ as in Lemma \ref{lemma-prepare-tame-symbol}.
Then
$$
C = (B \otimes_A B')/(\mathfrak m\text{-power torsion})
$$
will be a third one. Thus we may assume we have
$A \subset B \subset C$ and factorizations over the
local rings of $B$ and we have to show that using
the same factorizations over the local rings of $C$
gives the same element of $\kappa(\mathfrak m)$.
By transitivity of norms
(Fields, Lemma \ref{fields-lemma-trace-and-norm-tower})
this comes down to the following problem:
if $B$ is a Noetherian local ring of dimension $1$
and $\pi \in B$ is a nonzerodivisor, then
$$
\lambda^m = \prod \text{Norm}_{\kappa_k/\kappa}(\lambda)^{m_k}
$$
Here we have used the following notation:
(1) $\kappa$ is the residue field of $B$,
(2) $\lambda$ is an element of $\kappa$,
(3) $\mathfrak m_k \subset C$ are the maximal ideals of $C$,
(4) $\kappa_k = \kappa(\mathfrak m_k)$ is the residue field of
$C_k = C_{\mathfrak m_k}$,
(5) $m = \text{length}_B(B/\pi B)$, and
(6) $m_k = \text{length}_{C_k}(C_k/\pi C_k)$.
The displayed equality holds because
$\text{Norm}_{\kappa_k/\kappa}(\lambda) = \lambda^{[\kappa_k : \kappa]}$
as $\lambda \in \kappa$ and because $m = \sum m_k[\kappa_k:\kappa]$.
First, we have $m = \text{length}_B(B/xB) = \text{length}_B(C/\pi C)$
by Lemma \ref{lemma-compare-periodic-lengths} 
and (\ref{equation-multiplicity-coker-ker}).
Finally, we have $\text{length}_B(C/\pi C) = \sum m_k[\kappa_k:\kappa]$
by Algebra, Lemma \ref{algebra-lemma-pushdown-module}.
\end{proof}

\begin{lemma}
\label{lemma-tame-symbol}
The tame symbol (\ref{equation-tame-symbol}) satisfies
(\ref{item-bilinear-better}), (\ref{item-skew-better}),
(\ref{item-normalization}), (\ref{item-1-x-better}) and hence
gives a map $\partial_A : Q(A)^* \times Q(A)^* \to \kappa(\mathfrak m)^*$
satisfying (\ref{item-bilinear}), (\ref{item-skew}), (\ref{item-1-x}).
\end{lemma}

\begin{proof}
Let us prove (\ref{item-bilinear-better}).
Let $a_1, a_2, a_3 \in A$ be nonzerodivisors.
Choose $A \subset B$ as in Lemma \ref{lemma-prepare-tame-symbol}
for $a_1, a_2, a_3$. Then the equality
$$
\partial_A(a_1a_2, a_3) = \partial_A(a_1, a_3) \partial_A(a_2, a_3)
$$
follows from the equality
$$
(-1)^{(e_{1, j} + e_{2, j})e_{3, j}}
(u_{1, j}u_{2, j})^{e_{3, j}}u_{3, j}^{-e_{1, j} - e_{2, j}} =
(-1)^{e_{1, j}e_{3, j}}
u_{1, j}^{e_{3, j}}u_{3, j}^{-e_{1, j}}
(-1)^{e_{2, j}e_{3, j}}
u_{2, j}^{e_{3, j}}u_{3, j}^{-e_{2, j}}
$$
in $B_j$. Properties (\ref{item-skew-better}) and
(\ref{item-normalization}) are equally immediate.

\medskip\noindent
Let us prove (\ref{item-1-x-better}). Let $a_1, a_2, a_1 - a_2 \in A$
be nonzerodivisors and set $a_3 = a_1 - a_2$.
Choose $A \subset B$ as in Lemma \ref{lemma-prepare-tame-symbol}
for $a_1, a_2, a_3$. Then it suffices to show
$$
(-1)^{e_{1, j}e_{2, j} + e_{1, j}e_{3, j} + e_{2, j}e_{3, j} + e_{2, j}}
u_{1, j}^{e_{2, j} - e_{3, j}}
u_{2, j}^{e_{3, j} - e_{1, j}}
u_{3, j}^{e_{1, j} - e_{2, j}} \bmod \mathfrak m_j = 1
$$
This is clear if $e_{1, j} = e_{2, j} = e_{3, j}$.
Say $e_{1, j} > e_{2, j}$. Then we see that $e_{3, j} = e_{2, j}$
because $a_3 = a_1 - a_2$ and we see that $u_{3, j}$
has the same residue class as $-u_{2, j}$. Hence
the formula is true -- the signs work out as well
and this verification is the reason for the choice of signs
in (\ref{equation-tame-symbol}).
The other cases are handled in exactly the same manner.
\end{proof}

\begin{lemma}
\label{lemma-norm-down-tame-symbol}
Let $(A, \mathfrak m)$ be a Noetherian local ring of dimension $1$.
Let $A \subset B$ be a finite ring extension with $B/A$
annihilated by a power of $\mathfrak m$ and $\mathfrak m$ not
an associated prime of $B$.
For $a, b \in A$ nonzerodivisors we have
$$
\partial_A(a, b) = \prod
\text{Norm}_{\kappa(\mathfrak m_j)/\kappa(\mathfrak m)}(\partial_{B_j}(a, b))
$$
where the product is over the maximal ideals $\mathfrak m_j$ of $B$
and $B_j = B_{\mathfrak m_j}$.
\end{lemma}

\begin{proof}
Choose $B_j \subset C_j$ as in
Lemma \ref{lemma-prepare-tame-symbol} for $a, b$.
By Lemma \ref{lemma-glue-at-max} we can choose a finite ring
extension $B \subset C$ with $C_j \cong C_{\mathfrak m_j}$ for all $j$.
Let $\mathfrak m_{j, k} \subset C$ be the maximal ideals of $C$
lying over $\mathfrak m_j$. Let
$$
a = u_{j, k}\pi_{j, k}^{f_{j, k}},\quad
b = v_{j, k}\pi_{j, k}^{g_{j, k}}
$$
be the local factorizations which exist by our choice of
$C_j \cong C_{\mathfrak m_j}$. By definition we have
$$
\partial_A(a, b) = 
\prod\nolimits_{j, k}
\text{Norm}_{\kappa(\mathfrak m_{j, k})/\kappa(\mathfrak m)}
((-1)^{f_{j, k}g_{j, k}}u_{j, k}^{g_{j, k}}v_{j, k}^{-f_{j, k}}
\bmod \mathfrak m_{j, k})^{m_{j, k}}
$$
and
$$
\partial_{B_j}(a, b) = 
\prod\nolimits_k
\text{Norm}_{\kappa(\mathfrak m_{j, k})/\kappa(\mathfrak m_j)}
((-1)^{f_{j, k}g_{j, k}}u_{j, k}^{g_{j, k}}v_{j, k}^{-f_{j, k}}
\bmod \mathfrak m_{j, k})^{m_{j, k}}
$$
The result follows by transitivity of norms
for $\kappa(\mathfrak m_{j, k})/\kappa(\mathfrak m_j)/\kappa(\mathfrak m)$, see
Fields, Lemma \ref{fields-lemma-trace-and-norm-tower}.
\end{proof}

\begin{lemma}
\label{lemma-tame-symbol-formally-smooth}
Let $(A, \mathfrak m, \kappa) \to (A', \mathfrak m', \kappa')$
be a local homomorphism of Noetherian local rings. Assume $A \to A'$
is flat and $\dim(A) = \dim(A') = 1$. Set
$m = \text{length}_{A'}(A'/\mathfrak mA')$.
For $a_1, a_2 \in A$ nonzerodivisors
$\partial_A(a_1, a_2)^m$ maps to $\partial_{A'}(a_1, a_2)$
via $\kappa \to \kappa'$.
\end{lemma}

\begin{proof}
If $a_1, a_2$ are both units, then $\partial_A(a_1, a_2) = 1$
and $\partial_{A'}(a_1, a_2) = 1$ and the result is true.
If not, then we can choose a ring extension $A \subset B$ and
local factorizations as in Lemma \ref{lemma-prepare-tame-symbol}.
Denote $\mathfrak m_1, \ldots, \mathfrak m_m$
be the maximal ideals of $B$. Let $\mathfrak m_1, \ldots, \mathfrak m_m$
be the maximal ideals of $B$ with residue fields $\kappa_1, \ldots, \kappa_m$.
For each $j \in \{1, \ldots, m\}$ denote $\pi_j \in B_j = B_{\mathfrak m_j}$
a nonzerodivisor such that we have factorizations
$a_i = u_{i, j}\pi_j^{e_{i, j}}$ as in the lemma.
By definition we have
$$
\partial_A(a_1, a_2) = \prod\nolimits_j
\text{Norm}_{\kappa_j/\kappa}
((-1)^{e_{1, j}e_{2, j}}u_{1, j}^{e_{2, j}}u_{2, j}^{-e_{1, j}}
\bmod \mathfrak m_j)^{m_j}
$$
where $m_j = \text{length}_{B_j}(B_j/\pi_j B_j)$.

\medskip\noindent
Set $B' = A' \otimes_A B$. Since $A'$ is flat over $A$ we see
that $A' \subset B'$ is a ring extension with $B'/A'$ annihilated
by a power of $\mathfrak m'$. Let
$$
\mathfrak m'_{j, l},\quad l = 1, \ldots, n_j
$$
be the maximal ideals of $B'$ lying over $\mathfrak m_j$. Denote
$\kappa'_{j, l}$ the residue field of $\mathfrak m'_{j, l}$. Denote
$B'_{j, l}$ the localization of $B'$ at $\mathfrak m'_{j, l}$.
As factorizations of $a_1$ and $a_2$ in $B'_{j, l}$
we use the image of the factorizations
$a_i = u_{i, j} \pi_j^{e_{i, j}}$ given to us in $B_j$.
By definition we have
$$
\partial_{A'}(a_1, a_2) = \prod\nolimits_{j, l}
\text{Norm}_{\kappa'_{j, l}/\kappa'}
((-1)^{e_{1, j}e_{2, j}}u_{1, j}^{e_{2, j}}u_{2, j}^{-e_{1, j}}
\bmod \mathfrak m'_{j, l})^{m'_{j, l}}
$$
where $m'_{j, l} = \text{length}_{B'_{j, l}}(B'_{j, l}/\pi_j B'_{j, l})$.

\medskip\noindent
Comparing the formulae above we see that it suffices to show that
for each $j$ and for any unit $u \in B_j$ we have
\begin{equation}
\label{equation-to-prove}
\left(\text{Norm}_{\kappa_j/\kappa}(u \bmod \mathfrak m_j)^{m_j}\right)^m
=
\prod\nolimits_l
\text{Norm}_{\kappa'_{j, l}/\kappa'}(u \bmod \mathfrak m'_{j, l})^{m'_{j, l}}
\end{equation}
in $\kappa'$. We are going to use the construction of determinants of
endomorphisms of finite length modules in
More on Algebra, Section \ref{more-algebra-section-determinants-finite-length}
to prove this. Set $M = B_j/\pi_j B_j$. By
More on Algebra, Lemma \ref{more-algebra-lemma-multiplication} we have
$$
\text{Norm}_{\kappa_j/\kappa}(u \bmod \mathfrak m_j)^{m_j} =
\det\nolimits_\kappa(u : M \to M)
$$
Thus, by
More on Algebra, Lemma \ref{more-algebra-lemma-flat-base-change-det},
the left hand side of (\ref{equation-to-prove}) is equal to
$\det_{\kappa'}(u : M \otimes_A A' \to M \otimes_A A')$.
We have an isomorphism
$$
M \otimes_A A' = (B_j/\pi_j B_j) \otimes_A A' =
\bigoplus\nolimits_l B'_{j, l}/\pi_j B'_{j, l}
$$
of $A'$-modules. Setting $M'_l = B'_{j, l}/\pi_j B'_{j, l}$ we see that
$\text{Norm}_{\kappa'_{j, l}/\kappa'}(u \bmod \mathfrak m'_{j, l})^{m'_{j, l}}
= \det_{\kappa'}(u_j : M'_l \to M'_l)$ by
More on Algebra, Lemma \ref{more-algebra-lemma-multiplication} again.
Hence (\ref{equation-to-prove}) holds by multiplicativity of the determinant
construction, see More on Algebra, Lemma \ref{more-algebra-lemma-ses}.
\end{proof}






\section{A key lemma}
\label{section-key-lemma}

\noindent
In this section we apply the results above to prove
Lemma \ref{lemma-milnor-gersten-low-degree}.
This lemma is a low degree case of the statement
that there is a complex for Milnor K-theory similar
to the Gersten-Quillen complex in Quillen's K-theory.
See Remark \ref{remark-gersten-complex-milnor}.

\begin{lemma}
\label{lemma-perpare-key}
Let $(A, \mathfrak m)$ be a $2$-dimensional Noetherian local ring.
Let $t \in \mathfrak m$ be a nonzerodivisor. Say
$V(t) = \{\mathfrak m, \mathfrak q_1, \ldots, \mathfrak q_r\}$.
Let $A_{\mathfrak q_i} \subset B_i$ be a finite ring
extension with $B_i/A_{\mathfrak q_i}$ annihilated by a power of
$t$. Then there exists a finite extension $A \subset B$ of
local rings identifying residue fields
with $B_i \cong B_{\mathfrak q_i}$ and $B/A$ annihilated
by a power of $t$.
\end{lemma}

\begin{proof}
Choose $n > 0$ such that $B_i \subset t^{-n}A_{\mathfrak q_i}$.
Let $M \subset t^{-n}A$, resp.\ $M' \subset t^{-2n}A$ be the
$A$-submodule consisting of elements mapping to $B_i$ in
$t^{-n}A_{\mathfrak q_i}$, resp.\ $t^{-2n}A_{\mathfrak q_i}$.
Then $M \subset M'$ are finite $A$-modules as $A$ is Noetherian
and $M_{\mathfrak q_i} = M'_{\mathfrak q_i} = B_i$ as localization
is exact. Thus $M'/M$ is annihilated by $\mathfrak m^c$ for some
$c > 0$. Observe that $M \cdot M \subset M'$ under the multiplication
$t^{-n}A \times t^{-n}A \to t^{-2n}A$. Hence
$B = A + \mathfrak m^{c + 1}M$ is a finite $A$-algebra with the correct
localizations. We omit the verification that $B$ is local with
maximal ideal $\mathfrak m + \mathfrak m^{c + 1}M$.
\end{proof}

\begin{lemma}
\label{lemma-key-nonzerodivisors}
Let $(A, \mathfrak m)$ be a $2$-dimensional Noetherian local ring.
Let $a, b \in A$ be nonzerodivisors.
Then we have
$$
\sum
\text{ord}_{A/\mathfrak q}(\partial_{A_{\mathfrak q}}(a, b))
=
0
$$
where the sum is over the height $1$ primes $\mathfrak q$ of $A$.
\end{lemma}

\begin{proof}
If $\mathfrak q$ is a height $1$ prime of $A$ such that $a, b$
map to a unit of $A_\mathfrak q$, then $\partial_{A_\mathfrak q}(a, b) = 1$.
Thus the sum is finite. In fact, if
$V(ab) = \{\mathfrak m, \mathfrak q_1, \ldots, \mathfrak q_r\}$
then the sum is over $i = 1, \ldots, r$.
For each $i$ we pick an extension $A_{\mathfrak q_i} \subset B_i$
as in Lemma \ref{lemma-prepare-tame-symbol} for $a, b$.
By Lemma \ref{lemma-perpare-key} with $t = ab$ and the given list of primes
we may assume we have a finite local extension $A \subset B$
with $B/A$ annihilated by a power of $ab$ and such that
for each $i$ the $B_{\mathfrak q_i} \cong B_i$.
Observe that if $\mathfrak q_{i, j}$ are the primes of $B$
lying over $\mathfrak q_i$ then we have
$$
\text{ord}_{A/\mathfrak q_i}(\partial_{A_{\mathfrak q_i}}(a, b))
=
\sum\nolimits_j
\text{ord}_{B/\mathfrak q_{i, j}}(\partial_{B_{\mathfrak q_{i, j}}}(a, b))
$$
by Lemma \ref{lemma-norm-down-tame-symbol} and
Algebra, Lemma \ref{algebra-lemma-finite-extension-dim-1}.
Thus we may replace $A$ by $B$ and
reduce to the case discussed in the next paragraph.

\medskip\noindent
Assume for each $i$ there is a nonzerodivisor
$\pi_i \in A_{\mathfrak q_i}$ and units $u_i, v_i \in A_{\mathfrak q_i}$
such that for some integers $e_i, f_i \geq 0$ we have
$$
a = u_i \pi_i^{e_i},\quad b = v_i \pi_i^{f_i}
$$
in $A_{\mathfrak q_i}$. Setting
$m_i = \text{length}_{A_{\mathfrak q_i}}(A_{\mathfrak q_i}/\pi_i)$
we have
$\partial_{A_{\mathfrak q_i}}(a, b) =
((-1)^{e_if_i}u_i^{f_i}v_i^{-e_i})^{m_i}$ by definition.
Since $a, b$ are nonzerodivisors the
$(2, 1)$-periodic complex $(A/(ab), a, b)$ has vanishing cohomology.
Denote $M_i$ the image of $A/(ab)$ in $A_{\mathfrak q_i}/(ab)$.
Then we have a map
$$
A/(ab) \longrightarrow \bigoplus M_i
$$
whose kernel and cokernel are supported in $\{\mathfrak m\}$
and hence have finite length. Thus we see that
$$
\sum e_A(M_i, a, b) = 0
$$
by Lemma \ref{lemma-compare-periodic-lengths}. Hence it suffices to show
$e_A(M_i, a, b) =
- \text{ord}_{A/\mathfrak q_i}(\partial_{A_{\mathfrak q_i}}(a, b))$.

\medskip\noindent
Let us prove this first, in case $\pi_i, u_i, v_i$ are the images
of elements $\pi_i, u_i, v_i \in A$ (using the same symbols
should not cause any confusion). In this case we get
\begin{align*}
e_A(M_i, a, b) & =
e_A(M_i, u_i\pi_i^{e_i}, v_i\pi_i^{f_i}) \\
& =
e_A(M_i, \pi_i^{e_i}, \pi_i^{f_i}) -
e_A(\pi_i^{e_i}M_i, 0, u_i) +
e_A(\pi_i^{f_i}M_i, 0, v_i) \\
& =
0 -
f_im_i\text{ord}_{A/\mathfrak q_i}(u_i) +
e_im_i\text{ord}_{A/\mathfrak q_i}(v_i) \\
& =
-m_i\text{ord}_{A/\mathfrak q_i}(u_i^{f_i}v_i^{-e_i}) =
-\text{ord}_{A/\mathfrak q_i}(\partial_{A_{\mathfrak q_i}}(a, b))
\end{align*}
The second equality holds by Lemma \ref{lemma-multiply-period-length}.
Observe that
$M_i \subset (M_i)_{\mathfrak q_i} = A_{\mathfrak q_i}/(\pi_i^{e_i + f_i})$
and
$(\pi_i^{e_i}M_i)_{\mathfrak q_i} \cong A_{\mathfrak q_i}/\pi_i^{f_i}$ and
$(\pi_i^{f_i}M_i)_{\mathfrak q_i} \cong A_{\mathfrak q_i}/\pi_i^{e_i}$.
The $0$ in the third equality comes from
Lemma \ref{lemma-powers-period-length-zero}
and the other two terms come from
Lemma \ref{lemma-length-multiplication}.
The last two equalities follow from multiplicativity of
the order function and from the definition of our tame symbol.

\medskip\noindent
In general, we may first choose $c \in A$, $c \not \in \mathfrak q_i$
such that $c\pi_i \in A$. After replacing $\pi_i$ by $c\pi_i$
and $u_i$ by $c^{-e_i}u_i$ and $v_i$ by $c^{-f_i}v_i$
we may and do assume $\pi_i$ is in $A$.
Next, choose an $c \in A$, $c \not \in \mathfrak q_i$
with $cu_i, cv_i \in A$. Then we observe that
$$
e_A(M_i, ca, cb) = e_A(M_i, a, b) - e_A(aM_i, 0, c) + e_A(bM_i, 0, c)
$$
by Lemma \ref{lemma-length-multiplication}.
On the other hand, we have
$$
\partial_{A_{\mathfrak q_i}}(ca, cb) =
c^{m_i(f_i - e_i)}\partial_{A_{\mathfrak q_i}}(a, b)
$$
in $\kappa(\mathfrak q_i)^*$ because $c$ is a unit in $A_{\mathfrak q_i}$.
The arguments in the previous paragraph show that
$e_A(M_i, ca, cb) = -
\text{ord}_{A/\mathfrak q_i}(\partial_{A_{\mathfrak q_i}}(ca, cb))$.
Thus it suffices to prove
$$
e_A(aM_i, 0, c) = \text{ord}_{A/\mathfrak q_i}(c^{m_if_i})
\quad\text{and}\quad
e_A(bM_i, 0, c) = \text{ord}_{A/\mathfrak q_i}(c^{m_ie_i})
$$
and this follows from Lemma \ref{lemma-length-multiplication}
by the description (see above)
of what happens when we localize at $\mathfrak q_i$.
\end{proof}

\begin{lemma}[Key Lemma]
\label{lemma-milnor-gersten-low-degree}
\begin{reference}
When $A$ is an excellent ring this is \cite[Proposition 1]{Kato-Milnor-K}.
\end{reference}
Let $A$ be a $2$-dimensional Noetherian local domain with fraction field $K$.
Let $f, g \in K^*$.
Let $\mathfrak q_1, \ldots, \mathfrak q_t$ be the height
$1$ primes $\mathfrak q$ of $A$ such that either $f$ or $g$ is not an
element of $A^*_{\mathfrak q}$.
Then we have
$$
\sum\nolimits_{i = 1, \ldots, t}
\text{ord}_{A/\mathfrak q_i}(\partial_{A_{\mathfrak q_i}}(f, g))
=
0
$$
We can also write this as
$$
\sum\nolimits_{\text{height}(\mathfrak q) = 1}
\text{ord}_{A/\mathfrak q}(\partial_{A_{\mathfrak q}}(f, g))
=
0
$$
since at any height $1$ prime $\mathfrak q$
of $A$ where $f, g \in A^*_{\mathfrak q}$
we have $\partial_{A_{\mathfrak q}}(f, g) = 1$.
\end{lemma}

\begin{proof}
Since the tame symbols $\partial_{A_{\mathfrak q}}(f, g)$ are
bilinear and the order functions $\text{ord}_{A/\mathfrak q}$
are additive it suffices to prove the formula when
$f$ and $g$ are elements of $A$. This case is proven in
Lemma \ref{lemma-key-nonzerodivisors}.
\end{proof}

\begin{remark}[Milnor K-theory]
\label{remark-gersten-complex-milnor}
For a field $k$ let us denote $K^M_*(k)$ the quotient of
the tensor algebra on $k^*$ divided by the two-sided ideal
generated by the elements $x \otimes 1 - x$ for $x \in k \setminus \{0, 1\}$.
Thus $K^M_0(k) = \mathbf{Z}$, $K_1^M(k) = k^*$, and
$$
K^M_2(k) = k^* \otimes_\mathbf{Z} k^* / \langle x \otimes 1 - x \rangle
$$
If $A$ is a discrete valuation ring with fraction field $F = \text{Frac}(A)$
and residue field $\kappa$, there is a tame symbol
$$
\partial_A : K_{i + 1}^M(F) \to K_i^M(\kappa)
$$
defined as in Section \ref{section-tame-symbol}; see \cite{Kato-Milnor-K}.
More generally, this map can be extended to the case where $A$ is an
excellent local domain of dimension $1$ using normalization and norm
maps on $K_i^M$, see \cite{Kato-Milnor-K}; presumably the method in
Section \ref{section-tame-symbol} can be used to extend the construction
of the tame symbol $\partial_A$ to arbitrary Noetherian local domains $A$
of dimension $1$. Next, let $X$ be a Noetherian scheme with a
dimension function $\delta$. Then we can use these tame symbols to get
the arrows in the following:
$$
\bigoplus\nolimits_{\delta(x) = j + 1} K^M_{i + 1}(\kappa(x))
\longrightarrow
\bigoplus\nolimits_{\delta(x) = j} K^M_i(\kappa(x))
\longrightarrow
\bigoplus\nolimits_{\delta(x) = j - 1} K^M_{i - 1}(\kappa(x))
$$
However, it is not clear, that the composition is zero, i.e., that
we obtain a complex of abelian groups.
For excellent $X$ this is shown in \cite{Kato-Milnor-K}.
When $i = 1$ and $j$ arbitrary, this follows from
Lemma \ref{lemma-milnor-gersten-low-degree}.
\end{remark}





















\section{Setup}
\label{section-setup}

\noindent
We will throughout work over a locally Noetherian universally
catenary base $S$ endowed with a dimension function $\delta$.
Although it is likely possible to generalize (parts of) the
discussion in the chapter, it seems that this is a good first
approximation. It is exactly the generality discussed in \cite{Thorup}.
We usually do not assume our schemes are
separated or quasi-compact. Many interesting algebraic stacks
are non-separated and/or non-quasi-compact and this is a good
case study to see how to develop a reasonable theory for those as well.
In order to reference these hypotheses we give it a number.

\begin{situation}
\label{situation-setup}
Here $S$ is a locally Noetherian, and universally catenary scheme.
Moreover, we assume $S$ is endowed with a dimension function
$\delta : S \longrightarrow \mathbf{Z}$.
\end{situation}

\noindent
See Morphisms, Definition \ref{morphisms-definition-universally-catenary}
for the notion of a universally catenary scheme, and see
Topology, Definition \ref{topology-definition-dimension-function}
for the notion of a dimension function. Recall that any locally
Noetherian catenary scheme locally has a dimension function, see
Properties, Lemma \ref{properties-lemma-catenary-dimension-function}.
Moreover, there are lots of schemes which are universally catenary,
see Morphisms, Lemma \ref{morphisms-lemma-ubiquity-uc}.

\medskip\noindent
Let $(S, \delta)$ be as in Situation \ref{situation-setup}.
Any scheme $X$ locally of finite type over $S$ is locally Noetherian
and catenary. In fact, $X$ has a canonical dimension function
$$
\delta = \delta_{X/S} : X \longrightarrow \mathbf{Z}
$$
associated to $(f : X \to S, \delta)$ given by the rule
$\delta_{X/S}(x) = \delta(f(x)) + \text{trdeg}_{\kappa(f(x))}\kappa(x)$.
See Morphisms, Lemma \ref{morphisms-lemma-dimension-function-propagates}.
Moreover, if $h : X \to Y$ is a morphism of schemes locally of finite
type over $S$, and $x \in X$, $y = h(x)$,
then obviously
$\delta_{X/S}(x) = \delta_{Y/S}(y) + \text{trdeg}_{\kappa(y)}\kappa(x)$.
We will freely use this function and its properties in the following.

\medskip\noindent
Here are the basic examples of setups as above.
In fact, the main interest lies in the case where the base
is the spectrum of a field, or the case where the base
is the spectrum of a Dedekind ring (e.g.\ $\mathbf{Z}$,
or a discrete valuation ring).

\begin{example}
\label{example-field}
Here $S = \Spec(k)$ and $k$ is a field.
We set $\delta(pt) = 0$ where $pt$ indicates the unique point of $S$.
The pair $(S, \delta)$ is an example of a situation as in
Situation \ref{situation-setup} by
Morphisms, Lemma \ref{morphisms-lemma-ubiquity-uc}.
\end{example}

\begin{example}
\label{example-domain-dimension-1}
Here $S = \Spec(A)$, where $A$ is a Noetherian domain
of dimension $1$.
For example we could consider $A = \mathbf{Z}$.
We set $\delta(\mathfrak p) = 0$ if
$\mathfrak p$ is a maximal ideal and $\delta(\mathfrak p) = 1$
if $\mathfrak p = (0)$ corresponds to the generic point.
This is an example of Situation \ref{situation-setup} by
Morphisms, Lemma \ref{morphisms-lemma-ubiquity-uc}.
\end{example}

\begin{example}
\label{example-CM-irreducible}
Here $S$ is a Cohen-Macaulay scheme. Then $S$ is universally catenary by
Morphisms, Lemma \ref{morphisms-lemma-ubiquity-uc}.
We set $\delta(s) = -\dim(\mathcal{O}_{S, s})$.
If $s' \leadsto s$ is a nontrivial specialization of points of $S$,
then $\mathcal{O}_{S, s'}$ is the localization of $\mathcal{O}_{S, s}$
at a nonmaximal prime ideal $\mathfrak p \subset \mathcal{O}_{S, s}$, see
Schemes, Lemma \ref{schemes-lemma-specialize-points}.
Thus $\dim(\mathcal{O}_{S, s}) =
\dim(\mathcal{O}_{S, s'}) + \dim(\mathcal{O}_{S, s}/\mathfrak p) >
\dim(\mathcal{O}_{S, s'})$ by
Algebra, Lemma \ref{algebra-lemma-CM-dim-formula}.
Hence $\delta(s') > \delta(s)$. If $s' \leadsto s$ is
an immediate specialization, then there is no prime
ideal strictly between $\mathfrak p$ and $\mathfrak m_s$
and we find $\delta(s') = \delta(s) + 1$. Thus $\delta$
is a dimension function. In other words, the pair $(S, \delta)$
is an example of Situation \ref{situation-setup}.
\end{example}

\noindent
If $S$ is Jacobson and $\delta$ sends closed points to zero, then $\delta$
is the function sending a point to the dimension of its closure.

\begin{lemma}
\label{lemma-delta-is-dimension}
Let $(S, \delta)$ be as in Situation \ref{situation-setup}.
Assume in addition $S$ is a Jacobson scheme, and $\delta(s) = 0$ for every
closed point $s$ of $S$. Let $X$ be locally of finite type over $S$.
Let $Z \subset X$ be an integral closed subscheme and let
$\xi \in Z$ be its generic point. The following integers are the same:
\begin{enumerate}
\item $\delta_{X/S}(\xi)$,
\item $\dim(Z)$, and
\item $\dim(\mathcal{O}_{Z, z})$ where $z$ is a closed point of $Z$.
\end{enumerate}
\end{lemma}

\begin{proof}
Let $X \to S$, $\xi \in Z \subset X$ be as in the lemma.
Since $X$ is locally of finite type over $S$ we see that
$X$ is Jacobson, see
Morphisms, Lemma \ref{morphisms-lemma-Jacobson-universally-Jacobson}.
Hence closed points of $X$ are dense in every closed subset of $Z$
and map to closed points of $S$. Hence given any chain
of irreducible closed subsets of $Z$ we can end it with a closed point of $Z$.
It follows that $\dim(Z) = \sup_z(\dim(\mathcal{O}_{Z, z})$
(see Properties, Lemma \ref{properties-lemma-codimension-local-ring})
where $z \in Z$ runs over the closed points of $Z$.
Note that $\dim(\mathcal{O}_{Z, z}) = \delta(\xi) - \delta(z)$
by the properties of a dimension function.
For each closed $z \in Z$ the field extension
$\kappa(z)/\kappa(f(z))$ is finite, see Morphisms,
Lemma \ref{morphisms-lemma-jacobson-finite-type-points}.
Hence $\delta_{X/S}(z) = \delta(f(z)) = 0$ for $z \in Z$ closed.
It follows that all three integers are equal.
\end{proof}

\noindent
In the situation of the lemma above the
value of $\delta$ at the generic point of a closed irreducible subset
is the dimension of the irreducible closed subset.
However, in general we cannot expect the equality to hold.
For example if $S = \Spec(\mathbf{C}[[t]])$ and
$X = \Spec(\mathbf{C}((t)))$ then we would get
$\delta(x) = 1$ for the unique point of $X$, but $\dim(X) = 0$.
Still we want to think of $\delta_{X/S}$ as giving the
dimension of the irreducible closed subschemes. Thus we introduce
the following terminology.

\begin{definition}
\label{definition-delta-dimension}
Let $(S, \delta)$ as in Situation \ref{situation-setup}.
For any scheme $X$ locally of finite type over $S$
and any irreducible closed subset $Z \subset X$ we define
$$
\dim_\delta(Z) = \delta(\xi)
$$
where $\xi \in Z$ is the generic point of $Z$.
We will call this the {\it $\delta$-dimension of $Z$}.
If $Z$ is a closed subscheme of $X$, then we define
$\dim_\delta(Z)$ as the supremum of the $\delta$-dimensions
of its irreducible components.
\end{definition}







\section{Cycles}
\label{section-cycles}

\noindent
Since we are not assuming our schemes are quasi-compact we have
to be a little careful when defining cycles. We have to allow
infinite sums because a rational function may have infinitely many
poles for example. In any case, if $X$ is quasi-compact then a
cycle is a finite sum as usual.

\begin{definition}
\label{definition-cycles}
Let $(S, \delta)$ be as in Situation \ref{situation-setup}.
Let $X$ be locally of finite type over $S$.
Let $k \in \mathbf{Z}$.
\begin{enumerate}
\item A {\it cycle on $X$} is a formal sum
$$
\alpha = \sum n_Z [Z]
$$
where the sum is over integral closed subschemes $Z \subset X$,
each $n_Z \in \mathbf{Z}$, and the collection
$\{Z; n_Z \not = 0\}$ is locally finite
(Topology, Definition \ref{topology-definition-locally-finite}).
\item A {\it $k$-cycle} on $X$ is a cycle
$$
\alpha = \sum n_Z [Z]
$$
where $n_Z \not = 0 \Rightarrow \dim_\delta(Z) = k$.
\item The abelian group of all $k$-cycles on $X$ is denoted $Z_k(X)$.
\end{enumerate}
\end{definition}

\noindent
In other words, a $k$-cycle on $X$
is a locally finite formal $\mathbf{Z}$-linear
combination of integral closed subschemes of $\delta$-dimension $k$.
Addition of $k$-cycles $\alpha = \sum n_Z[Z]$ and
$\beta = \sum m_Z[Z]$ is given by
$$
\alpha + \beta = \sum (n_Z + m_Z)[Z],
$$
i.e., by adding the coefficients.

\begin{remark}
\label{remark-cycles-pointwise}
Let $(S, \delta)$ be as in Situation \ref{situation-setup}.
Let $X$ be locally of finite type over $S$. Let $k \in \mathbf{Z}$.
Then we can write
$$
Z_k(X) = \bigoplus\nolimits_{\delta(x) = k}' K_0^M(\kappa(x))
\quad\subset\quad
\bigoplus\nolimits_{\delta(x) = k} K_0^M(\kappa(x))
$$
with the following notation and conventions:
\begin{enumerate}
\item $K_0^M(\kappa(x)) = \mathbf{Z}$ is the degree $0$ part of
the Milnor K-theory of the residue field $\kappa(x)$ of the point
$x \in X$ (see Remark \ref{remark-gersten-complex-milnor}), and
\item the direct sum on the right is over all points $x \in X$
with $\delta(x) = k$,
\item the notation $\bigoplus'_x$ signifies that we consider the
subgroup consisting of locally finite elements; namely, elements
$\sum_x n_x$ such that for every quasi-compact open $U \subset X$
the set of $x \in U$ with $n_x \not = 0$ is finite.
\end{enumerate}
\end{remark}

\begin{definition}
\label{definition-support-cycle}
Let $(S, \delta)$ be as in Situation \ref{situation-setup}.
Let $X$ be locally of finite type over $S$.
The {\it support} of a cycle $\alpha = \sum n_Z [Z]$ on $X$
is
$$
\text{Supp}(\alpha) = \bigcup\nolimits_{n_Z \not = 0} Z \subset X
$$
\end{definition}

\noindent
Since the collection $\{Z; n_Z \not = 0\}$ is locally finite
we see that $\text{Supp}(\alpha)$ is a closed subset of $X$.
If $\alpha$ is a $k$-cycle, then every irreducible component $Z$
of $\text{Supp}(\alpha)$ has $\delta$-dimension $k$.

\begin{definition}
\label{definition-effective-cycle}
Let $(S, \delta)$ be as in Situation \ref{situation-setup}.
Let $X$ be locally of finite type over $S$.
A cycle $\alpha$ on $X$ is {\it effective} if it
can be written as $\alpha =\sum n_Z [Z]$ with $n_Z \geq 0$ for all $Z$.
\end{definition}

\noindent
The set of all effective cycles is a monoid because the sum of
two effective cycles is effective, but it is not a group
(unless $X = \emptyset$).




\section{Cycle associated to a closed subscheme}
\label{section-cycle-of-closed-subscheme}

\begin{lemma}
\label{lemma-multiplicity-finite}
Let $(S, \delta)$ be as in Situation \ref{situation-setup}.
Let $X$ be locally of finite type over $S$.
Let $Z \subset X$ be a closed subscheme.
\begin{enumerate}
\item Let $Z' \subset Z$ be an irreducible component and
let $\xi \in Z'$ be its generic point.
Then
$$
\text{length}_{\mathcal{O}_{X, \xi}} \mathcal{O}_{Z, \xi} < \infty
$$
\item If $\dim_\delta(Z) \leq k$ and $\xi \in Z$ with
$\delta(\xi) = k$, then $\xi$ is a generic point of an
irreducible component of $Z$.
\end{enumerate}
\end{lemma}

\begin{proof}
Let $Z' \subset Z$, $\xi \in Z'$ be as in (1).
Then $\dim(\mathcal{O}_{Z, \xi}) = 0$ (for example by
Properties, Lemma \ref{properties-lemma-codimension-local-ring}).
Hence $\mathcal{O}_{Z, \xi}$ is Noetherian
local ring of dimension zero, and hence has finite length over
itself (see
Algebra, Proposition \ref{algebra-proposition-dimension-zero-ring}).
Hence, it also has finite length over $\mathcal{O}_{X, \xi}$, see
Algebra, Lemma \ref{algebra-lemma-length-independent}.

\medskip\noindent
Assume $\xi \in Z$ and $\delta(\xi) = k$.
Consider the closure $Z' = \overline{\{\xi\}}$. It is an irreducible
closed subscheme with $\dim_\delta(Z') = k$ by definition.
Since $\dim_\delta(Z) = k$ it must be an irreducible component
of $Z$. Hence we see (2) holds.
\end{proof}

\begin{definition}
\label{definition-cycle-associated-to-closed-subscheme}
Let $(S, \delta)$ be as in Situation \ref{situation-setup}.
Let $X$ be locally of finite type over $S$.
Let $Z \subset X$ be a closed subscheme.
\begin{enumerate}
\item For any irreducible component $Z' \subset Z$ with generic point $\xi$
the integer
$m_{Z', Z} = \text{length}_{\mathcal{O}_{X, \xi}} \mathcal{O}_{Z, \xi}$
(Lemma \ref{lemma-multiplicity-finite})
is called the {\it multiplicity of $Z'$ in $Z$}.
\item Assume $\dim_\delta(Z) \leq k$.
The {\it $k$-cycle associated to $Z$} is
$$
[Z]_k
=
\sum m_{Z', Z}[Z']
$$
where the sum is over the irreducible components of $Z$
of $\delta$-dimension $k$. (This is a $k$-cycle by
Divisors, Lemma \ref{divisors-lemma-components-locally-finite}.)
\end{enumerate}
\end{definition}

\noindent
It is important to note that we only define $[Z]_k$ if the $\delta$-dimension
of $Z$ does not exceed $k$. In other words, by convention, if we write
$[Z]_k$ then this implies that $\dim_\delta(Z) \leq k$.



\section{Cycle associated to a coherent sheaf}
\label{section-cycle-of-coherent-sheaf}



\begin{lemma}
\label{lemma-length-finite}
Let $(S, \delta)$ be as in Situation \ref{situation-setup}.
Let $X$ be locally of finite type over $S$.
Let $\mathcal{F}$ be a coherent $\mathcal{O}_X$-module.
\begin{enumerate}
\item The collection of irreducible components of the support of
$\mathcal{F}$ is locally finite.
\item Let $Z' \subset \text{Supp}(\mathcal{F})$
be an irreducible component and
let $\xi \in Z'$ be its generic point.
Then
$$
\text{length}_{\mathcal{O}_{X, \xi}} \mathcal{F}_\xi < \infty
$$
\item If $\dim_\delta(\text{Supp}(\mathcal{F})) \leq k$
and $\xi \in \text{Supp}(\mathcal{F})$ with $\delta(\xi) = k$, then $\xi$ is a
generic point of an irreducible component of $\text{Supp}(\mathcal{F})$.
\end{enumerate}
\end{lemma}

\begin{proof}
By Cohomology of Schemes, Lemma \ref{coherent-lemma-coherent-support-closed}
the support $Z$ of $\mathcal{F}$ is a closed subset of $X$.
We may think of $Z$ as a reduced closed subscheme of $X$
(Schemes, Lemma \ref{schemes-lemma-reduced-closed-subscheme}).
Hence (1) follows from
Divisors, Lemma \ref{divisors-lemma-components-locally-finite} applied to $Z$
and (3) follows from
Lemma \ref{lemma-multiplicity-finite} applied to $Z$.

\medskip\noindent
Let $\xi \in Z'$ be as in (2). In this case for any specialization
$\xi' \leadsto \xi$ in $X$ we have $\mathcal{F}_{\xi'} = 0$.
Recall that the non-maximal primes of $\mathcal{O}_{X, \xi}$ correspond
to the points of $X$ specializing to $\xi$
(Schemes, Lemma \ref{schemes-lemma-specialize-points}).
Hence $\mathcal{F}_\xi$ is a finite $\mathcal{O}_{X, \xi}$-module
whose support is $\{\mathfrak m_\xi\}$. Hence it has finite length
by Algebra, Lemma \ref{algebra-lemma-support-point}.
\end{proof}

\begin{definition}
\label{definition-cycle-associated-to-coherent-sheaf}
Let $(S, \delta)$ be as in Situation \ref{situation-setup}.
Let $X$ be locally of finite type over $S$.
Let $\mathcal{F}$ be a coherent $\mathcal{O}_X$-module.
\begin{enumerate}
\item For any irreducible component $Z' \subset \text{Supp}(\mathcal{F})$
with generic point $\xi$ the integer
$m_{Z', \mathcal{F}} = \text{length}_{\mathcal{O}_{X, \xi}} \mathcal{F}_\xi$
(Lemma \ref{lemma-length-finite})
is called the {\it multiplicity of $Z'$ in $\mathcal{F}$}.
\item Assume $\dim_\delta(\text{Supp}(\mathcal{F})) \leq k$.
The {\it $k$-cycle associated to $\mathcal{F}$} is
$$
[\mathcal{F}]_k
=
\sum m_{Z', \mathcal{F}}[Z']
$$
where the sum is over the irreducible components of
$\text{Supp}(\mathcal{F})$ of $\delta$-dimension $k$.
(This is a $k$-cycle by Lemma \ref{lemma-length-finite}.)
\end{enumerate}
\end{definition}

\noindent
It is important to note that we only define $[\mathcal{F}]_k$
if $\mathcal{F}$ is coherent and the $\delta$-dimension
of $\text{Supp}(\mathcal{F})$ does not exceed $k$. In other words,
by convention, if we write $[\mathcal{F}]_k$ then this implies that
$\mathcal{F}$ is coherent on $X$ and
$\dim_\delta(\text{Supp}(\mathcal{F})) \leq k$.

\begin{lemma}
\label{lemma-cycle-closed-coherent}
Let $(S, \delta)$ be as in Situation \ref{situation-setup}.
Let $X$ be locally of finite type over $S$.
Let $Z \subset X$ be a closed subscheme.
If $\dim_\delta(Z) \leq k$, then $[Z]_k = [{\mathcal O}_Z]_k$.
\end{lemma}

\begin{proof}
This is because in this case the multiplicities $m_{Z', Z}$ and
$m_{Z', \mathcal{O}_Z}$ agree by definition.
\end{proof}

\begin{lemma}
\label{lemma-additivity-sheaf-cycle}
Let $(S, \delta)$ be as in Situation \ref{situation-setup}.
Let $X$ be locally of finite type over $S$.
Let $0 \to \mathcal{F} \to \mathcal{G} \to \mathcal{H} \to 0$
be a short exact sequence of coherent sheaves on $X$.
Assume that the $\delta$-dimension of the supports
of $\mathcal{F}$, $\mathcal{G}$, and $\mathcal{H}$ is $\leq k$.
Then $[\mathcal{G}]_k = [\mathcal{F}]_k + [\mathcal{H}]_k$.
\end{lemma}

\begin{proof}
Follows immediately from additivity of lengths, see
Algebra, Lemma \ref{algebra-lemma-length-additive}.
\end{proof}









\section{Preparation for proper pushforward}
\label{section-preparation-pushforward}

\begin{lemma}
\label{lemma-equal-dimension}
Let $(S, \delta)$ be as in Situation \ref{situation-setup}.
Let $X$, $Y$ be locally of finite type over $S$.
Let $f : X \to Y$ be a morphism.
Assume $X$, $Y$ integral and $\dim_\delta(X) = \dim_\delta(Y)$.
Then either $f(X)$ is contained in a proper closed subscheme
of $Y$, or $f$ is dominant and the extension of function fields
$R(X)/R(Y)$ is finite.
\end{lemma}

\begin{proof}
The closure $\overline{f(X)} \subset Y$ is irreducible as $X$
is irreducible (Topology, Lemmas
\ref{topology-lemma-image-irreducible-space} and
\ref{topology-lemma-irreducible}).
If $\overline{f(X)} \not = Y$, then we are done.
If $\overline{f(X)} = Y$, then $f$ is dominant and by
Morphisms,
Lemma \ref{morphisms-lemma-dominant-finite-number-irreducible-components}
we see that the generic point $\eta_Y$ of $Y$ is in the image of $f$.
Of course this implies that $f(\eta_X) = \eta_Y$, where $\eta_X \in X$
is the generic point of $X$. Since $\delta(\eta_X) = \delta(\eta_Y)$
we see that $R(Y) = \kappa(\eta_Y) \subset \kappa(\eta_X) = R(X)$
is an extension of transcendence degree $0$.
Hence $R(Y) \subset R(X)$ is a finite extension by
Morphisms, Lemma \ref{morphisms-lemma-finite-degree}
(which applies by
Morphisms, Lemma \ref{morphisms-lemma-permanence-finite-type}).
\end{proof}

\begin{lemma}
\label{lemma-quasi-compact-locally-finite}
Let $(S, \delta)$ be as in Situation \ref{situation-setup}.
Let $X$, $Y$ be locally of finite type over $S$.
Let $f : X \to Y$ be a morphism.
Assume $f$ is quasi-compact, and $\{Z_i\}_{i \in I}$ is a locally
finite collection of closed subsets of $X$.
Then $\{\overline{f(Z_i)}\}_{i \in I}$ is a locally finite
collection of closed subsets of $Y$.
\end{lemma}

\begin{proof}
Let $V \subset Y$ be a quasi-compact open subset.
Since $f$ is quasi-compact the open $f^{-1}(V)$ is
quasi-compact. Hence the set
$\{i \in I \mid Z_i \cap f^{-1}(V) \not = \emptyset \}$
is finite by a simple topological argument which we omit.
Since this is the same as the set
$$
\{i \in I \mid f(Z_i) \cap V \not = \emptyset \} =
\{i \in I \mid \overline{f(Z_i)} \cap V \not = \emptyset \}
$$
the lemma is proved.
\end{proof}









\section{Proper pushforward}
\label{section-proper-pushforward}

\begin{definition}
\label{definition-proper-pushforward}
Let $(S, \delta)$ be as in Situation \ref{situation-setup}.
Let $X$, $Y$ be locally of finite type over $S$.
Let $f : X \to Y$ be a morphism.
Assume $f$ is proper.
\begin{enumerate}
\item Let $Z \subset X$ be an integral closed subscheme
with $\dim_\delta(Z) = k$. We define
$$
f_*[Z] =
\left\{
\begin{matrix}
0 & \text{if} & \dim_\delta(f(Z))< k, \\
\deg(Z/f(Z)) [f(Z)] & \text{if} & \dim_\delta(f(Z)) = k.
\end{matrix}
\right.
$$
Here we think of $f(Z) \subset Y$ as an integral closed subscheme.
The degree of $Z$ over $f(Z)$ is finite if
$\dim_\delta(f(Z)) = \dim_\delta(Z)$
by Lemma \ref{lemma-equal-dimension}.
\item Let $\alpha = \sum n_Z [Z]$ be a $k$-cycle on $X$. The
{\it pushforward} of $\alpha$ as the sum
$$
f_* \alpha = \sum n_Z f_*[Z]
$$
where each $f_*[Z]$ is defined as above. The sum is locally finite
by Lemma \ref{lemma-quasi-compact-locally-finite} above.
\end{enumerate}
\end{definition}

\noindent
By definition the proper pushforward of cycles
$$
f_* : Z_k(X) \longrightarrow Z_k(Y)
$$
is a homomorphism of abelian groups. It turns $X \mapsto Z_k(X)$
into a covariant functor on the category of schemes locally of
finite type over $S$ with morphisms equal to proper morphisms.

\begin{lemma}
\label{lemma-compose-pushforward}
Let $(S, \delta)$ be as in Situation \ref{situation-setup}.
Let $X$, $Y$, and $Z$ be locally of finite type over $S$.
Let $f : X \to Y$ and $g : Y \to Z$ be proper morphisms.
Then $g_* \circ f_* = (g \circ f)_*$ as maps $Z_k(X) \to Z_k(Z)$.
\end{lemma}

\begin{proof}
Let $W \subset X$ be an integral closed subscheme of dimension $k$.
Consider $W' = f(W) \subset Y$ and $W'' = g(f(W)) \subset Z$.
Since $f$, $g$ are proper we see that $W'$ (resp.\ $W''$) is
an integral closed subscheme of $Y$ (resp.\ $Z$).
We have to show that $g_*(f_*[W]) = (g \circ f)_*[W]$.
If $\dim_\delta(W'') < k$, then both sides are zero.
If $\dim_\delta(W'') = k$, then we see the induced morphisms
$$
W \longrightarrow
W' \longrightarrow
W''
$$
both satisfy the hypotheses of Lemma \ref{lemma-equal-dimension}. Hence
$$
g_*(f_*[W]) = \deg(W/W')\deg(W'/W'')[W''],
\quad
(g \circ f)_*[W] = \deg(W/W'')[W''].
$$
Then we can apply
Morphisms, Lemma \ref{morphisms-lemma-degree-composition}
to conclude.
\end{proof}

\noindent
A closed immersion is proper. If $i : Z \to X$ is a closed immersion
then the maps
$$
i_* : Z_k(Z) \longrightarrow Z_k(X)
$$
are all {\it injective}.

\begin{lemma}
\label{lemma-exact-sequence-closed}
Let $(S, \delta)$ be as in Situation \ref{situation-setup}.
Let $X$ be locally of finite type over $S$. Let $X_1, X_2 \subset X$
be closed subschemes such that $X = X_1 \cup X_2$ set theoretically.
For every $k \in \mathbf{Z}$ the sequence of abelian groups
$$
\xymatrix{
Z_k(X_1 \cap X_2) \ar[r] &
Z_k(X_1) \oplus Z_k(X_2) \ar[r] &
Z_k(X) \ar[r] &
0
}
$$
is exact. Here $X_1 \cap X_2$ is the scheme theoretic intersection and
the maps are the pushforward maps with one multiplied by $-1$.
\end{lemma}

\begin{proof}
First assume $X$ is quasi-compact. Then $Z_k(X)$ is a free $\mathbf{Z}$-module
with basis given by the elements $[Z]$ where $Z \subset X$ is integral
closed of $\delta$-dimension $k$. The groups
$Z_k(X_1)$, $Z_k(X_2)$, $Z_k(X_1 \cap X_2)$ are free on the subset of these
$Z$ such that $Z \subset X_1$, $Z \subset X_2$, $Z \subset X_1 \cap X_2$.
This immediately proves the lemma in this case. The general case is similar
and the proof is omitted.
\end{proof}

\begin{lemma}
\label{lemma-cycle-push-sheaf}
Let $(S, \delta)$ be as in Situation \ref{situation-setup}.
Let $f : X \to Y$ be a proper morphism of schemes which are
locally of finite type over $S$.
\begin{enumerate}
\item Let $Z \subset X$ be a closed subscheme with $\dim_\delta(Z) \leq k$.
Then
$$
f_*[Z]_k = [f_*{\mathcal O}_Z]_k.
$$
\item Let $\mathcal{F}$ be a coherent sheaf on $X$ such that
$\dim_\delta(\text{Supp}(\mathcal{F})) \leq k$. Then
$$
f_*[\mathcal{F}]_k = [f_*{\mathcal F}]_k.
$$
\end{enumerate}
Note that the statement makes sense since $f_*\mathcal{F}$ and
$f_*\mathcal{O}_Z$ are coherent $\mathcal{O}_Y$-modules by
Cohomology of Schemes, Proposition
\ref{coherent-proposition-proper-pushforward-coherent}.
\end{lemma}

\begin{proof}
Part (1) follows from (2) and Lemma \ref{lemma-cycle-closed-coherent}.
Let $\mathcal{F}$ be a coherent sheaf on $X$.
Assume that $\dim_\delta(\text{Supp}(\mathcal{F})) \leq k$.
By Cohomology of Schemes, Lemma \ref{coherent-lemma-coherent-support-closed}
there exists a closed subscheme $i : Z \to X$ and a coherent
$\mathcal{O}_Z$-module $\mathcal{G}$ such that
$i_*\mathcal{G} \cong \mathcal{F}$ and such that the support
of $\mathcal{F}$ is $Z$. Let $Z' \subset Y$ be the scheme theoretic image
of $f|_Z : Z \to Y$. Consider the commutative diagram of schemes
$$
\xymatrix{
Z \ar[r]_i \ar[d]_{f|_Z} &
X \ar[d]^f \\
Z' \ar[r]^{i'} & Y
}
$$
We have $f_*\mathcal{F} = f_*i_*\mathcal{G} = i'_*(f|_Z)_*\mathcal{G}$
by going around the diagram in two ways. Suppose we know the result holds
for closed immersions and for $f|_Z$. Then we see that
$$
f_*[\mathcal{F}]_k = f_*i_*[\mathcal{G}]_k
= (i')_*(f|_Z)_*[\mathcal{G}]_k =
(i')_*[(f|_Z)_*\mathcal{G}]_k =
[(i')_*(f|_Z)_*\mathcal{G}]_k = [f_*\mathcal{F}]_k
$$
as desired. The case of a closed immersion is straightforward (omitted).
Note that $f|_Z : Z \to Z'$ is a dominant morphism (see
Morphisms, Lemma \ref{morphisms-lemma-quasi-compact-scheme-theoretic-image}).
Thus we have reduced to the case where
$\dim_\delta(X) \leq k$ and $f : X \to Y$ is proper and dominant.

\medskip\noindent
Assume $\dim_\delta(X) \leq k$ and $f : X \to Y$ is proper and dominant.
Since $f$ is dominant, for every irreducible component $Z \subset Y$
with generic point $\eta$ there exists a point $\xi \in X$ such
that $f(\xi) = \eta$. Hence $\delta(\eta) \leq \delta(\xi) \leq k$.
Thus we see that in the expressions
$$
f_*[\mathcal{F}]_k = \sum n_Z[Z],
\quad
\text{and}
\quad
[f_*\mathcal{F}]_k = \sum m_Z[Z].
$$
whenever $n_Z \not = 0$, or $m_Z \not = 0$ the integral closed
subscheme $Z$ is actually an irreducible component of $Y$ of
$\delta$-dimension $k$. Pick such an integral closed subscheme
$Z \subset Y$ and denote $\eta$ its generic point. Note that for
any $\xi \in X$ with $f(\xi) = \eta$ we have $\delta(\xi) \geq k$
and hence $\xi$ is a generic point of an irreducible component
of $X$ of $\delta$-dimension $k$ as well
(see Lemma \ref{lemma-multiplicity-finite}). Since $f$ is quasi-compact
and $X$ is locally Noetherian, there can be only finitely many of
these and hence $f^{-1}(\{\eta\})$ is finite.
By Morphisms, Lemma \ref{morphisms-lemma-generically-finite} there exists
an open neighbourhood $\eta \in V \subset Y$ such that $f^{-1}(V) \to V$
is finite. Replacing $Y$ by $V$ and $X$ by $f^{-1}(V)$ we reduce to the
case where $Y$ is affine, and $f$ is finite.

\medskip\noindent
Write $Y = \Spec(R)$ and $X = \Spec(A)$ (possible as
a finite morphism is affine).
Then $R$ and $A$ are Noetherian rings and $A$ is finite over $R$.
Moreover $\mathcal{F} = \widetilde{M}$ for some finite $A$-module
$M$. Note that $f_*\mathcal{F}$ corresponds to $M$ viewed as an $R$-module.
Let $\mathfrak p \subset R$ be the minimal prime corresponding
to $\eta \in Y$. The coefficient of $Z$ in $[f_*\mathcal{F}]_k$
is clearly $\text{length}_{R_{\mathfrak p}}(M_{\mathfrak p})$.
Let $\mathfrak q_i$, $i = 1, \ldots, t$ be the primes of $A$
lying over $\mathfrak p$. Then $A_{\mathfrak p} = \prod A_{\mathfrak q_i}$
since $A_{\mathfrak p}$ is an Artinian ring being finite over the
dimension zero local Noetherian ring $R_{\mathfrak p}$.
Clearly the coefficient of $Z$ in $f_*[\mathcal{F}]_k$ is
$$
\sum\nolimits_{i = 1, \ldots, t}
[\kappa(\mathfrak q_i) : \kappa(\mathfrak p)]
\text{length}_{A_{\mathfrak q_i}}(M_{\mathfrak q_i})
$$
Hence the desired equality follows from
Algebra, Lemma \ref{algebra-lemma-pushdown-module}.
\end{proof}




















\section{Preparation for flat pullback}
\label{section-preparation-flat-pullback}


\noindent
Recall that a morphism $f : X \to Y$ which is locally of finite type
is said to have relative dimension $r$ if every nonempty fibre
is equidimensional of dimension $r$. See
Morphisms, Definition \ref{morphisms-definition-relative-dimension-d}.

\begin{lemma}
\label{lemma-flat-inverse-image-dimension}
Let $(S, \delta)$ be as in Situation \ref{situation-setup}.
Let $X$, $Y$ be locally of finite type over $S$.
Let $f : X \to Y$ be a morphism.
Assume $f$ is flat of relative dimension $r$.
For any closed subset $Z \subset Y$ we have
$$
\dim_\delta(f^{-1}(Z)) = \dim_\delta(Z) + r.
$$
provided $f^{-1}(Z)$ is nonempty.
If $Z$ is irreducible and $Z' \subset f^{-1}(Z)$ is an irreducible
component, then $Z'$ dominates $Z$ and
$\dim_\delta(Z') = \dim_\delta(Z) + r$.
\end{lemma}

\begin{proof}
It suffices to prove the final statement.
We may replace $Y$ by the integral closed subscheme $Z$ and
$X$ by the scheme theoretic inverse image $f^{-1}(Z) = Z \times_Y X$.
Hence we may assume $Z = Y$ is integral and $f$ is a flat morphism
of relative dimension $r$. Since $Y$ is locally Noetherian the
morphism $f$ which is locally of finite type,
is actually locally of finite presentation. Hence
Morphisms, Lemma \ref{morphisms-lemma-fppf-open}
applies and we see that $f$ is open.
Let $\xi \in X$ be a generic point of an irreducible component
of $X$. By the openness of $f$ we see that $f(\xi)$ is the
generic point $\eta$ of $Z = Y$. Note that $\dim_\xi(X_\eta) = r$
by assumption that $f$ has relative dimension $r$. On the other
hand, since $\xi$ is a generic point of $X$ we see that
$\mathcal{O}_{X, \xi} = \mathcal{O}_{X_\eta, \xi}$ has only one
prime ideal and hence has dimension $0$. Thus by
Morphisms, Lemma \ref{morphisms-lemma-dimension-fibre-at-a-point}
we conclude that the transcendence
degree of $\kappa(\xi)$ over $\kappa(\eta)$ is $r$.
In other words, $\delta(\xi) = \delta(\eta) + r$ as desired.
\end{proof}

\noindent
Here is the lemma that we will use to prove that the flat pullback
of a locally finite collection of closed subschemes is locally finite.

\begin{lemma}
\label{lemma-inverse-image-locally-finite}
Let $(S, \delta)$ be as in Situation \ref{situation-setup}.
Let $X$, $Y$ be locally of finite type over $S$.
Let $f : X \to Y$ be a morphism.
Assume $\{Z_i\}_{i \in I}$ is a locally
finite collection of closed subsets of $Y$.
Then $\{f^{-1}(Z_i)\}_{i \in I}$ is a locally finite
collection of closed subsets of $X$.
\end{lemma}

\begin{proof}
Let $U \subset X$ be a quasi-compact open subset.
Since the image $f(U) \subset Y$ is a quasi-compact subset
there exists a quasi-compact open $V \subset Y$ such that
$f(U) \subset V$. Note that
$$
\{i \in I \mid f^{-1}(Z_i) \cap U \not = \emptyset \}
\subset
\{i \in I \mid Z_i \cap V \not = \emptyset \}.
$$
Since the right hand side is finite by assumption we win.
\end{proof}



\section{Flat pullback}
\label{section-flat-pullback}

\noindent
In the following we use $f^{-1}(Z)$ to denote the
{\it scheme theoretic inverse image} of a closed subscheme
$Z \subset Y$ for a morphism of schemes $f : X \to Y$.
We recall that the scheme theoretic inverse image is the fibre product
$$
\xymatrix{
f^{-1}(Z) \ar[r] \ar[d] & X \ar[d] \\
Z \ar[r] & Y
}
$$
and it is also the closed subscheme of $X$ cut out by the
quasi-coherent sheaf of ideals $f^{-1}(\mathcal{I})\mathcal{O}_X$, if
$\mathcal{I} \subset \mathcal{O}_Y$ is the quasi-coherent sheaf of ideals
corresponding to $Z$ in $Y$.
(This is discussed in
Schemes, Section \ref{schemes-section-closed-immersion} and
Lemma \ref{schemes-lemma-fibre-product-immersion}
and Definition \ref{schemes-definition-inverse-image-closed-subscheme}.)

\begin{definition}
\label{definition-flat-pullback}
Let $(S, \delta)$ be as in Situation \ref{situation-setup}.
Let $X$, $Y$ be locally of finite type over $S$.
Let $f : X \to Y$ be a morphism.
Assume $f$ is flat of relative dimension $r$.
\begin{enumerate}
\item Let $Z \subset Y$ be an integral closed subscheme of
$\delta$-dimension $k$. We define $f^*[Z]$ to be the
$(k+r)$-cycle on $X$ to the scheme theoretic inverse image
$$
f^*[Z] = [f^{-1}(Z)]_{k+r}.
$$
This makes sense since $\dim_\delta(f^{-1}(Z)) = k + r$
by Lemma \ref{lemma-flat-inverse-image-dimension}.
\item Let $\alpha = \sum n_i [Z_i]$ be
a $k$-cycle on $Y$. The {\it flat pullback of $\alpha$ by $f$}
is the sum
$$
f^* \alpha = \sum n_i f^*[Z_i]
$$
where each $f^*[Z_i]$ is defined as above.
The sum is locally finite by Lemma \ref{lemma-inverse-image-locally-finite}.
\item We denote $f^* : Z_k(Y) \to Z_{k + r}(X)$ the map of abelian
groups so obtained.
\end{enumerate}
\end{definition}

\noindent
An open immersion is flat. This is an important though trivial special
case of a flat morphism. If $U \subset X$ is open then sometimes the
pullback by $j : U \to X$ of a cycle is called the {\it restriction} of the
cycle to $U$. Note that in this case the maps
$$
j^* : Z_k(X) \longrightarrow Z_k(U)
$$
are all {\it surjective}. The reason is that given any integral closed
subscheme $Z' \subset U$, we can take the closure of $Z$ of $Z'$ in $X$
and think of it as a reduced closed subscheme of $X$ (see
Schemes, Lemma \ref{schemes-lemma-reduced-closed-subscheme}).
And clearly $Z \cap U = Z'$, in other words
$j^*[Z] = [Z']$ whence the surjectivity. In fact a little bit more
is true.

\begin{lemma}
\label{lemma-exact-sequence-open}
Let $(S, \delta)$ be as in Situation \ref{situation-setup}.
Let $X$ be locally of finite type over $S$.
Let $U \subset X$ be an open subscheme, and denote
$i : Y = X \setminus U \to X$ as a reduced closed subscheme of $X$.
For every $k \in \mathbf{Z}$ the sequence
$$
\xymatrix{
Z_k(Y) \ar[r]^{i_*} & Z_k(X) \ar[r]^{j^*} & Z_k(U) \ar[r] & 0
}
$$
is an exact complex of abelian groups.
\end{lemma}

\begin{proof}
First assume $X$ is quasi-compact. Then $Z_k(X)$ is a free $\mathbf{Z}$-module
with basis given by the elements $[Z]$ where $Z \subset X$ is integral
closed of $\delta$-dimension $k$. Such a basis element maps
either to the basis element $[Z \cap U]$ or to zero if $Z \subset Y$.
Hence the lemma is clear in this case. The general case is similar
and the proof is omitted.
\end{proof}

\begin{lemma}
\label{lemma-compose-flat-pullback}
Let $(S, \delta)$ be as in Situation \ref{situation-setup}.
Let $X, Y, Z$ be locally of finite type over $S$.
Let $f : X \to Y$ and $g : Y \to Z$ be flat morphisms of relative dimensions
$r$ and $s$. Then $g \circ f$ is flat of relative dimension
$r + s$ and
$$
f^* \circ g^* = (g \circ f)^*
$$
as maps $Z_k(Z) \to Z_{k + r + s}(X)$.
\end{lemma}

\begin{proof}
The composition is flat of relative dimension $r + s$ by
Morphisms, Lemma \ref{morphisms-lemma-composition-relative-dimension-d}.
Suppose that
\begin{enumerate}
\item $W \subset Z$ is a closed integral subscheme of $\delta$-dimension $k$,
\item $W' \subset Y$ is a closed integral subscheme of $\delta$-dimension
$k + s$ with $W' \subset g^{-1}(W)$, and
\item $W'' \subset Y$ is a closed integral subscheme of $\delta$-dimension
$k + s + r$ with $W'' \subset f^{-1}(W')$.
\end{enumerate}
We have to show that the coefficient $n$ of $[W'']$ in
$(g \circ f)^*[W]$ agrees with the coefficient $m$ of
$[W'']$ in $f^*(g^*[W])$. That it suffices to check the lemma in these
cases follows from Lemma \ref{lemma-flat-inverse-image-dimension}.
Let $\xi'' \in W''$, $\xi' \in W'$
and $\xi \in W$ be the generic points. Consider the local rings
$A = \mathcal{O}_{Z, \xi}$, $B = \mathcal{O}_{Y, \xi'}$
and $C = \mathcal{O}_{X, \xi''}$. Then we have local flat ring maps
$A \to B$, $B \to C$ and moreover
$$
n = \text{length}_C(C/\mathfrak m_AC),
\quad
\text{and}
\quad
m = \text{length}_C(C/\mathfrak m_BC) \text{length}_B(B/\mathfrak m_AB)
$$
Hence the equality follows from
Algebra, Lemma \ref{algebra-lemma-pullback-transitive}.
\end{proof}

\begin{lemma}
\label{lemma-pullback-coherent}
Let $(S, \delta)$ be as in Situation \ref{situation-setup}.
Let $X, Y$ be locally of finite type over $S$.
Let $f : X \to Y$ be a flat morphism of relative dimension $r$.
\begin{enumerate}
\item Let $Z \subset Y$ be a closed subscheme with
$\dim_\delta(Z) \leq k$. Then we have
$\dim_\delta(f^{-1}(Z)) \leq k + r$
and $[f^{-1}(Z)]_{k + r} = f^*[Z]_k$ in $Z_{k + r}(X)$.
\item Let $\mathcal{F}$ be a coherent sheaf on $Y$ with
$\dim_\delta(\text{Supp}(\mathcal{F})) \leq k$.
Then we have $\dim_\delta(\text{Supp}(f^*\mathcal{F})) \leq k + r$
and
$$
f^*[{\mathcal F}]_k = [f^*{\mathcal F}]_{k+r}
$$
in $Z_{k + r}(X)$.
\end{enumerate}
\end{lemma}

\begin{proof}
The statements on dimensions follow immediately from
Lemma \ref{lemma-flat-inverse-image-dimension}.
Part (1) follows from part (2) by Lemma \ref{lemma-cycle-closed-coherent}
and the fact that $f^*\mathcal{O}_Z = \mathcal{O}_{f^{-1}(Z)}$.

\medskip\noindent
Proof of (2).
As $X$, $Y$ are locally Noetherian we may apply
Cohomology of Schemes, Lemma \ref{coherent-lemma-coherent-Noetherian} to see
that $\mathcal{F}$ is of finite type, hence $f^*\mathcal{F}$ is
of finite type (Modules, Lemma \ref{modules-lemma-pullback-finite-type}),
hence $f^*\mathcal{F}$ is coherent
(Cohomology of Schemes, Lemma \ref{coherent-lemma-coherent-Noetherian} again).
Thus the lemma makes sense. Let $W \subset Y$ be an integral closed
subscheme of $\delta$-dimension $k$, and let $W' \subset X$ be
an integral closed subscheme of dimension $k + r$ mapping into $W$
under $f$. We have to show that the coefficient $n$ of
$[W']$ in $f^*[{\mathcal F}]_k$ agrees with the coefficient
$m$ of $[W']$ in $[f^*{\mathcal F}]_{k+r}$. Let $\xi \in W$ and
$\xi' \in W'$ be the generic points. Let
$A = \mathcal{O}_{Y, \xi}$, $B = \mathcal{O}_{X, \xi'}$
and set $M = \mathcal{F}_\xi$ as an $A$-module. (Note that
$M$ has finite length by our dimension assumptions, but we
actually do not need to verify this. See
Lemma \ref{lemma-length-finite}.)
We have $f^*\mathcal{F}_{\xi'} = B \otimes_A M$.
Thus we see that
$$
n = \text{length}_B(B \otimes_A M)
\quad
\text{and}
\quad
m = \text{length}_A(M) \text{length}_B(B/\mathfrak m_AB)
$$
Thus the equality follows from
Algebra, Lemma \ref{algebra-lemma-pullback-module}.
\end{proof}



\section{Push and pull}
\label{section-push-pull}

\noindent
In this section we verify that proper pushforward and flat pullback
are compatible when this makes sense. By the work we did above this
is a consequence of cohomology and base change.

\begin{lemma}
\label{lemma-flat-pullback-proper-pushforward}
Let $(S, \delta)$ be as in Situation \ref{situation-setup}.
Let
$$
\xymatrix{
X' \ar[r]_{g'} \ar[d]_{f'} & X \ar[d]^f \\
Y' \ar[r]^g & Y
}
$$
be a fibre product diagram of schemes locally of finite type over $S$.
Assume $f : X \to Y$ proper and $g : Y' \to Y$ flat of relative dimension $r$.
Then also $f'$ is proper and $g'$ is flat of relative dimension $r$.
For any $k$-cycle $\alpha$ on $X$ we have
$$
g^*f_*\alpha = f'_*(g')^*\alpha
$$
in $Z_{k + r}(Y')$.
\end{lemma}

\begin{proof}
The assertion that $f'$ is proper follows from
Morphisms, Lemma \ref{morphisms-lemma-base-change-proper}.
The assertion that $g'$ is flat of relative dimension $r$ follows from
Morphisms, Lemmas \ref{morphisms-lemma-base-change-relative-dimension-d}
and \ref{morphisms-lemma-base-change-flat}.
It suffices to prove the equality of cycles when $\alpha = [W]$
for some integral closed subscheme $W \subset X$ of $\delta$-dimension $k$.
Note that in this case we have $\alpha = [\mathcal{O}_W]_k$, see
Lemma \ref{lemma-cycle-closed-coherent}.
By Lemmas \ref{lemma-cycle-push-sheaf} and
\ref{lemma-pullback-coherent} it therefore suffices
to show that $f'_*(g')^*\mathcal{O}_W$ is isomorphic to
$g^*f_*\mathcal{O}_W$. This follows from cohomology and
base change, see
Cohomology of Schemes, Lemma \ref{coherent-lemma-flat-base-change-cohomology}.
\end{proof}

\begin{lemma}
\label{lemma-finite-flat}
Let $(S, \delta)$ be as in Situation \ref{situation-setup}.
Let $X$, $Y$ be locally of finite type over $S$.
Let $f : X \to Y$ be a finite locally free morphism
of degree $d$ (see
Morphisms, Definition \ref{morphisms-definition-finite-locally-free}).
Then $f$ is both proper and flat of relative dimension $0$, and
$$
f_*f^*\alpha = d\alpha
$$
for every $\alpha \in Z_k(Y)$.
\end{lemma}

\begin{proof}
A finite locally free morphism is flat and finite by
Morphisms, Lemma \ref{morphisms-lemma-finite-flat},
and a finite morphism is proper
by Morphisms, Lemma \ref{morphisms-lemma-finite-proper}.
We omit showing that a finite
morphism has relative dimension $0$. Thus the formula makes sense.
To prove it, let $Z \subset Y$ be an integral closed subscheme
of $\delta$-dimension $k$. It suffices to prove the formula
for $\alpha = [Z]$. Since the base change of a finite locally free
morphism is finite locally free
(Morphisms, Lemma \ref{morphisms-lemma-base-change-finite-locally-free})
we see that $f_*f^*\mathcal{O}_Z$ is a finite locally free sheaf of
rank $d$ on $Z$. Hence
$$
f_*f^*[Z] = f_*f^*[\mathcal{O}_Z]_k =
[f_*f^*\mathcal{O}_Z]_k = d[Z]
$$
where we have used Lemmas \ref{lemma-pullback-coherent} and
\ref{lemma-cycle-push-sheaf}.
\end{proof}








\section{Preparation for principal divisors}
\label{section-preparation-principal-divisors}

\noindent
Some of the material in this section partially overlaps with the
discussion in Divisors, Section \ref{divisors-section-Weil-divisors}.

\begin{lemma}
\label{lemma-divisor-delta-dimension}
Let $(S, \delta)$ be as in Situation \ref{situation-setup}.
Let $X$ be locally of finite type over $S$. Assume $X$ is
integral.
\begin{enumerate}
\item If $Z \subset X$ is an integral closed subscheme, then
the following are equivalent:
\begin{enumerate}
\item $Z$ is a prime divisor,
\item $Z$ has codimension $1$ in $X$, and
\item $\dim_\delta(Z) = \dim_\delta(X) - 1$.
\end{enumerate}
\item If $Z$ is an irreducible component of an effective Cartier
divisor on $X$, then $\dim_\delta(Z) = \dim_\delta(X) - 1$.
\end{enumerate}
\end{lemma}

\begin{proof}
Part (1) follows from the definition of a prime divisor
(Divisors, Definition \ref{divisors-definition-Weil-divisor})
and the definition of a dimension function
(Topology, Definition \ref{topology-definition-dimension-function}).
Let $\xi \in Z$ be the generic point of an irreducible component $Z$ of
an effective Cartier divisor $D \subset X$.
Then $\dim(\mathcal{O}_{D, \xi}) = 0$ and
$\mathcal{O}_{D, \xi} = \mathcal{O}_{X, \xi}/(f)$ for some
nonzerodivisor $f \in \mathcal{O}_{X, \xi}$ (Divisors,
Lemma \ref{divisors-lemma-effective-Cartier-in-points}).
Then $\dim(\mathcal{O}_{X, \xi}) = 1$ by
Algebra, Lemma \ref{algebra-lemma-one-equation}. Hence $Z$ is as in (1) by
Properties, Lemma \ref{properties-lemma-codimension-local-ring}
and the proof is complete.
\end{proof}

\begin{lemma}
\label{lemma-finite-in-codimension-one}
Let $f : X \to Y$ be a morphism of schemes.
Let $\xi \in Y$ be a point.
Assume that
\begin{enumerate}
\item $X$, $Y$ are integral,
\item $Y$ is locally Noetherian
\item $f$ is proper, dominant and $R(Y) \subset R(X)$ is finite, and
\item $\dim(\mathcal{O}_{Y, \xi}) = 1$.
\end{enumerate}
Then there exists an open neighbourhood $V \subset Y$ of $\xi$
such that $f|_{f^{-1}(V)} : f^{-1}(V) \to V$ is finite.
\end{lemma}

\begin{proof}
This lemma is a special case of
Varieties, Lemma \ref{varieties-lemma-finite-in-codim-1}.
Here is a direct argument in this case.
By Cohomology of Schemes,
Lemma \ref{coherent-lemma-proper-finite-fibre-finite-in-neighbourhood}
it suffices to prove that $f^{-1}(\{\xi\})$ is finite.
We replace $Y$ by an affine open, say $Y = \Spec(R)$.
Note that $R$ is Noetherian, as $Y$ is assumed locally Noetherian.
Since $f$ is proper it is quasi-compact. Hence we can find a finite
affine open covering $X = U_1 \cup \ldots \cup U_n$ with
each $U_i = \Spec(A_i)$. Note that $R \to A_i$ is a
finite type injective homomorphism of domains such that
the induced extension of fraction fields is finite.
Thus the lemma follows
from Algebra, Lemma \ref{algebra-lemma-finite-in-codim-1}.
\end{proof}


\section{Principal divisors}
\label{section-principal-divisors}

\noindent
The following definition is the analogue of
Divisors, Definition \ref{divisors-definition-principal-divisor}
in our current setup.

\begin{definition}
\label{definition-principal-divisor}
Let $(S, \delta)$ be as in Situation \ref{situation-setup}.
Let $X$ be locally of finite type over $S$. Assume $X$ is
integral with $\dim_\delta(X) = n$.
Let $f \in R(X)^*$. The {\it principal divisor
associated to $f$} is the $(n - 1)$-cycle
$$
\text{div}(f) = \text{div}_X(f) = \sum \text{ord}_Z(f) [Z]
$$
defined in Divisors, Definition \ref{divisors-definition-principal-divisor}.
This makes sense because prime divisors have $\delta$-dimension $n - 1$ by
Lemma \ref{lemma-divisor-delta-dimension}.
\end{definition}

\noindent
In the situation of the definition for $f, g \in R(X)^*$ we have
$$
\text{div}_X(fg) = \text{div}_X(f) + \text{div}_X(g)
$$
in $Z_{n - 1}(X)$. See Divisors, Lemma \ref{divisors-lemma-div-additive}.
The following lemma will be superseded by the more general
Lemma \ref{lemma-flat-pullback-rational-equivalence}.

\begin{lemma}
\label{lemma-flat-pullback-principal-divisor}
Let $(S, \delta)$ be as in Situation \ref{situation-setup}.
Let $X$, $Y$ be locally of finite type over $S$. Assume $X$, $Y$
are integral and $n = \dim_\delta(Y)$.
Let $f : X \to Y$ be a flat morphism of relative dimension $r$.
Let $g \in R(Y)^*$. Then
$$
f^*(\text{div}_Y(g)) = \text{div}_X(g)
$$
in $Z_{n + r - 1}(X)$.
\end{lemma}

\begin{proof}
Note that since $f$ is flat it is dominant so that
$f$ induces an embedding $R(Y) \subset R(X)$, and hence
we may think of $g$ as an element of $R(X)^*$.
Let $Z \subset X$ be an integral closed subscheme of
$\delta$-dimension $n + r - 1$. Let $\xi \in Z$
be its generic point. If $\dim_\delta(f(Z)) > n - 1$,
then we see that the coefficient of $[Z]$ in the left and
right hand side of the equation is zero.
Hence we may assume that $Z' = \overline{f(Z)}$ is an
integral closed subscheme of $Y$ of $\delta$-dimension $n - 1$.
Let $\xi' = f(\xi)$. It is the generic point of $Z'$.
Set $A = \mathcal{O}_{Y, \xi'}$, $B = \mathcal{O}_{X, \xi}$.
The ring map $A \to B$ is a flat local homomorphism of
Noetherian local domains of dimension $1$.
We have $g$ in the fraction field of $A$. What we have to show is that
$$
\text{ord}_A(g) \text{length}_B(B/\mathfrak m_AB)
=
\text{ord}_B(g).
$$
This follows from Algebra, Lemma \ref{algebra-lemma-pullback-module}
(details omitted).
\end{proof}











\section{Principal divisors and pushforward}
\label{section-two-fun}

\noindent
The first lemma implies that the pushforward of a principal
divisor along a generically finite morphism is a principal divisor.

\begin{lemma}
\label{lemma-proper-pushforward-alteration}
Let $(S, \delta)$ be as in Situation \ref{situation-setup}.
Let $X$, $Y$ be locally of finite type over $S$. Assume $X$, $Y$
are integral and $n = \dim_\delta(X) = \dim_\delta(Y)$.
Let $p : X \to Y$ be a dominant proper morphism.
Let $f \in R(X)^*$. Set
$$
g = \text{Nm}_{R(X)/R(Y)}(f).
$$
Then we have
$p_*\text{div}(f) = \text{div}(g)$.
\end{lemma}

\begin{proof}
Let $Z \subset Y$ be an integral closed subscheme of $\delta$-dimension
$n - 1$. We want to show that the coefficient of $[Z]$ in
$p_*\text{div}(f)$ and $\text{div}(g)$ are equal. We may apply
Lemma \ref{lemma-finite-in-codimension-one}
to the morphism $p : X \to Y$ and the generic point $\xi \in Z$.
Hence we may replace $Y$ by an
affine open neighbourhood of $\xi$ and assume that $p : X \to Y$ is finite.
Write $Y = \Spec(R)$ and $X = \Spec(A)$ with $p$ induced
by a finite homomorphism $R \to A$ of Noetherian domains which induces
an finite field extension $L/K$ of fraction fields.
Now we have $f \in L$, $g = \text{Nm}(f) \in K$,
and a prime $\mathfrak p \subset R$ with $\dim(R_{\mathfrak p}) = 1$.
The coefficient of $[Z]$ in $\text{div}_Y(g)$ is
$\text{ord}_{R_\mathfrak p}(g)$.
The coefficient of $[Z]$ in $p_*\text{div}_X(f)$ is
$$
\sum\nolimits_{\mathfrak q\text{ lying over }\mathfrak p}
[\kappa(\mathfrak q) : \kappa(\mathfrak p)]
\text{ord}_{A_{\mathfrak q}}(f)
$$
The desired equality therefore follows from
Algebra, Lemma \ref{algebra-lemma-finite-extension-dim-1}.
\end{proof}

\noindent
An important role in the discussion of principal divisors
is played by the ``universal'' principal divisor $[0] - [\infty]$
on $\mathbf{P}^1_S$. To make this more precise, let us denote
\begin{equation}
\label{equation-zero-infty}
D_0, D_\infty \subset
\mathbf{P}^1_S = \underline{\text{Proj}}_S(\mathcal{O}_S[T_0, T_1])
\end{equation}
the closed subscheme cut out by the section $T_1$, resp.\ $T_0$
of $\mathcal{O}(1)$. These are effective Cartier divisors, see
Divisors, Definition \ref{divisors-definition-effective-Cartier-divisor}
and Lemma \ref{divisors-lemma-characterize-OD}.
The following lemma says that loosely speaking we have
``$\text{div}(T_1/T_0) = [D_0] - [D_1]$'' and that this is the
universal principal divisor.

\begin{lemma}
\label{lemma-rational-function}
Let $(S, \delta)$ be as in Situation \ref{situation-setup}.
Let $X$ be locally of finite type over $S$. Assume $X$ is
integral and $n = \dim_\delta(X)$. Let $f \in R(X)^*$.
Let $U \subset X$ be a nonempty open such that $f$
corresponds to a section $f \in \Gamma(U, \mathcal{O}_X^*)$.
Let $Y \subset X \times_S \mathbf{P}^1_S$ be the
closure of the graph of $f : U \to \mathbf{P}^1_S$.
Then
\begin{enumerate}
\item the projection morphism $p : Y \to X$ is proper,
\item $p|_{p^{-1}(U)} : p^{-1}(U) \to U$ is an isomorphism,
\item the pullbacks $Y_0 = q^{-1}D_0$ and $Y_\infty = q^{-1}D_\infty$
via the morphism $q : Y \to \mathbf{P}^1_S$ are defined
(Divisors, Definition
\ref{divisors-definition-pullback-effective-Cartier-divisor}),
\item we have
$$
\text{div}_Y(f) = [Y_0]_{n - 1} - [Y_\infty]_{n - 1}
$$
\item we have
$$
\text{div}_X(f) = p_*\text{div}_Y(f)
$$
\item if we view $Y_0$ and $Y_\infty$ as closed subschemes of $X$
via the morphism $p$ then we have
$$
\text{div}_X(f) = [Y_0]_{n - 1} - [Y_\infty]_{n - 1}
$$
\end{enumerate}
\end{lemma}

\begin{proof}
Since $X$ is integral, we see that $U$ is integral.
Hence $Y$ is integral, and $(1, f)(U) \subset Y$ is an open dense subscheme.
Also, note that the closed subscheme $Y \subset X \times_S \mathbf{P}^1_S$
does not depend on the choice of the open $U$, since after all it is
the closure of the one point set $\{\eta'\} = \{(1, f)(\eta)\}$
where $\eta \in X$ is the generic point. Having said this let us
prove the assertions of the lemma.

\medskip\noindent
For (1) note that $p$ is the composition of the closed immersion
$Y \to X \times_S \mathbf{P}^1_S = \mathbf{P}^1_X$ with the proper
morphism $\mathbf{P}^1_X \to X$. As a composition of proper morphisms
is proper (Morphisms, Lemma \ref{morphisms-lemma-composition-proper})
we conclude.

\medskip\noindent
It is clear that $Y \cap U \times_S \mathbf{P}^1_S = (1, f)(U)$.
Thus (2) follows. It also follows that $\dim_\delta(Y) = n$.

\medskip\noindent
Note that $q(\eta') = f(\eta)$ is not contained in $D_0$ or $D_\infty$
since $f \in R(X)^*$. Hence (3) by
Divisors, Lemma \ref{divisors-lemma-pullback-effective-Cartier-defined}.
We obtain $\dim_\delta(Y_0) = n - 1$
and $\dim_\delta(Y_\infty) = n - 1$ from
Lemma \ref{lemma-divisor-delta-dimension}.

\medskip\noindent
Consider the effective Cartier divisor $Y_0$.
At every point $\xi \in Y_0$ we have $f \in \mathcal{O}_{Y, \xi}$ and
the local equation for $Y_0$ is given by $f$.
In particular, if $\delta(\xi) = n - 1$ so $\xi$ is the generic point
of an integral closed subscheme $Z$ of $\delta$-dimension $n - 1$,
then we see that the coefficient of $[Z]$ in $\text{div}_Y(f)$ is
$$
\text{ord}_Z(f) =
\text{length}_{\mathcal{O}_{Y, \xi}}
(\mathcal{O}_{Y, \xi}/f\mathcal{O}_{Y, \xi}) =
\text{length}_{\mathcal{O}_{Y, \xi}}
(\mathcal{O}_{Y_0, \xi})
$$
which is the coefficient of $[Z]$ in $[Y_0]_{n - 1}$. A similar
argument using the rational function $1/f$ shows that
$-[Y_\infty]$ agrees with the terms with negative coefficients in
the expression for $\text{div}_Y(f)$. Hence (4) follows.

\medskip\noindent
Note that $D_0 \to S$ is an isomorphism. Hence we see that
$X \times_S D_0 \to X$ is an isomorphism as well. Clearly
we have $Y_0 = Y \cap X \times_S D_0$ (scheme theoretic intersection)
inside $X \times_S \mathbf{P}^1_S$. Hence it is really the case that
$Y_0 \to X$ is a closed immersion. It follows that
$$
p_*\mathcal{O}_{Y_0} = \mathcal{O}_{Y'_0}
$$
where $Y'_0 \subset X$ is the image of $Y_0 \to X$.
By Lemma \ref{lemma-cycle-push-sheaf} we
have $p_*[Y_0]_{n - 1} = [Y'_0]_{n - 1}$. The same
is true for $D_\infty$ and $Y_\infty$. Hence (6) is a consequence of (5).
Finally, (5) follows immediately from
Lemma \ref{lemma-proper-pushforward-alteration}.
\end{proof}

\noindent
The following lemma says that the degree of a principal divisor on
a proper curve is zero.

\begin{lemma}
\label{lemma-curve-principal-divisor}
Let $K$ be any field. Let $X$ be a $1$-dimensional integral scheme
endowed with a proper morphism $c : X \to \Spec(K)$.
Let $f \in K(X)^*$ be an invertible rational function.
Then
$$
\sum\nolimits_{x \in X \text{ closed}}
[\kappa(x) : K] \text{ord}_{\mathcal{O}_{X, x}}(f)
=
0
$$
where $\text{ord}$ is as in
Algebra, Definition \ref{algebra-definition-ord}.
In other words, $c_*\text{div}(f) = 0$.
\end{lemma}

\begin{proof}
Consider the diagram
$$
\xymatrix{
Y \ar[r]_p \ar[d]_q & X \ar[d]^c \\
\mathbf{P}^1_K \ar[r]^-{c'} & \Spec(K)
}
$$
that we constructed in Lemma \ref{lemma-rational-function}
starting with $X$ and the rational function $f$ over $S = \Spec(K)$.
We will use all the results of this lemma without further mention.
We have to show that $c_*\text{div}_X(f) = c_*p_*\text{div}_Y(f) = 0$.
This is the same as proving that $c'_*q_*\text{div}_Y(f) = 0$.
If $q(Y)$ is a closed point of $\mathbf{P}^1_K$ then we
see that $\text{div}_X(f) = 0$ and the lemma holds.
Thus we may assume that $q$ is dominant.
Suppose we can show that $q : Y \to \mathbf{P}^1_K$ is finite
locally free of degree $d$ (see
Morphisms, Definition \ref{morphisms-definition-finite-locally-free}).
Since $\text{div}_Y(f) = [q^{-1}D_0]_0 - [q^{-1}D_\infty]_0$
we see (by definition of flat pullback) that
$\text{div}_Y(f) = q^*([D_0]_0 - [D_\infty]_0)$.
Then by Lemma \ref{lemma-finite-flat} we get
$q_*\text{div}_Y(f) = d([D_0]_0 - [D_\infty]_0)$.
Since clearly $c'_*[D_0]_0 = c'_*[D_\infty]_0$ we win.

\medskip\noindent
It remains to show that $q$ is finite locally free.
(It will automatically have some given degree as $\mathbf{P}^1_K$
is connected.)
Since $\dim(\mathbf{P}^1_K) = 1$ we see that $q$ is finite for example
by Lemma \ref{lemma-finite-in-codimension-one}.
All local rings of $\mathbf{P}^1_K$ at
closed points are regular local rings of dimension $1$
(in other words discrete valuation rings), since they are
localizations of $K[T]$ (see
Algebra, Lemma \ref{algebra-lemma-dim-affine-space}).
Hence for $y\in Y$ closed the local ring $\mathcal{O}_{Y, y}$
will be flat over $\mathcal{O}_{\mathbf{P}^1_K, q(y)}$ as soon as
it is torsion free (More on Algebra, Lemma
\ref{more-algebra-lemma-dedekind-torsion-free-flat}).
This is obviously the case as
$\mathcal{O}_{Y, y}$ is a domain and $q$ is dominant.
Thus $q$ is flat. Hence $q$ is finite locally free by
Morphisms, Lemma \ref{morphisms-lemma-finite-flat}.
\end{proof}





\section{Rational equivalence}
\label{section-rational-equivalence}

\noindent
In this section we define {\it rational equivalence} on $k$-cycles.
We will allow locally finite sums of images of
principal divisors (under closed immersions). This leads to some
pretty strange phenomena, see Example \ref{example-weird}.
However, if we do not allow these then we do not know how to prove that
capping with Chern classes of line bundles factors through rational
equivalence.

\begin{definition}
\label{definition-rational-equivalence}
Let $(S, \delta)$ be as in Situation \ref{situation-setup}.
Let $X$ be a scheme locally of finite type over $S$.
Let $k \in \mathbf{Z}$.
\begin{enumerate}
\item Given any locally finite collection $\{W_j \subset X\}$
of integral closed subschemes with $\dim_\delta(W_j) = k + 1$,
and any $f_j \in R(W_j)^*$ we may consider
$$
\sum (i_j)_*\text{div}(f_j) \in Z_k(X)
$$
where $i_j : W_j \to X$ is the inclusion morphism.
This makes sense as the morphism
$\coprod i_j : \coprod W_j \to X$ is proper.
\item We say that $\alpha \in Z_k(X)$ is {\it rationally equivalent to zero}
if $\alpha$ is a cycle of the form displayed above.
\item We say $\alpha, \beta \in Z_k(X)$ are
{\it rationally equivalent} and we write $\alpha \sim_{rat} \beta$
if $\alpha - \beta$ is rationally equivalent to zero.
\item We define
$$
\CH_k(X) = Z_k(X) / \sim_{rat}
$$
to be the {\it Chow group of $k$-cycles on $X$}. This is sometimes called
the {\it Chow group of $k$-cycles modulo rational equivalence on $X$}.
\end{enumerate}
\end{definition}

\noindent
There are many other interesting (adequate) equivalence relations.
Rational equivalence is the coarsest one of them all.

\begin{remark}
\label{remark-chow-group-pointwise}
Let $(S, \delta)$ be as in Situation \ref{situation-setup}.
Let $X$ be locally of finite type over $S$. Let $k \in \mathbf{Z}$.
Let us show that we have a presentation
$$
\bigoplus\nolimits_{\delta(x) = k + 1}' K_1^M(\kappa(x))
\xrightarrow{\partial}
\bigoplus\nolimits_{\delta(x) = k}' K_0^M(\kappa(x)) \to
\CH_k(X) \to 0
$$
Here we use the notation and conventions introduced in
Remark \ref{remark-cycles-pointwise} and in addition
\begin{enumerate}
\item $K_1^M(\kappa(x)) = \kappa(x)^*$ is the degree $1$ part of
the Milnor K-theory of the residue field $\kappa(x)$ of the point
$x \in X$ (see Remark \ref{remark-gersten-complex-milnor}), and
\item the differential $\partial$ is defined as follows:
given an element $\xi = \sum_x f_x$ we denote $W_x = \overline{x}$
the integral closed subscheme of $X$ with generic point $x$ and we set
$$
\partial(\xi) = \sum (W_x \to X)_*\text{div}(f_x)
$$
in $Z_k(X)$ which makes sense as we have seen that the second
term of the complex is equal to $Z_k(X)$ by
Remark \ref{remark-cycles-pointwise}.
\end{enumerate}
The fact that we obtain a presentation of $\CH_k(X)$ follows
immediately by comparing with Definition \ref{definition-rational-equivalence}.
\end{remark}

\noindent
A very simple but important lemma is the following.

\begin{lemma}
\label{lemma-restrict-to-open}
Let $(S, \delta)$ be as in Situation \ref{situation-setup}.
Let $X$ be a scheme locally of finite type over $S$.
Let $U \subset X$ be an open subscheme, and denote
$i : Y = X \setminus U \to X$ as a reduced closed subscheme of $X$.
Let $k \in \mathbf{Z}$.
Suppose $\alpha, \beta \in Z_k(X)$.
If $\alpha|_U \sim_{rat} \beta|_U$ then there exist a cycle
$\gamma \in Z_k(Y)$ such that
$$
\alpha \sim_{rat} \beta + i_*\gamma.
$$
In other words, the sequence
$$
\xymatrix{
\CH_k(Y) \ar[r]^{i_*} & \CH_k(X) \ar[r]^{j^*} & \CH_k(U) \ar[r] & 0
}
$$
is an exact complex of abelian groups.
\end{lemma}

\begin{proof}
Let $\{W_j\}_{j \in J}$ be a locally finite collection of integral closed
subschemes of $U$ of $\delta$-dimension $k + 1$, and let $f_j \in R(W_j)^*$
be elements such that $(\alpha - \beta)|_U = \sum (i_j)_*\text{div}(f_j)$
as in the definition. Set $W_j' \subset X$ equal
to the closure of $W_j$. Suppose that $V \subset X$ is a quasi-compact
open. Then also $V \cap U$ is quasi-compact open in $U$ as
$V$ is Noetherian. Hence the set
$\{j \in J \mid W_j \cap V \not = \emptyset\}
= \{j \in J \mid W'_j \cap V \not = \emptyset\}$
is finite since $\{W_j\}$ is locally finite. In other words we see that
$\{W'_j\}$ is also locally finite. Since $R(W_j) = R(W'_j)$ we see
that
$$
\alpha - \beta - \sum (i'_j)_*\text{div}(f_j)
$$
is a cycle supported on $Y$ and the lemma follows (see
Lemma \ref{lemma-exact-sequence-open}).
\end{proof}

\begin{lemma}
\label{lemma-exact-sequence-closed-chow}
Let $(S, \delta)$ be as in Situation \ref{situation-setup}.
Let $X$ be locally of finite type over $S$. Let $X_1, X_2 \subset X$
be closed subschemes such that $X = X_1 \cup X_2$ set theoretically.
For every $k \in \mathbf{Z}$ the sequence of abelian groups
$$
\xymatrix{
\CH_k(X_1 \cap X_2) \ar[r] &
\CH_k(X_1) \oplus \CH_k(X_2) \ar[r] &
\CH_k(X) \ar[r] &
0
}
$$
is exact. Here $X_1 \cap X_2$ is the scheme theoretic intersection and the
maps are the pushforward maps with one multiplied by $-1$.
\end{lemma}

\begin{proof}
By Lemma \ref{lemma-exact-sequence-closed} the arrow
$\CH_k(X_1) \oplus \CH_k(X_2) \to \CH_k(X)$ is surjective.
Suppose that $(\alpha_1, \alpha_2)$ maps to zero under this map.
Write $\alpha_1 = \sum n_{1, i}[W_{1, i}]$ and
$\alpha_2 = \sum n_{2, i}[W_{2, i}]$. Then we obtain a locally
finite collection $\{W_j\}_{j \in J}$ of integral closed
subschemes of $X$ of $\delta$-dimension $k + 1$ and $f_j \in R(W_j)^*$
such that
$$
\sum n_{1, i}[W_{1, i}] + \sum n_{2, i}[W_{2, i}] = \sum (i_j)_*\text{div}(f_j)
$$
as cycles on $X$ where $i_j : W_j \to X$ is the inclusion morphism.
Choose a disjoint union decomposition $J = J_1 \amalg J_2$ such that
$W_j \subset X_1$ if $j \in J_1$ and $W_j \subset X_2$ if $j \in J_2$.
(This is possible because the $W_j$ are integral.) Then we can write
the equation above as
$$
\sum n_{1, i}[W_{1, i}] - \sum\nolimits_{j \in J_1} (i_j)_*\text{div}(f_j) =
- \sum n_{2, i}[W_{2, i}] + \sum\nolimits_{j \in J_2} (i_j)_*\text{div}(f_j)
$$
Hence this expression is a cycle (!) on $X_1 \cap X_2$. In other words
the element $(\alpha_1, \alpha_2)$ is in the image of the first arrow
and the proof is complete.
\end{proof}

\begin{example}
\label{example-weird}
Here is a ``strange'' example.
Suppose that $S$ is the spectrum of a field $k$
with $\delta$ as in Example \ref{example-field}.
Suppose that $X = C_1 \cup C_2 \cup \ldots$ is an infinite
union of curves $C_j \cong \mathbf{P}^1_k$ glued together
in the following way: The point $\infty \in C_j$ is glued
transversally to the point $0 \in C_{j + 1}$ for $j = 1, 2, 3, \ldots$.
Take the point $0 \in C_1$. This gives a zero cycle
$[0] \in Z_0(X)$. The ``strangeness'' in this situation is
that actually $[0] \sim_{rat} 0$! Namely we can choose
the rational function $f_j \in R(C_j)$ to be the function
which has a simple zero at $0$ and a simple pole at $\infty$
and no other zeros or poles. Then we see that the sum
$\sum (i_j)_*\text{div}(f_j)$ is exactly the $0$-cycle
$[0]$. In fact it turns out that $\CH_0(X) = 0$ in this example.
If you find this too bizarre, then you can just
make sure your spaces are always quasi-compact
(so $X$ does not even exist for you).
\end{example}

\begin{remark}
\label{remark-infinite-sums-rational-equivalences}
Let $(S, \delta)$ be as in Situation \ref{situation-setup}.
Let $X$ be a scheme locally of finite type over $S$.
Suppose we have infinite collections $\alpha_i, \beta_i \in Z_k(X)$,
$i \in I$ of $k$-cycles on $X$. Suppose that the supports
of $\alpha_i$ and $\beta_i$ form locally finite collections
of closed subsets of $X$ so that $\sum \alpha_i$
and $\sum \beta_i$ are defined as cycles. Moreover, assume that
$\alpha_i \sim_{rat} \beta_i$ for each $i$. Then it is not
clear that $\sum \alpha_i \sim_{rat} \sum \beta_i$. Namely,
the problem is that the rational equivalences may be
given by locally finite
families $\{W_{i, j}, f_{i, j} \in R(W_{i, j})^*\}_{j \in J_i}$
but the union $\{W_{i, j}\}_{i \in I, j\in J_i}$ may not
be locally finite.

\medskip\noindent
In many cases in practice, one has a locally finite family of closed
subsets $\{T_i\}_{i \in I}$ such that $\alpha_i, \beta_i$
are supported on $T_i$ and such that $\alpha_i = \beta_i$
in $\CH_k(T_i)$, in other words, the families
$\{W_{i, j}, f_{i, j} \in R(W_{i, j})^*\}_{j \in J_i}$
consist of subschemes $W_{i, j} \subset T_i$. In this case it is true that
$\sum \alpha_i \sim_{rat} \sum \beta_i$ on $X$, simply because
the family $\{W_{i, j}\}_{i \in I, j\in J_i}$ is automatically
locally finite in this case.
\end{remark}






\section{Rational equivalence and push and pull}
\label{section-properties-rational-equivalence}

\noindent
In this section we show that flat pullback and proper pushforward
commute with rational equivalence.

\begin{lemma}
\label{lemma-prepare-flat-pullback-rational-equivalence}
Let $(S, \delta)$ be as in Situation \ref{situation-setup}.
Let $X$, $Y$ be schemes locally of finite type over $S$.
Assume $Y$ integral with $\dim_\delta(Y) = k$.
Let $f : X \to Y$ be a flat morphism of
relative dimension $r$. Then for $g \in R(Y)^*$ we have
$$
f^*\text{div}_Y(g) =
\sum n_j i_{j, *}\text{div}_{X_j}(g \circ f|_{X_j})
$$
as $(k + r - 1)$-cycles on $X$ where the sum is over the irreducible
components $X_j$ of $X$ and $n_j$ is the multiplicity of $X_j$ in $X$.
\end{lemma}

\begin{proof}
Let $Z \subset X$ be an integral closed subscheme of $\delta$-dimension
$k + r - 1$. We have to show that the coefficient $n$ of $[Z]$ in
$f^*\text{div}(g)$ is equal to the coefficient
$m$ of $[Z]$ in $\sum i_{j, *} \text{div}(g \circ f|_{X_j})$.
Let $Z'$ be the closure of $f(Z)$ which is an integral closed
subscheme of $Y$. By Lemma \ref{lemma-flat-inverse-image-dimension}
we have $\dim_\delta(Z') \geq k - 1$. Thus either $Z' = Y$
or $Z'$ is a prime divisor on $Y$. If $Z' = Y$, then the coefficients
$n$ and $m$ are both zero: this is clear for $n$ by definition
of $f^*$ and follows for $m$ because $g \circ f|_{X_j}$ is
a unit in any point of $X_j$ mapping to the generic point of $Y$.
Hence we may assume that $Z' \subset Y$ is a prime divisor.

\medskip\noindent
We are going to translate the equality of $n$ and $m$ into algebra.
Namely, let $\xi' \in Z'$ and $\xi \in Z$ be the generic points.
Set $A = \mathcal{O}_{Y, \xi'}$ and $B = \mathcal{O}_{X, \xi}$.
Note that $A$, $B$ are Noetherian, $A \to B$ is flat, local,
$A$ is a domain, and $\mathfrak m_AB$ is an ideal of definition
of the local ring $B$. The rational function $g$ is an element
of the fraction field $Q(A)$ of $A$.
By construction, the closed subschemes $X_j$
which meet $\xi$ correspond $1$-to-$1$ with minimal primes
$$
\mathfrak q_1, \ldots, \mathfrak q_s \subset B
$$
The integers $n_j$ are the corresponding lengths
$$
n_i = \text{length}_{B_{\mathfrak q_i}}(B_{\mathfrak q_i})
$$
The rational functions $g \circ f|_{X_j}$ correspond to the image
$g_i \in \kappa(\mathfrak q_i)^*$ of $g \in Q(A)$.
Putting everything together we see that
$$
n = \text{ord}_A(g) \text{length}_B(B/\mathfrak m_AB)
$$
and that
$$
m = \sum \text{ord}_{B/\mathfrak q_i}(g_i)
\text{length}_{B_{\mathfrak q_i}}(B_{\mathfrak q_i})
$$
Writing $g = x/y$ for some nonzero $x, y \in A$ we see that it suffices
to prove
$$
\text{length}_A(A/(x)) \text{length}_B(B/\mathfrak m_AB) =
\text{length}_B(B/xB)
$$
(equality uses Algebra, Lemma \ref{algebra-lemma-pullback-module})
equals
$$
\sum\nolimits_{i = 1, \ldots, s}
\text{length}_{B/\mathfrak q_i}(B/(x, \mathfrak q_i))
\text{length}_{B_{\mathfrak q_i}}(B_{\mathfrak q_i})
$$
and similarly for $y$. As $A \to B$ is flat it follows that $x$
is a nonzerodivisor in $B$. Hence the desired equality follows from
Lemma \ref{lemma-additivity-divisors-restricted}.
\end{proof}

\begin{lemma}
\label{lemma-flat-pullback-rational-equivalence}
Let $(S, \delta)$ be as in Situation \ref{situation-setup}.
Let $X$, $Y$ be schemes locally of finite type over $S$.
Let $f : X \to Y$ be a flat morphism of relative dimension $r$.
Let $\alpha \sim_{rat} \beta$ be rationally equivalent $k$-cycles on $Y$.
Then $f^*\alpha \sim_{rat} f^*\beta$ as $(k + r)$-cycles on $X$.
\end{lemma}

\begin{proof}
What do we have to show? Well, suppose we are given a collection
$$
i_j : W_j \longrightarrow Y
$$
of closed immersions, with each $W_j$ integral of $\delta$-dimension $k + 1$
and rational functions $g_j \in R(W_j)^*$. Moreover, assume that
the collection $\{i_j(W_j)\}_{j \in J}$ is locally finite on $Y$.
Then we have to show that
$$
f^*(\sum i_{j, *}\text{div}(g_j)) = \sum f^*i_{j, *}\text{div}(g_j)
$$
is rationally equivalent to zero on $X$. The sum on the right
makes sense as $\{W_j\}$ is locally finite in $X$ by
Lemma \ref{lemma-inverse-image-locally-finite}.

\medskip\noindent
Consider the fibre products
$$
i'_j : W'_j = W_j \times_Y X \longrightarrow X.
$$
and denote $f_j : W'_j \to W_j$ the first projection.
By Lemma \ref{lemma-flat-pullback-proper-pushforward}
we can write the sum above as
$$
\sum i'_{j, *}(f_j^*\text{div}(g_j))
$$
By Lemma \ref{lemma-prepare-flat-pullback-rational-equivalence}
we see that each $f_j^*\text{div}(g_j)$ is rationally equivalent
to zero on $W'_j$. Hence each $i'_{j, *}(f_j^*\text{div}(g_j))$
is rationally equivalent to zero. Then the same is true for
the displayed sum by the discussion in
Remark \ref{remark-infinite-sums-rational-equivalences}.
\end{proof}

\begin{lemma}
\label{lemma-proper-pushforward-rational-equivalence}
Let $(S, \delta)$ be as in Situation \ref{situation-setup}.
Let $X$, $Y$ be schemes locally of finite type over $S$.
Let $p : X \to Y$ be a proper morphism.
Suppose $\alpha, \beta \in Z_k(X)$ are rationally equivalent.
Then $p_*\alpha$ is rationally equivalent to $p_*\beta$.
\end{lemma}

\begin{proof}
What do we have to show? Well, suppose we are given a collection
$$
i_j : W_j \longrightarrow X
$$
of closed immersions, with each $W_j$ integral of $\delta$-dimension $k + 1$
and rational functions $f_j \in R(W_j)^*$.
Moreover, assume that
the collection $\{i_j(W_j)\}_{j \in J}$ is locally finite on $X$.
Then we have to show that
$$
p_*\left(\sum i_{j, *}\text{div}(f_j)\right)
$$
is rationally equivalent to zero on $X$.

\medskip\noindent
Note that the sum is equal to
$$
\sum p_*i_{j, *}\text{div}(f_j).
$$
Let $W'_j \subset Y$ be the integral closed subscheme which is the
image of $p \circ i_j$. The collection $\{W'_j\}$ is locally finite
in $Y$ by Lemma \ref{lemma-quasi-compact-locally-finite}.
Hence it suffices to show, for a given $j$, that either
$p_*i_{j, *}\text{div}(f_j) = 0$ or that it
is equal to $i'_{j, *}\text{div}(g_j)$ for some $g_j \in R(W'_j)^*$.

\medskip\noindent
The arguments above therefore reduce us to the case of a single
integral closed subscheme $W \subset X$ of $\delta$-dimension $k + 1$.
Let $f \in R(W)^*$. Let $W' = p(W)$ as above.
We get a commutative diagram of morphisms
$$
\xymatrix{
W \ar[r]_i \ar[d]_{p'} & X \ar[d]^p \\
W' \ar[r]^{i'} & Y
}
$$
Note that $p_*i_*\text{div}(f) = i'_*(p')_*\text{div}(f)$ by
Lemma \ref{lemma-compose-pushforward}. As explained above
we have to show that $(p')_*\text{div}(f)$
is the divisor of a rational function on $W'$ or zero.
There are three cases to distinguish.

\medskip\noindent
The case $\dim_\delta(W') < k$. In this case automatically
$(p')_*\text{div}(f) = 0$ and there is nothing to prove.

\medskip\noindent
The case $\dim_\delta(W') = k$. Let us show that $(p')_*\text{div}(f) = 0$
in this case. Let $\eta \in W'$ be the generic point.
Note that $c : W_\eta \to \Spec(K)$
is a proper integral curve over $K = \kappa(\eta)$
whose function field $K(W_\eta)$ is identified with $R(W)$.
Here is a diagram
$$
\xymatrix{
W_\eta \ar[r] \ar[d]_c & W \ar[d]^{p'} \\
\Spec(K) \ar[r] & W'
}
$$
Let us denote $f_\eta \in K(W_\eta)^*$ the rational function
corresponding to $f \in R(W)^*$.
Moreover, the closed points $\xi$ of $W_\eta$ correspond $1 - 1$ to the
closed integral subschemes $Z = Z_\xi \subset W$ of $\delta$-dimension $k$
with $p'(Z) = W'$. Note that the multiplicity
of $Z_\xi$ in $\text{div}(f)$ is equal to
$\text{ord}_{\mathcal{O}_{W_\eta, \xi}}(f_\eta)$ simply because the
local rings $\mathcal{O}_{W_\eta, \xi}$ and $\mathcal{O}_{W, \xi}$
are identified (as subrings of their fraction fields).
Hence we see that the multiplicity of $[W']$ in
$(p')_*\text{div}(f)$ is equal to the multiplicity of
$[\Spec(K)]$ in $c_*\text{div}(f_\eta)$.
By Lemma \ref{lemma-curve-principal-divisor} this is zero.

\medskip\noindent
The case $\dim_\delta(W') = k + 1$. In this case
Lemma \ref{lemma-proper-pushforward-alteration} applies,
and we see that indeed $p'_*\text{div}(f) = \text{div}(g)$
for some $g \in R(W')^*$ as desired.
\end{proof}















\section{Rational equivalence and the projective line}
\label{section-different-rational-equivalence}

\noindent
Let $(S, \delta)$ be as in Situation \ref{situation-setup}.
Let $X$ be a scheme locally of finite type over $S$.
Given any closed subscheme
$Z \subset X \times_S \mathbf{P}^1_S = X \times \mathbf{P}^1$
we let $Z_0$, resp.\ $Z_\infty$ be the scheme theoretic closed
subscheme $Z_0 = \text{pr}_2^{-1}(D_0)$,
resp.\ $Z_\infty = \text{pr}_2^{-1}(D_\infty)$.
Here $D_0$, $D_\infty$ are as in (\ref{equation-zero-infty}).

\begin{lemma}
\label{lemma-rational-equivalence-family}
Let $(S, \delta)$ be as in Situation \ref{situation-setup}.
Let $X$ be a scheme locally of finite type over $S$.
Let $W \subset X \times_S \mathbf{P}^1_S$ be an integral
closed subscheme of $\delta$-dimension $k + 1$.
Assume $W \not = W_0$, and $W \not = W_\infty$. Then
\begin{enumerate}
\item $W_0$, $W_\infty$ are effective Cartier divisors of $W$,
\item $W_0$, $W_\infty$ can be viewed as closed subschemes
of $X$ and
$$
[W_0]_k \sim_{rat} [W_\infty]_k,
$$
\item for any locally finite family of
integral closed subschemes
$W_i \subset X \times_S \mathbf{P}^1_S$
of $\delta$-dimension $k + 1$ with $W_i \not = (W_i)_0$ and
$W_i \not = (W_i)_\infty$ we have
$\sum ([(W_i)_0]_k - [(W_i)_\infty]_k) \sim_{rat} 0$
on $X$, and
\item for any $\alpha \in Z_k(X)$ with $\alpha \sim_{rat} 0$
there exists a locally finite family of
integral closed subschemes $W_i \subset X \times_S \mathbf{P}^1_S$
as above such that $\alpha = \sum ([(W_i)_0]_k - [(W_i)_\infty]_k)$.
\end{enumerate}
\end{lemma}

\begin{proof}
Part (1) follows from
Divisors, Lemma \ref{divisors-lemma-pullback-effective-Cartier-defined}
since the generic point
of $W$ is not mapped into $D_0$ or $D_\infty$ under the projection
$X \times_S \mathbf{P}^1_S \to \mathbf{P}^1_S$ by assumption.

\medskip\noindent
Since $X \times_S D_0 \to X$ is a closed immersion, we see that $W_0$
is isomorphic to a closed subscheme of $X$. Similarly for $W_\infty$.
The morphism $p : W \to X$ is proper as a composition of
the closed immersion $W \to X \times_S \mathbf{P}^1_S$ and the
proper morphism $X \times_S \mathbf{P}^1_S \to X$. By
Lemma \ref{lemma-rational-function} we have
$[W_0]_k \sim_{rat} [W_\infty]_k$ as cycles on $W$. Hence part (2) follows from
Lemma \ref{lemma-proper-pushforward-rational-equivalence} as clearly
$p_*[W_0]_k = [W_0]_k$ and similarly for $W_\infty$.

\medskip\noindent
The only content of statement (3) is, given parts (1) and (2), that
the collection $\{(W_i)_0, (W_i)_\infty\}$ is a locally finite collection
of closed subschemes of $X$. This is clear.

\medskip\noindent
Suppose that $\alpha \sim_{rat} 0$.
By definition this means there exist integral closed subschemes
$V_i \subset X$ of $\delta$-dimension $k + 1$ and rational
functions $f_i \in R(V_i)^*$ such that the family
$\{V_i\}_{i \in I}$ is locally finite in $X$ and such that
$\alpha = \sum (V_i \to X)_*\text{div}(f_i)$.
Let
$$
W_i \subset V_i \times_S \mathbf{P}^1_S \subset X \times_S \mathbf{P}^1_S
$$
be the closure of the graph of the rational map $f_i$ as in
Lemma \ref{lemma-rational-function}.
Then we have that $(V_i \to X)_*\text{div}(f_i)$
is equal to $[(W_i)_0]_k - [(W_i)_\infty]_k$ by that same lemma.
Hence the result is clear.
\end{proof}

\begin{lemma}
\label{lemma-closed-subscheme-cross-p1}
Let $(S, \delta)$ be as in Situation \ref{situation-setup}.
Let $X$ be a scheme locally of finite type over $S$.
Let $Z$ be a closed subscheme of $X \times \mathbf{P}^1$.
Assume
\begin{enumerate}
\item $\dim_\delta(Z) \leq k + 1$,
\item $\dim_\delta(Z_0) \leq k$, $\dim_\delta(Z_\infty) \leq k$, and
\item for any embedded point $\xi$ (Divisors, Definition
\ref{divisors-definition-embedded}) of $Z$ either
$\xi \not \in Z_0 \cup Z_\infty$ or $\delta(\xi) < k$.
\end{enumerate}
Then $[Z_0]_k \sim_{rat} [Z_\infty]_k$ as $k$-cycles on $X$.
\end{lemma}

\begin{proof}
Let $\{W_i\}_{i \in I}$ be the collection of irreducible
components of $Z$ which have $\delta$-dimension $k + 1$.
Write
$$
[Z]_{k + 1} = \sum n_i[W_i]
$$
with $n_i > 0$ as per definition. Note that $\{W_i\}$
is a locally finite collection of closed subsets of
$X \times_S \mathbf{P}^1_S$ by
Divisors, Lemma \ref{divisors-lemma-components-locally-finite}.
We claim that
$$
[Z_0]_k = \sum n_i[(W_i)_0]_k
$$
and similarly for $[Z_\infty]_k$. If we prove this then the lemma
follows from Lemma \ref{lemma-rational-equivalence-family}.

\medskip\noindent
Let $Z' \subset X$ be an integral closed subscheme of $\delta$-dimension $k$.
To prove the equality above it suffices to show that the coefficient $n$
of $[Z']$ in $[Z_0]_k$ is the same as the coefficient $m$ of
$[Z']$ in $\sum n_i[(W_i)_0]_k$. Let $\xi' \in Z'$ be the generic point.
Set $\xi = (\xi', 0) \in  X \times_S \mathbf{P}^1_S$.
Consider the local ring $A = \mathcal{O}_{X \times_S \mathbf{P}^1_S, \xi}$.
Let $I \subset A$ be the ideal cutting out $Z$, in other words so that
$A/I = \mathcal{O}_{Z, \xi}$. Let $t \in A$ be the element cutting
out $X \times_S D_0$ (i.e., the coordinate of $\mathbf{P}^1$ at zero
pulled back). By our choice of $\xi' \in Z'$ we have $\delta(\xi) = k$
and hence $\dim(A/I) = 1$. Since $\xi$ is not an embedded point by
assumption (3) we see that $A/I$ is Cohen-Macaulay. Since $\dim_\delta(Z_0)
= k$ we see that $\dim(A/(t, I)) = 0$ which implies that $t$
is a nonzerodivisor on $A/I$. Finally, the irreducible closed subschemes
$W_i$ passing through $\xi$ correspond to the minimal primes
$I \subset \mathfrak q_i$ over $I$. The multiplicities $n_i$ correspond
to the lengths $\text{length}_{A_{\mathfrak q_i}}(A/I)_{\mathfrak q_i}$.
Hence we see that
$$
n = \text{length}_A(A/(t, I))
$$
and
$$
m = \sum
\text{length}_A(A/(t, \mathfrak q_i))
\text{length}_{A_{\mathfrak q_i}}(A/I)_{\mathfrak q_i}
$$
Thus the result follows from
Lemma \ref{lemma-additivity-divisors-restricted}.
\end{proof}

\begin{lemma}
\label{lemma-coherent-sheaf-cross-p1}
Let $(S, \delta)$ be as in Situation \ref{situation-setup}.
Let $X$ be a scheme locally of finite type over $S$.
Let $\mathcal{F}$ be a coherent sheaf on $X \times \mathbf{P}^1$.
Let $i_0, i_\infty : X \to X \times \mathbf{P}^1$ be the closed immersion
such that $i_t(x) = (x, t)$. Denote $\mathcal{F}_0 = i_0^*\mathcal{F}$ and
$\mathcal{F}_\infty = i_\infty^*\mathcal{F}$.
Assume
\begin{enumerate}
\item $\dim_\delta(\text{Supp}(\mathcal{F})) \leq k + 1$,
\item $\dim_\delta(\text{Supp}(\mathcal{F}_0)) \leq k$,
$\dim_\delta(\text{Supp}(\mathcal{F}_\infty)) \leq k$, and
\item for any embedded associated point $\xi$ of $\mathcal{F}$ either
$\xi \not \in (X \times \mathbf{P}^1)_0 \cup (X \times \mathbf{P}^1)_\infty$
or $\delta(\xi) < k$.
\end{enumerate}
Then $[\mathcal{F}_0]_k \sim_{rat} [\mathcal{F}_\infty]_k$ as $k$-cycles on $X$.
\end{lemma}

\begin{proof}
Let $\{W_i\}_{i \in I}$ be the collection of irreducible
components of $\text{Supp}(\mathcal{F})$
which have $\delta$-dimension $k + 1$.
Write
$$
[\mathcal{F}]_{k + 1} = \sum n_i[W_i]
$$
with $n_i > 0$ as per definition. Note that $\{W_i\}$
is a locally finite collection of closed subsets of
$X \times_S \mathbf{P}^1_S$ by Lemma \ref{lemma-length-finite}.
We claim that
$$
[\mathcal{F}_0]_k = \sum n_i[(W_i)_0]_k
$$
and similarly for $[\mathcal{F}_\infty]_k$. If we prove this then the lemma
follows from Lemma \ref{lemma-rational-equivalence-family}.

\medskip\noindent
Let $Z' \subset X$ be an integral closed subscheme of $\delta$-dimension $k$.
To prove the equality above it suffices to show that the coefficient $n$
of $[Z']$ in $[\mathcal{F}_0]_k$ is the same as the coefficient $m$ of
$[Z']$ in $\sum n_i[(W_i)_0]_k$. Let $\xi' \in Z'$ be the generic point.
Set $\xi = (\xi', 0) \in  X \times_S \mathbf{P}^1_S$.
Consider the local ring $A = \mathcal{O}_{X \times_S \mathbf{P}^1_S, \xi}$.
Let $M = \mathcal{F}_\xi$ as an $A$-module.
Let $t \in A$ be the element cutting out $X \times_S D_0$
(i.e., the coordinate of $\mathbf{P}^1$ at zero pulled back).
By our choice of $\xi' \in Z'$ we have $\delta(\xi) = k$
and hence $\dim(\text{Supp}(M)) = 1$. Since $\xi$ is not an associated point
of $\mathcal{F}$ by assumption (3) we see that $M$ is a Cohen-Macaulay module.
Since $\dim_\delta(\text{Supp}(\mathcal{F}_0)) = k$
we see that $\dim(\text{Supp}(M/tM)) = 0$ which implies that $t$
is a nonzerodivisor on $M$. Finally, the irreducible closed subschemes
$W_i$ passing through $\xi$ correspond to the minimal primes
$\mathfrak q_i$ of $\text{Ass}(M)$. The multiplicities $n_i$ correspond
to the lengths $\text{length}_{A_{\mathfrak q_i}}M_{\mathfrak q_i}$.
Hence we see that
$$
n = \text{length}_A(M/tM)
$$
and
$$
m = \sum
\text{length}_A(A/(t, \mathfrak q_i)A)
\text{length}_{A_{\mathfrak q_i}}M_{\mathfrak q_i}
$$
Thus the result follows from
Lemma \ref{lemma-additivity-divisors-restricted}.
\end{proof}







\section{Chow groups and envelopes}
\label{section-envelopes}

\noindent
Here is the definition.

\begin{definition}
\label{definition-envelope}
\begin{reference}
\cite[Definition 18.3]{F}
\end{reference}
Let $X$ be a scheme. An {\it envelope} is a proper morphism $f : Y \to X$
which is completely decomposed
(More on Morphisms, Definition \ref{more-morphisms-definition-cd-morphism}).
\end{definition}

\noindent
The exact sequence of Lemma \ref{lemma-envelope}
is the main motivation for the definition.

\begin{lemma}
\label{lemma-composition-envelope}
Let $(S, \delta)$ be as in Situation \ref{situation-setup}.
Let $X$ be a scheme locally of finite type over $S$.
If $f : Y \to X$ and $g : Z \to Y$ are envelopes, then
$f \circ g$ is an envelope.
\end{lemma}

\begin{proof}
Follows from Morphisms, Lemma \ref{morphisms-lemma-composition-proper}
and More on Morphisms, Lemma \ref{more-morphisms-lemma-composition-cd}.
\end{proof}

\begin{lemma}
\label{lemma-base-change-envelope}
Let $(S, \delta)$ be as in Situation \ref{situation-setup}.
Let $X' \to X$ be a morphism of schemes locally of finite type over $S$.
If $f : Y \to X$ is an envelope, then the base change $f' : Y' \to X'$
of $f$ is an envelope too.
\end{lemma}

\begin{proof}
Follows from Morphisms, Lemma \ref{morphisms-lemma-base-change-proper}
and More on Morphisms, Lemma \ref{more-morphisms-lemma-base-change-cd}.
\end{proof}

\begin{lemma}
\label{lemma-envelope}
Let $(S, \delta)$ be as in Situation \ref{situation-setup}.
Let $X$ be a scheme locally of finite type over $S$.
Let $f : Y \to X$ be an envelope. Then
we have an exact sequence
$$
\CH_k(Y \times_X Y) \xrightarrow{p_* - q_*}
\CH_k(Y) \xrightarrow{f_*}
\CH_k(X) \to 0
$$
for all $k \in \mathbf{Z}$. Here $p, q : Y \times_X Y \to Y$ are
the projections.
\end{lemma}

\begin{proof}
Since $f$ is an envelope, $f$ is proper and hence pushforward on
cycles and cycle classes is defined, see
Sections \ref{section-proper-pushforward} and \ref{section-push-pull}.
Similarly, the morphisms $p$ and $q$ are proper as base changes of $f$.
The composition of the arrows is zero as
$f_* \circ p_* = (p \circ f)_* = (q \circ f)_* = f_* \circ q_*$, see
Lemma \ref{lemma-compose-pushforward}.

\medskip\noindent
Let us show that $f_* : Z_k(Y) \to Z_k(X)$ is surjective.
Namely, suppose that we have $\alpha = \sum n_i[Z_i] \in Z_k(X)$
where $Z_i \subset X$ is a locally finite family of integral
closed subschemes. Let $x_i \in Z_i$ be the generic point.
Since $f$ is an envelope and hence completely decomposed,
there exists a point $y_i \in Y$ with $f(y_i) = x_i$
and with $\kappa(y_i)/\kappa(x_i)$ trivial. Let $W_i \subset Y$
be the integral closed subscheme with generic point $y_i$.
Since $f$ is closed, we see that $f(W_i) = Z_i$.
It follows that the family of closed subschemes $W_i$ is locally finite
on $Y$. Since $\kappa(y_i)/\kappa(x_i)$ is trivial we see
that $\dim_\delta(W_i) = \dim_\delta(Z_i) = k$. Hence
$\beta = \sum n_i[W_i]$ is in $Z_k(Y)$. Finally, since
$\kappa(y_i)/\kappa(x_i)$ is trivial, the degree of the dominant
morphism $f|_{W_i} : W_i \to Z_i$ is $1$ and we conclude
that $f_*\beta = \alpha$.

\medskip\noindent
Since $f_* : Z_k(Y) \to Z_k(X)$ is surjective, a fortiori the map
$f_* : \CH_k(Y) \to \CH_k(X)$ is surjective.

\medskip\noindent
Let $\beta \in Z_k(Y)$ be an element such that $f_*\beta$ is zero in
$\CH_k(X)$. This means we can find a locally finite family of
integral closed subschemes $Z_j \subset X$ with $\dim_\delta(Z_j) = k + 1$
and $f_j \in R(Z_j)^*$ such that
$$
f_*\beta = \sum (Z_j \to X)_*\text{div}(f_j)
$$
as cycles where $i_j : Z_j \to X$ is the given closed immersion.
Arguing exactly as above, we can find a locally finite
family of integral closed subschemes $W_j \subset Y$
with $f(W_j) = Z_j$ and such that $W_j \to Z_j$ is birational, i.e.,
induces an isomorphism $R(Z_j) = R(W_j)$. Denote $g_j \in R(W_j)^*$
the element corresponding to $f_j$. Observe that $W_j \to Z_j$
is proper and that $(W_j \to Z_j)_*\text{div}(g_j) = \text{div}(f_j)$
as cycles on $Z_j$. It follows from this that if we replace
$\beta$ by the rationally equivalent cycle
$$
\beta' = \beta - \sum (W_j \to Y)_*\text{div}(g_j)
$$
then we find that $f_*\beta' = 0$.
(This uses Lemma \ref{lemma-compose-pushforward}.)
Thus to finish the proof
of the lemma it suffices to show the claim in the following paragraph.

\medskip\noindent
Claim: if $\beta \in Z_k(Y)$ and $f_*\beta = 0$, then
$\beta = \delta + p_*\gamma - q_*\gamma$ in $Z_k(Y)$ for some
$\gamma \in Z_k(Y \times_X Y)$. Namely, write $\beta = \sum_{j \in J} n_j[W_j]$
with $\{W_j\}_{j \in J}$ a locally finite family of integral closed
subschemes of $Y$ with $\dim_\delta(W_j) = k$.
Fix an integral closed subscheme $Z \subset X$. Consider the subset
$J_Z = \{j \in J : f(W_j) = Z\}$. This is a finite set. There are three
cases:
\begin{enumerate}
\item $J_Z = \emptyset$. In this case we set $\gamma_Z = 0$.
\item $J_Z \not = \emptyset$ and $\dim_\delta(Z) = k$.
The condition $f_*\beta = 0$ implies by looking at the
coefficient of $Z$ that $\sum_{j \in J_Z} n_j\deg(W_j/Z) = 0$.
In this case we choose an integral closed subscheme $W \subset Y$
which maps birationally onto $Z$ (see above). Looking at generic
points, we see that $W_j \times_Z W$ has a unique irreducible
component $W'_j \subset W_j \times_Z W \subset Y \times_X Y$
mapping birationally to $W_j$. Then $W'_j \to W$ is dominant
and $\deg(W'_j/W) = \deg(W_j/W)$. Thus if we set
$\gamma_Z = \sum_{j \in J_Z} n_j[W'_j]$
then we see that
$p_*\gamma_Z = \sum_{j \in J_Z} n_j[W_j]$ and
$q_*\gamma_Z = \sum_{j \in J_Z} n_j\deg(W'_j/W)[W] = 0$.
\item $J_Z \not = \emptyset$ and $\dim_\delta(Z) < k$.
In this case we choose an integral closed subscheme $W \subset Y$
which maps birationally onto $Z$ (see above). Looking at generic
points, we see that $W_j \times_Z W$ has a unique irreducible
component $W'_j \subset W_j \times_Z W \subset Y \times_X Y$
mapping birationally to $W_j$. Then $W'_j \to W$ is dominant
and $k = \dim_\delta(W'_j) > \dim_\delta(W) = \dim_\delta(Z)$.
Thus if we set $\gamma_Z = \sum_{j \in J_Z} n_j[W'_j]$
then we see that
$p_*\gamma_Z = \sum_{j \in J_Z} n_j[W_j]$ and
$q_*\gamma_Z = 0$.
\end{enumerate}
Since the family of integral closed subschemes $\{f(W_j)\}$
is locally finite on $X$
(Lemma \ref{lemma-quasi-compact-locally-finite})
we see that the $k$-cycle
$$
\gamma = \sum\nolimits_{Z \subset X\text{ integral closed}} \gamma_Z
$$
on $Y \times_X Y$ is well defined. By our computations above it follows that
$p_*\gamma_Z = \beta$ and $q_*\gamma_Z = 0$ which implies
what we wanted to prove.
\end{proof}













\section{Chow groups and K-groups}
\label{section-chow-and-K}

\noindent
In this section we are going to compare $K_0$ of the
category of coherent sheaves to the chow groups.

\medskip\noindent
Let $(S, \delta)$ be as in Situation \ref{situation-setup}.
Let $X$ be a scheme locally of finite type over $S$.
We denote $\textit{Coh}(X) = \textit{Coh}(\mathcal{O}_X)$
the category of coherent sheaves on $X$.
It is an abelian category, see
Cohomology of Schemes, Lemma \ref{coherent-lemma-coherent-abelian-Noetherian}.
For any $k \in \mathbf{Z}$ we let $\textit{Coh}_{\leq k}(X)$
be the full subcategory of $\textit{Coh}(X)$
consisting of those coherent sheaves $\mathcal{F}$
having $\dim_\delta(\text{Supp}(\mathcal{F})) \leq k$.

\begin{lemma}
\label{lemma-Serre-subcategories}
Let $(S, \delta)$ be as in Situation \ref{situation-setup}.
Let $X$ be a scheme locally of finite type over $S$.
The categories $\textit{Coh}_{\leq k}(X)$ are Serre subcategories
of the abelian category $\textit{Coh}(X)$.
\end{lemma}

\begin{proof}
The definition of a Serre subcategory is
Homology, Definition \ref{homology-definition-serre-subcategory}.
The proof of the lemma is straightforward and omitted.
\end{proof}

\begin{lemma}
\label{lemma-cycles-k-group}
Let $(S, \delta)$ be as in Situation \ref{situation-setup}.
Let $X$ be a scheme locally of finite type over $S$.
The maps
$$
Z_k(X)
\longrightarrow
K_0(\textit{Coh}_{\leq k}(X)/\textit{Coh}_{\leq k - 1}(X)),
\quad
\sum n_Z[Z] \mapsto
\left[\bigoplus\nolimits_{n_Z > 0} \mathcal{O}_Z^{\oplus n_Z}\right]
-
\left[\bigoplus\nolimits_{n_Z < 0} \mathcal{O}_Z^{\oplus -n_Z}\right]
$$
and
$$
K_0(\textit{Coh}_{\leq k}(X)/\textit{Coh}_{\leq k - 1}(X))
\longrightarrow
Z_k(X),\quad
\mathcal{F} \longmapsto [\mathcal{F}]_k
$$
are mutually inverse isomorphisms.
\end{lemma}

\begin{proof}
Note that if $\sum n_Z[Z]$ is in $Z_k(X)$, then
the direct sums
$\bigoplus\nolimits_{n_Z > 0} \mathcal{O}_Z^{\oplus n_Z}$ and
$\bigoplus\nolimits_{n_Z < 0} \mathcal{O}_Z^{\oplus -n_Z}$
are coherent sheaves on $X$ since the family $\{Z \mid n_Z > 0\}$
is locally finite on $X$.
The map $\mathcal{F} \to [\mathcal{F}]_k$ is additive
on $\textit{Coh}_{\leq k}(X)$, see
Lemma \ref{lemma-additivity-sheaf-cycle}. And $[\mathcal{F}]_k = 0$
if $\mathcal{F} \in \textit{Coh}_{\leq k - 1}(X)$. By part (1)
of Homology, Lemma \ref{homology-lemma-serre-subcategory-K-groups}
this implies that the second map is well defined too.
It is clear that the composition of the first map with the second
map is the identity.

\medskip\noindent
Conversely, say we start with a coherent sheaf $\mathcal{F}$
on $X$. Write $[\mathcal{F}]_k = \sum_{i \in I} n_i[Z_i]$
with $n_i > 0$ and $Z_i \subset X$, $i \in I$
pairwise distinct integral closed subschemes of $\delta$-dimension $k$.
We have to show that
$$
[\mathcal{F}] = [\bigoplus\nolimits_{i \in I} \mathcal{O}_{Z_i}^{\oplus n_i}]
$$
in $K_0(\textit{Coh}_{\leq k}(X)/\textit{Coh}_{\leq k - 1}(X))$.
Denote $\xi_i \in Z_i$ the generic point.
If we set
$$
\mathcal{F}' = \Ker(\mathcal{F} \to \bigoplus \xi_{i, *}\mathcal{F}_{\xi_i})
$$
then $\mathcal{F}'$ is the maximal coherent submodule of $\mathcal{F}$
whose support has dimension $\leq k - 1$. In particular $\mathcal{F}$
and $\mathcal{F}/\mathcal{F}'$ have the same class in
$K_0(\textit{Coh}_{\leq k}(X)/\textit{Coh}_{\leq k - 1}(X))$.
Thus after replacing $\mathcal{F}$ by $\mathcal{F}/\mathcal{F}'$
we may and do assume that the kernel $\mathcal{F}'$ displayed
above is zero.

\medskip\noindent
For each $i \in I$ we choose a filtration
$$
\mathcal{F}_{\xi_i} = \mathcal{F}_i^0 \supset \mathcal{F}_i^1 \supset
\ldots \supset \mathcal{F}_i^{n_i} = 0
$$
such that the successive quotients are of dimension $1$ over the residue
field at $\xi_i$. This is possible as the length of $\mathcal{F}_{\xi_i}$
over $\mathcal{O}_{X, \xi_i}$ is $n_i$.
For $p > n_i$ set $\mathcal{F}_i^p = 0$. For $p \geq 0$ we denote
$$
\mathcal{F}^p =
\Ker\left(\mathcal{F} \longrightarrow \bigoplus
\xi_{i, *}(\mathcal{F}_{\xi_i}/\mathcal{F}_i^p)\right)
$$
Then $\mathcal{F}^p$ is coherent, $\mathcal{F}^0 = \mathcal{F}$, and
$\mathcal{F}^p/\mathcal{F}^{p + 1}$ is isomorphic to a free
$\mathcal{O}_{Z_i}$-module of rank $1$ (if $n_i > p$) or $0$
(if $n_i \leq p$) in an open neighbourhood of $\xi_i$. Moreover,
$\mathcal{F}' = \bigcap \mathcal{F}^p = 0$. Since every quasi-compact
open $U \subset X$ contains only a finite number of $\xi_i$
we conclude that $\mathcal{F}^p|_U$ is zero for $p \gg 0$.
Hence $\bigoplus_{p \geq 0} \mathcal{F}^p$ is a coherent
$\mathcal{O}_X$-module. Consider the short exact sequences
$$
0 \to
\bigoplus\nolimits_{p > 0} \mathcal{F}^p \to
\bigoplus\nolimits_{p \geq 0} \mathcal{F}^p \to
\bigoplus\nolimits_{p > 0} \mathcal{F}^p/\mathcal{F}^{p + 1} \to 0
$$
and
$$
0 \to
\bigoplus\nolimits_{p > 0} \mathcal{F}^p \to
\bigoplus\nolimits_{p \geq 0} \mathcal{F}^p \to
\mathcal{F} \to 0
$$
of coherent $\mathcal{O}_X$-modules. This already shows that
$$
[\mathcal{F}] = [\bigoplus \mathcal{F}^p/\mathcal{F}^{p + 1}]
$$
in $K_0(\textit{Coh}_{\leq k}(X)/\textit{Coh}_{\leq k - 1}(X))$.
Next, for every $p \geq 0$ and $i \in I$ such that $n_i > p$
we choose a nonzero ideal sheaf $\mathcal{I}_{i, p} \subset \mathcal{O}_{Z_i}$
and a map $\mathcal{I}_{i, p} \to \mathcal{F}^p/\mathcal{F}^{p + 1}$ on $X$
which is an isomorphism over the open neighbourhood of $\xi_i$
mentioned above. This is possible by
Cohomology of Schemes, Lemma \ref{coherent-lemma-extend-coherent}.
Then we consider the short exact sequence
$$
0 \to
\bigoplus\nolimits_{p \geq 0, i \in I, n_i > p} \mathcal{I}_{i, p}
\to
\bigoplus \mathcal{F}^p/\mathcal{F}^{p + 1} \to
\mathcal{Q} \to 0
$$
and the short exact sequence
$$
0 \to
\bigoplus\nolimits_{p \geq 0, i \in I, n_i > p} \mathcal{I}_{i, p}
\to
\bigoplus\nolimits_{p \geq 0, i \in I, n_i > p} \mathcal{O}_{Z_i}
\to
\mathcal{Q}' \to 0
$$
Observe that both $\mathcal{Q}$ and $\mathcal{Q}'$ are zero in a neighbourhood
of the points $\xi_i$ and that they are supported on $\bigcup Z_i$.
Hence $\mathcal{Q}$ and $\mathcal{Q}'$ are in
$\textit{Coh}_{\leq k - 1}(X)$.
Since
$$
\bigoplus\nolimits_{i \in I} \mathcal{O}_{Z_i}^{\oplus n_i} \cong
\bigoplus\nolimits_{p \geq 0, i \in I, n_i > p} \mathcal{O}_{Z_i}
$$
this concludes the proof.
\end{proof}

\begin{lemma}
\label{lemma-finite-cycles-k-group}
Let $\pi : X \to Y$ be a finite morphism of schemes locally of finite type
over $(S, \delta)$ as in Situation \ref{situation-setup}. Then
$\pi_* : \textit{Coh}(X) \to \textit{Coh}(Y)$ is an exact functor
which sends $\textit{Coh}_{\leq k}(X)$ into $\textit{Coh}_{\leq k}(Y)$
and induces homomorphisms on $K_0$ of these categories and
their quotients. The maps of Lemma \ref{lemma-cycles-k-group}
fit into a commutative diagram
$$
\xymatrix{
Z_k(X) \ar[d]^{\pi_*} \ar[r] &
K_0(\textit{Coh}_{\leq k}(X)/\textit{Coh}_{\leq k - 1}(X))
\ar[d]^{\pi_*} \ar[r] &
Z_k(X) \ar[d]^{\pi_*} \\
Z_k(Y) \ar[r] &
K_0(\textit{Coh}_{\leq k}(Y)/\textit{Coh}_{\leq k - 1}(Y)) \ar[r] &
Z_k(Y)
}
$$
\end{lemma}

\begin{proof}
A finite morphism is affine, hence pushforward of quasi-coherent
modules along $\pi$ is an exact functor by
Cohomology of Schemes, Lemma \ref{coherent-lemma-relative-affine-vanishing}.
A finite morphism is proper, hence $\pi_*$ sends coherent sheaves
to coherent sheaves, see Cohomology of Schemes, Proposition
\ref{coherent-proposition-proper-pushforward-coherent}.
The statement on dimensions of supports is clear.
Commutativity on the right follows immediately from
Lemma \ref{lemma-cycle-push-sheaf}.
Since the horizontal arrows are bijections, we find that
we have commutativity on the left as well.
\end{proof}

\begin{lemma}
\label{lemma-from-chow-to-K}
Let $X$ be a scheme locally of finite type over $(S, \delta)$
as in Situation \ref{situation-setup}. There is a canonical map
$$
\CH_k(X)
\longrightarrow
K_0(\textit{Coh}_{\leq k + 1}(X)/\textit{Coh}_{\leq k - 1}(X))
$$
induced by the map
$Z_k(X) \to K_0(\textit{Coh}_{\leq k}(X)/\textit{Coh}_{\leq k - 1}(X))$
from Lemma \ref{lemma-cycles-k-group}.
\end{lemma}

\begin{proof}
We have to show that an element $\alpha$ of $Z_k(X)$ which is rationally
equivalent to zero, is mapped to zero in
$K_0(\textit{Coh}_{\leq k + 1}(X)/\textit{Coh}_{\leq k - 1}(X))$.
Write $\alpha = \sum (i_j)_*\text{div}(f_j)$ as in
Definition \ref{definition-rational-equivalence}.
Observe that
$$
\pi = \coprod i_j : W = \coprod W_j \longrightarrow X
$$
is a finite morphism as each $i_j : W_j \to X$ is a closed immersion
and the family of $W_j$ is locally finite in $X$. Hence we may use
Lemma \ref{lemma-finite-cycles-k-group} to reduce to the case of $W$.
Since $W$ is a disjoint union of integral scheme, we reduce
to the case discussed in the next paragraph.

\medskip\noindent
Assume $X$ is integral of $\delta$-dimension $k + 1$.
Let $f$ be a nonzero rational function on $X$.
Let $\alpha = \text{div}(f)$. We have to show that
$\alpha$ is mapped to zero in
$K_0(\textit{Coh}_{\leq k + 1}(X)/\textit{Coh}_{\leq k - 1}(X))$.
Let $\mathcal{I} \subset \mathcal{O}_X$ be the ideal of denominators
of $f$, see Divisors, Definition
\ref{divisors-definition-regular-meromorphic-ideal-denominators}.
Then we have short exact sequences
$$
0 \to \mathcal{I} \to \mathcal{O}_X \to \mathcal{O}_X/\mathcal{I} \to 0
$$
and
$$
0 \to \mathcal{I} \xrightarrow{f} \mathcal{O}_X \to
\mathcal{O}_X/f\mathcal{I} \to 0
$$
See Divisors, Lemma
\ref{divisors-lemma-regular-meromorphic-ideal-denominators}.
We claim that
$$
[\mathcal{O}_X/\mathcal{I}]_k - [\mathcal{O}_X/f\mathcal{I}]_k =
\text{div}(f)
$$
The claim implies the element $\alpha = \text{div}(f)$ is represented by
$[\mathcal{O}_X/\mathcal{I}] - [\mathcal{O}_X/f\mathcal{I}]$
in $K_0(\textit{Coh}_{\leq k}(X)/\textit{Coh}_{\leq k - 1}(X))$.
Then the short exact sequences show that this element maps to
zero in $K_0(\textit{Coh}_{\leq k + 1}(X)/\textit{Coh}_{\leq k - 1}(X))$.

\medskip\noindent
To prove the claim, let $Z \subset X$ be an integral closed subscheme
of $\delta$-dimension $k$ and let $\xi \in Z$ be its generic point.
Then $I = \mathcal{I}_\xi \subset A = \mathcal{O}_{X, \xi}$
is an ideal such that $fI \subset A$. Now the coefficient of
$[Z]$ in $\text{div}(f)$ is $\text{ord}_A(f)$. (Of course as usual
we identify the function field of $X$ with the fraction field of $A$.)
On the other hand, the coefficient of $[Z]$ in
$[\mathcal{O}_X/\mathcal{I}] - [\mathcal{O}_X/f\mathcal{I}]$
is
$$
\text{length}_A(A/I) - \text{length}_A(A/fI)
$$
Using the distance function of
Algebra, Definition \ref{algebra-definition-distance}
we can rewrite this as
$$
d(A, I) - d(A, fI) = d(I, fI) = \text{ord}_A(f)
$$
The equalities hold by Algebra, Lemmas
\ref{algebra-lemma-properties-distance-function} and
\ref{algebra-lemma-order-vanishing-determinant}.
(Using these lemmas isn't necessary, but convenient.)
\end{proof}

\begin{remark}
\label{remark-good-cases-K-A}
Let $(S, \delta)$ be as in Situation \ref{situation-setup}.
Let $X$ be a scheme locally of finite type over $S$.
We will see later (in Lemma \ref{lemma-cycles-rational-equivalence-K-group})
that the map
$$
\CH_k(X)
\longrightarrow
K_0(\textit{Coh}_{k + 1}(X)/\textit{Coh}_{\leq k - 1}(X))
$$
of Lemma \ref{lemma-from-chow-to-K} is injective.
Composing with the canonical map
$$
K_0(\textit{Coh}_{k + 1}(X)/\textit{Coh}_{\leq k - 1}(X))
\longrightarrow
K_0(\textit{Coh}(X)/\textit{Coh}_{\leq k - 1}(X))
$$
we obtain a canonical map
$$
\CH_k(X)
\longrightarrow
K_0(\textit{Coh}(X)/\textit{Coh}_{\leq k - 1}(X)).
$$
We have not been able to find a statement or conjecture in the
literature as to whether this map should be injective or not.
It seems reasonable to expect the kernel of this map to be torsion.
We will return to this question (insert future reference).
\end{remark}

\begin{lemma}
\label{lemma-K-coherent-supported-on-closed}
Let $X$ be a locally Noetherian scheme. Let $Z \subset X$ be a closed
subscheme. Denote $\textit{Coh}_Z(X) \subset \textit{Coh}(X)$
the Serre subcategory of coherent $\mathcal{O}_X$-modules whose
set theoretic support is contained in $Z$. Then the exact inclusion
functor $\textit{Coh}(Z) \to \textit{Coh}_Z(X)$ induces
an isomorphism
$$
K'_0(Z) = K_0(\textit{Coh}(Z)) \longrightarrow K_0(\textit{Coh}_Z(X))
$$
\end{lemma}

\begin{proof}
Let $\mathcal{F}$ be an object of $\textit{Coh}_Z(X)$.
Let $\mathcal{I} \subset \mathcal{O}_X$ be the quasi-coherent
ideal sheaf of $Z$. Consider the descending filtration
$$
\ldots \subset
\mathcal{F}^p = \mathcal{I}^p \mathcal{F} \subset
\mathcal{F}^{p - 1} \subset \ldots \subset \mathcal{F}^0 = \mathcal{F}
$$
Exactly as in the proof of Lemma \ref{lemma-from-chow-to-K} this filtration
is locally finite and hence
$\bigoplus_{p \geq 0} \mathcal{F}^p$,
$\bigoplus_{p \geq 1} \mathcal{F}^p$, and
$\bigoplus_{p \geq 0} \mathcal{F}^p/\mathcal{F}^{p + 1}$
are coherent $\mathcal{O}_X$-modules supported on $Z$.
Hence we get
$$
[\mathcal{F}] =
[\bigoplus\nolimits_{p \geq 0} \mathcal{F}^p/\mathcal{F}^{p + 1}]
$$
in $K_0(\textit{Coh}_Z(X))$ exactly as in the proof of
Lemma \ref{lemma-from-chow-to-K}. Since the coherent module
$\bigoplus_{p \geq 0} \mathcal{F}^p/\mathcal{F}^{p + 1}$
is annihilated by $\mathcal{I}$ we conclude that
$[\mathcal{F}]$ is in the image. Actually, we claim that the map
$$
\mathcal{F} \longmapsto 
c(\mathcal{F}) =
[\bigoplus\nolimits_{p \geq 0} \mathcal{F}^p/\mathcal{F}^{p + 1}]
$$
factors through $K_0(\textit{Coh}_Z(X))$ and is an inverse to
the map in the statement of the lemma. To see this all we have
to show is that if
$$
0 \to \mathcal{F} \to \mathcal{G} \to \mathcal{H} \to 0
$$
is a short exact sequence in $\textit{Coh}_Z(X)$, then we
get $c(\mathcal{G}) = c(\mathcal{F}) + c(\mathcal{H})$.
Observe that for all $q \geq 0$ we have a short exact sequence
$$
0 \to
(\mathcal{F} \cap \mathcal{I}^q\mathcal{G})/
(\mathcal{F} \cap \mathcal{I}^{q + 1}\mathcal{G}) \to
\mathcal{G}^q/\mathcal{G}^{q + 1} \to
\mathcal{H}^q/\mathcal{H}^{q + 1} \to 0
$$
For $p, q \geq 0$ consider the coherent submodule
$$
\mathcal{F}^{p, q} = \mathcal{I}^p\mathcal{F} \cap \mathcal{I}^q\mathcal{G}
$$
Arguing exactly as above and using that the filtrations
$\mathcal{F}^p = \mathcal{I}^p\mathcal{F}$ and
$\mathcal{F} \cap \mathcal{I}^q\mathcal{G}$ are locally finite,
we find that
$$
[\bigoplus\nolimits_{p \geq 0} \mathcal{F}^p/\mathcal{F}^{p + 1}] =
[\bigoplus\nolimits_{p, q \geq 0}
\mathcal{F}^{p, q}/(\mathcal{F}^{p + 1, q} + \mathcal{F}^{p, q + 1})] =
[\bigoplus\nolimits_{q \geq 0}
(\mathcal{F} \cap \mathcal{I}^q\mathcal{G})/
(\mathcal{F} \cap \mathcal{I}^{q + 1}\mathcal{G})]
$$
in $K_0(\textit{Coh}(Z))$. Combined with the exact sequences above we obtain
the desired result. Some details omitted.
\end{proof}













\section{The divisor associated to an invertible sheaf}
\label{section-divisor-invertible-sheaf}

\noindent
The following definition is the analogue of
Divisors, Definition \ref{divisors-definition-divisor-invertible-sheaf}
in our current setup.

\begin{definition}
\label{definition-divisor-invertible-sheaf}
Let $(S, \delta)$ be as in Situation \ref{situation-setup}.
Let $X$ be locally of finite type over $S$. Assume $X$ is
integral and $n = \dim_\delta(X)$.
Let $\mathcal{L}$ be an invertible $\mathcal{O}_X$-module.
\begin{enumerate}
\item For any nonzero meromorphic section $s$ of $\mathcal{L}$
we define the {\it Weil divisor associated to $s$} is the
$(n - 1)$-cycle
$$
\text{div}_\mathcal{L}(s) =
\sum \text{ord}_{Z, \mathcal{L}}(s) [Z]
$$
defined in Divisors, Definition
\ref{divisors-definition-divisor-invertible-sheaf}.
This makes sense because Weil divisors have $\delta$-dimension $n - 1$
by Lemma \ref{lemma-divisor-delta-dimension}.
\item We define {\it Weil divisor associated to $\mathcal{L}$} as
$$
c_1(\mathcal{L}) \cap [X] =
\text{class of }\text{div}_\mathcal{L}(s) \in \CH_{n - 1}(X)
$$
where $s$ is any nonzero meromorphic section of $\mathcal{L}$ over
$X$. This is well defined by
Divisors, Lemma \ref{divisors-lemma-divisor-meromorphic-well-defined}.
\end{enumerate}
\end{definition}

\noindent
Let $X$ and $S$ be as in Definition \ref{definition-divisor-invertible-sheaf}
above. Set $n = \dim_\delta(X)$. It is clear from the definitions that
$Cl(X) = \CH_{n - 1}(X)$ where $Cl(X)$ is the Weil divisor class group of $X$
as defined in Divisors, Definition \ref{divisors-definition-class-group}.
The map
$$
\Pic(X) \longrightarrow \CH_{n - 1}(X), \quad
\mathcal{L} \longmapsto c_1(\mathcal{L}) \cap [X]
$$
is the same as the map $\Pic(X) \to Cl(X)$ constructed in
Divisors, Equation (\ref{divisors-equation-c1}) for arbitrary
locally Noetherian integral schemes. In particular, this map
is a homomorphism of abelian groups, it is injective if $X$ is
a normal scheme, and an isomorphism if all local rings of $X$
are UFDs. See Divisors, Lemmas \ref{divisors-lemma-normal-c1-injective} and
\ref{divisors-lemma-local-rings-UFD-c1-bijective}.
There are some cases where it is easy to compute the
Weil divisor associated to an invertible sheaf.

\begin{lemma}
\label{lemma-compute-c1}
Let $(S, \delta)$ be as in Situation \ref{situation-setup}.
Let $X$ be locally of finite type over $S$. Assume $X$ is
integral and $n = \dim_\delta(X)$.
Let $\mathcal{L}$ be an invertible $\mathcal{O}_X$-module.
Let $s \in \Gamma(X, \mathcal{L})$ be a nonzero global section.
Then
$$
\text{div}_\mathcal{L}(s) = [Z(s)]_{n - 1}
$$
in $Z_{n - 1}(X)$ and
$$
c_1(\mathcal{L}) \cap [X] = [Z(s)]_{n - 1}
$$
in $\CH_{n - 1}(X)$.
\end{lemma}

\begin{proof}
Let $Z \subset X$ be an integral closed subscheme of
$\delta$-dimension $n - 1$. Let $\xi \in Z$ be its generic
point. Choose a generator $s_\xi \in \mathcal{L}_\xi$.
Write $s = fs_\xi$ for some $f \in \mathcal{O}_{X, \xi}$.
By definition of $Z(s)$, see
Divisors, Definition \ref{divisors-definition-zero-scheme-s}
we see that $Z(s)$ is cut out by a quasi-coherent
sheaf of ideals $\mathcal{I} \subset \mathcal{O}_X$ such
that $\mathcal{I}_\xi = (f)$. Hence
$\text{length}_{\mathcal{O}_{X, \xi}}(\mathcal{O}_{Z(s), \xi})
=
\text{length}_{\mathcal{O}_{X, \xi}}(\mathcal{O}_{X, \xi}/(f))
=
\text{ord}_{\mathcal{O}_{X, \xi}}(f)$ as desired.
\end{proof}

\noindent
The following lemma will be superseded by the more general
Lemma \ref{lemma-flat-pullback-cap-c1}.

\begin{lemma}
\label{lemma-flat-pullback-divisor-invertible-sheaf}
Let $(S, \delta)$ be as in Situation \ref{situation-setup}.
Let $X$, $Y$ be locally of finite type over $S$. Assume $X$, $Y$
are integral and $n = \dim_\delta(Y)$.
Let $\mathcal{L}$ be an invertible $\mathcal{O}_Y$-module.
Let $f : X \to Y$ be a flat morphism of relative dimension $r$. Then
$$
f^*(c_1(\mathcal{L}) \cap [Y]) = c_1(f^*\mathcal{L}) \cap [X]
$$
in $\CH_{n + r - 1}(X)$.
\end{lemma}

\begin{proof}
Let $s$ be a nonzero meromorphic section of $\mathcal{L}$.
We will show that actually
$f^*\text{div}_\mathcal{L}(s) = \text{div}_{f^*\mathcal{L}}(f^*s)$
and hence the lemma holds.
To see this let $\xi \in Y$ be a point and let $s_\xi \in \mathcal{L}_\xi$
be a generator. Write $s = gs_\xi$ with $g \in R(Y)^*$.
Then there is an open neighbourhood $V \subset Y$ of $\xi$
such that $s_\xi \in \mathcal{L}(V)$ and such that $s_\xi$ generates
$\mathcal{L}|_V$. Hence we see that
$$
\text{div}_\mathcal{L}(s)|_V = \text{div}_Y(g)|_V.
$$
In exactly the same way, since $f^*s_\xi$ generates $f^*\mathcal{L}$
over $f^{-1}(V)$ and since $f^*s = g f^*s_\xi$ we also
have
$$
\text{div}_\mathcal{L}(f^*s)|_{f^{-1}(V)}
=
\text{div}_X(g)|_{f^{-1}(V)}.
$$
Thus the desired equality of cycles over $f^{-1}(V)$ follows from the
corresponding result for pullbacks of principal divisors, see
Lemma \ref{lemma-flat-pullback-principal-divisor}.
\end{proof}



\section{Intersecting with an invertible sheaf}
\label{section-intersecting-with-divisors}

\noindent
In this section we study the following construction.

\begin{definition}
\label{definition-cap-c1}
Let $(S, \delta)$ be as in Situation \ref{situation-setup}.
Let $X$ be locally of finite type over $S$.
Let $\mathcal{L}$ be an invertible $\mathcal{O}_X$-module.
We define, for every integer $k$, an operation
$$
c_1(\mathcal{L}) \cap - :
Z_{k + 1}(X) \to \CH_k(X)
$$
called {\it intersection with the first Chern class of $\mathcal{L}$}.
\begin{enumerate}
\item Given an integral closed subscheme $i : W \to X$ with
$\dim_\delta(W) = k + 1$ we define
$$
c_1(\mathcal{L}) \cap [W] = i_*(c_1({i^*\mathcal{L}}) \cap [W])
$$
where the right hand side is defined in
Definition \ref{definition-divisor-invertible-sheaf}.
\item For a general $(k + 1)$-cycle $\alpha = \sum n_i [W_i]$ we set
$$
c_1(\mathcal{L}) \cap \alpha = \sum n_i c_1(\mathcal{L}) \cap [W_i]
$$
\end{enumerate}
\end{definition}

\noindent
Write each $c_1(\mathcal{L}) \cap W_i = \sum_j n_{i, j} [Z_{i, j}]$
with $\{Z_{i, j}\}_j$ a locally finite sum
of integral closed subschemes of $W_i$. Since $\{W_i\}$ is a locally
finite collection of integral closed subschemes on $X$, it follows
easily that $\{Z_{i, j}\}_{i, j}$ is a locally finite collection
of closed subschemes of $X$. Hence
$c_1(\mathcal{L}) \cap \alpha = \sum n_in_{i, j}[Z_{i, j}]$
is a cycle. Another, more convenient, way to think about this
is to observe that the morphism $\coprod W_i \to X$ is
proper. Hence $c_1(\mathcal{L}) \cap \alpha$ can be viewed
as the pushforward of a class in $\CH_k(\coprod W_i) = \prod \CH_k(W_i)$.
This also explains why the result is well defined up to rational
equivalence on $X$.

\medskip\noindent
The main goal for the next few sections is to show that intersecting with
$c_1(\mathcal{L})$ factors through rational equivalence.
This is not a triviality.

\begin{lemma}
\label{lemma-c1-cap-additive}
Let $(S, \delta)$ be as in Situation \ref{situation-setup}.
Let $X$ be locally of finite type over $S$.
Let $\mathcal{L}$, $\mathcal{N}$ be an invertible sheaves on $X$.
Then
$$
c_1(\mathcal{L}) \cap \alpha  + c_1(\mathcal{N}) \cap \alpha =
c_1(\mathcal{L} \otimes_{\mathcal{O}_X} \mathcal{N}) \cap \alpha
$$
in $\CH_k(X)$ for every $\alpha \in Z_{k + 1}(X)$. Moreover,
$c_1(\mathcal{O}_X) \cap \alpha = 0$ for all $\alpha$.
\end{lemma}

\begin{proof}
The additivity follows directly from
Divisors, Lemma \ref{divisors-lemma-c1-additive}
and the definitions. To see that $c_1(\mathcal{O}_X) \cap \alpha = 0$
consider the section $1 \in \Gamma(X, \mathcal{O}_X)$. This restricts
to an everywhere nonzero section on any integral closed subscheme
$W \subset X$. Hence $c_1(\mathcal{O}_X) \cap [W] = 0$ as desired.
\end{proof}

\noindent
Recall that $Z(s) \subset X$ denotes the zero scheme of a global section
$s$ of an invertible sheaf on a scheme $X$, see
Divisors, Definition \ref{divisors-definition-zero-scheme-s}.

\begin{lemma}
\label{lemma-prepare-geometric-cap}
Let $(S, \delta)$ be as in Situation \ref{situation-setup}.
Let $Y$ be locally of finite type over $S$.
Let $\mathcal{L}$ be an invertible $\mathcal{O}_Y$-module.
Let $s \in \Gamma(Y, \mathcal{L})$.
Assume
\begin{enumerate}
\item $\dim_\delta(Y) \leq k + 1$,
\item $\dim_\delta(Z(s)) \leq k$, and
\item for every generic point $\xi$ of an irreducible component of
$Z(s)$ of $\delta$-dimension $k$ the multiplication by $s$
induces an injection $\mathcal{O}_{Y, \xi} \to \mathcal{L}_\xi$.
\end{enumerate}
Write $[Y]_{k + 1} = \sum n_i[Y_i]$ where $Y_i \subset Y$ are the
irreducible components of $Y$ of $\delta$-dimension $k + 1$.
Set $s_i = s|_{Y_i} \in \Gamma(Y_i, \mathcal{L}|_{Y_i})$. Then
\begin{equation}
\label{equation-equal-as-cycles}
[Z(s)]_k =  \sum n_i[Z(s_i)]_k
\end{equation}
as $k$-cycles on $Y$.
\end{lemma}

\begin{proof}
Let $Z \subset Y$ be an integral closed subscheme of
$\delta$-dimension $k$. Let $\xi \in Z$ be its generic point.
We want to compare the coefficient $n$ of $[Z]$ in the expression
$\sum n_i[Z(s_i)]_k$ with the coefficient $m$ of $[Z]$ in the
expression $[Z(s)]_k$. Choose a generator $s_\xi \in \mathcal{L}_\xi$.
Write $A = \mathcal{O}_{Y, \xi}$, $L = \mathcal{L}_\xi$.
Then $L = As_\xi$. Write $s = f s_\xi$ for some (unique) $f \in A$.
Hypothesis (3) means that $f : A \to A$ is injective.
Since $\dim_\delta(Y) \leq k + 1$ and $\dim_\delta(Z) = k$
we have $\dim(A) = 0$ or $1$. We have
$$
m = \text{length}_A(A/(f))
$$
which is finite in either case.

\medskip\noindent
If $\dim(A) = 0$, then $f : A \to A$ being injective
implies that $f \in A^*$. Hence in this case $m$ is zero.
Moreover, the condition $\dim(A) = 0$ means that $\xi$
does not lie on any irreducible component of $\delta$-dimension
$k + 1$, i.e., $n = 0$ as well.

\medskip\noindent
Now, let $\dim(A) = 1$.
Since $A$ is a Noetherian local ring it has finitely
many minimal primes $\mathfrak q_1, \ldots, \mathfrak q_t$.
These correspond 1-1 with the $Y_i$ passing through $\xi'$.
Moreover $n_i = \text{length}_{A_{\mathfrak q_i}}(A_{\mathfrak q_i})$.
Also, the multiplicity of $[Z]$ in $[Z(s_i)]_k$ is
$\text{length}_A(A/(f, \mathfrak q_i))$.
Hence the equation to prove in this case is
$$
\text{length}_A(A/(f))
=
\sum \text{length}_{A_{\mathfrak q_i}}(A_{\mathfrak q_i})
\text{length}_A(A/(f, \mathfrak q_i))
$$
which follows from
Lemma \ref{lemma-additivity-divisors-restricted}.
\end{proof}

\noindent
The following lemma is a useful result in order to compute the intersection
product of the $c_1$ of an invertible sheaf and the cycle associated
to a closed subscheme.
Recall that $Z(s) \subset X$ denotes the zero scheme of a global section
$s$ of an invertible sheaf on a scheme $X$, see
Divisors, Definition \ref{divisors-definition-zero-scheme-s}.

\begin{lemma}
\label{lemma-geometric-cap}
Let $(S, \delta)$ be as in Situation \ref{situation-setup}.
Let $X$ be locally of finite type over $S$.
Let $\mathcal{L}$ be an invertible $\mathcal{O}_X$-module.
Let $Y \subset X$ be a closed subscheme.
Let $s \in \Gamma(Y, \mathcal{L}|_Y)$.
Assume
\begin{enumerate}
\item $\dim_\delta(Y) \leq k + 1$,
\item $\dim_\delta(Z(s)) \leq k$, and
\item for every generic point $\xi$ of an irreducible component of
$Z(s)$ of $\delta$-dimension $k$ the multiplication by $s$
induces an injection
$\mathcal{O}_{Y, \xi} \to (\mathcal{L}|_Y)_\xi$\footnote{For example,
this holds if $s$ is a regular section of $\mathcal{L}|_Y$.}.
\end{enumerate}
Then
$$
c_1(\mathcal{L}) \cap [Y]_{k + 1} = [Z(s)]_k
$$
in $\CH_k(X)$.
\end{lemma}

\begin{proof}
Write
$$
[Y]_{k + 1} = \sum n_i[Y_i]
$$
where $Y_i \subset Y$ are the irreducible components of
$Y$ of $\delta$-dimension $k + 1$ and $n_i > 0$.
By assumption the restriction
$s_i = s|_{Y_i} \in \Gamma(Y_i, \mathcal{L}|_{Y_i})$ is not
zero, and hence is a regular section. By Lemma \ref{lemma-compute-c1}
we see that $[Z(s_i)]_k$ represents $c_1(\mathcal{L}|_{Y_i})$.
Hence by definition
$$
c_1(\mathcal{L}) \cap [Y]_{k + 1} = \sum n_i[Z(s_i)]_k
$$
Thus the result follows from Lemma \ref{lemma-prepare-geometric-cap}.
\end{proof}




\section{Intersecting with an invertible sheaf and push and pull}
\label{section-intersecting-with-divisors-push-pull}

\noindent
In this section we prove that the operation $c_1(\mathcal{L}) \cap -$
commutes with flat pullback and proper pushforward.

\begin{lemma}
\label{lemma-prepare-flat-pullback-cap-c1}
Let $(S, \delta)$ be as in Situation \ref{situation-setup}.
Let $X$, $Y$ be locally of finite type over $S$.
Let $f : X \to Y$ be a flat morphism of relative dimension $r$.
Let $\mathcal{L}$ be an invertible sheaf on $Y$.
Assume $Y$ is integral and $n = \dim_\delta(Y)$.
Let $s$ be a nonzero meromorphic section of $\mathcal{L}$.
Then we have
$$
f^*\text{div}_\mathcal{L}(s) = \sum n_i\text{div}_{f^*\mathcal{L}|_{X_i}}(s_i)
$$
in $Z_{n + r - 1}(X)$. Here the sum is over the irreducible
components $X_i \subset X$ of $\delta$-dimension $n + r$,
the section $s_i = f|_{X_i}^*(s)$ is the pullback of $s$, and
$n_i = m_{X_i, X}$ is the multiplicity of $X_i$ in $X$.
\end{lemma}

\begin{proof}
To prove this equality of cycles, we may work locally on $Y$.
Hence we may assume $Y$ is affine and $s = p/q$ for some nonzero
sections $p \in \Gamma(Y, \mathcal{L})$ and $q \in \Gamma(Y, \mathcal{O})$.
If we can show both
$$
f^*\text{div}_\mathcal{L}(p) =
\sum n_i\text{div}_{f^*\mathcal{L}|_{X_i}}(p_i)
\quad\text{and}\quad
f^*\text{div}_\mathcal{O}(q) =
\sum n_i\text{div}_{\mathcal{O}_{X_i}}(q_i)
$$
(with obvious notations) then we win by the
additivity, see Divisors, Lemma \ref{divisors-lemma-c1-additive}.
Thus we may assume that $s \in \Gamma(Y, \mathcal{L})$.
In this case we may apply the equality
(\ref{equation-equal-as-cycles}) to see that
$$
[Z(f^*(s))]_{k + r - 1} =
\sum n_i\text{div}_{f^*\mathcal{L}|_{X_i}}(s_i)
$$
where $f^*(s) \in f^*\mathcal{L}$ denotes the pullback of $s$ to $X$.
On the other hand we have
$$
f^*\text{div}_\mathcal{L}(s) = f^*[Z(s)]_{k - 1}
= [f^{-1}(Z(s))]_{k + r - 1},
$$
by Lemmas \ref{lemma-compute-c1} and \ref{lemma-pullback-coherent}.
Since $Z(f^*(s)) = f^{-1}(Z(s))$ we win.
\end{proof}

\begin{lemma}
\label{lemma-flat-pullback-cap-c1}
Let $(S, \delta)$ be as in Situation \ref{situation-setup}.
Let $X$, $Y$ be locally of finite type over $S$.
Let $f : X \to Y$ be a flat morphism of relative dimension $r$.
Let $\mathcal{L}$ be an invertible sheaf on $Y$.
Let $\alpha$ be a $k$-cycle on $Y$.
Then
$$
f^*(c_1(\mathcal{L}) \cap \alpha) = c_1(f^*\mathcal{L}) \cap f^*\alpha
$$
in $\CH_{k + r - 1}(X)$.
\end{lemma}

\begin{proof}
Write $\alpha = \sum n_i[W_i]$. We will show that
$$
f^*(c_1(\mathcal{L}) \cap [W_i]) = c_1(f^*\mathcal{L}) \cap f^*[W_i]
$$
in $\CH_{k + r - 1}(X)$ by producing a rational equivalence
on the closed subscheme $f^{-1}(W_i)$ of $X$.
By the discussion in
Remark \ref{remark-infinite-sums-rational-equivalences}
this will prove the equality of the lemma is true.

\medskip\noindent
Let $W \subset Y$ be an integral closed subscheme of $\delta$-dimension $k$.
Consider the closed subscheme $W' = f^{-1}(W) = W \times_Y X$
so that we have the fibre product diagram
$$
\xymatrix{
W' \ar[r] \ar[d]_h & X \ar[d]^f \\
W \ar[r] & Y
}
$$
We have to show that
$f^*(c_1(\mathcal{L}) \cap [W]) = c_1(f^*\mathcal{L}) \cap f^*[W]$.
Choose a nonzero meromorphic section $s$ of $\mathcal{L}|_W$.
Let $W'_i \subset W'$ be the irreducible components of
$\delta$-dimension $k + r$. Write $[W']_{k + r} = \sum n_i[W'_i]$
with $n_i$ the multiplicity of $W'_i$ in $W'$ as per definition.
So $f^*[W] = \sum n_i[W'_i]$ in $Z_{k + r}(X)$.
Since each $W'_i \to W$ is dominant we
see that $s_i = s|_{W'_i}$ is a nonzero meromorphic section for
each $i$. By Lemma \ref{lemma-prepare-flat-pullback-cap-c1}
we have the following equality of cycles
$$
h^*\tex